% !TEX program = xelatex
\documentclass[aspectratio=169,10pt]{beamer}

\usepackage{fontspec}
\setmainfont{PT Serif}
\setsansfont{PT Sans}
\setmonofont{PT Mono}
\usepackage{polyglossia}
\setdefaultlanguage{russian}
\setotherlanguage{english}
\newfontfamily\cyrillicfont{PT Serif}
\newfontfamily\cyrillicfontsf{PT Sans}
\newfontfamily\cyrillicfonttt{PT Mono}

\usepackage{amsmath,amssymb,bm}
\usepackage{booktabs}
\usepackage{tabularx}
\usepackage{array}
\usepackage{graphicx}
\usepackage{tikz}
\usepackage{pgfplots}
\pgfplotsset{compat=1.18}
\usetikzlibrary{arrows.meta,positioning,calc,fit,shapes.geometric}
\graphicspath{{figures/molecules/}{figures/generated/}}

\usetheme{Madrid}
\setbeamertemplate{navigation symbols}{}
\setbeamertemplate{headline}{}
\setbeamertemplate{blocks}[rounded][shadow=false]
\setbeamerfont{page number in head/foot}{size=\scriptsize}
\setbeamertemplate{footline}{%
  \leavevmode%
  \hbox{\begin{beamercolorbox}[wd=\paperwidth,ht=2.2ex,dp=1ex,right]{page number in head/foot}%
    \usebeamerfont{page number in head/foot}\insertframenumber/\inserttotalframenumber\hspace*{2ex}%
  \end{beamercolorbox}}%
}
\setbeameroption{hide notes}

\definecolor{TGBlue}{HTML}{3B82F6}
\definecolor{TGBlueDark}{HTML}{1F5AA6}
\definecolor{TGOrange}{HTML}{F59E0B}
\definecolor{TGRed}{HTML}{EF4444}
\definecolor{TGGreen}{HTML}{14B8A6}
\definecolor{TGSlate}{HTML}{475569}
\definecolor{TGLight}{HTML}{CBD5E1}
\definecolor{TGSoft}{HTML}{F8FBFF}
\definecolor{TGPurple}{HTML}{6D28D9}
\definecolor{TGPaper}{HTML}{FFFDF7}

\setbeamercolor{structure}{fg=TGBlue}
\setbeamercolor{frametitle}{fg=white,bg=TGBlueDark}
\setbeamercolor{title}{fg=white,bg=TGBlueDark}
\setbeamercolor{block title}{fg=TGBlueDark,bg=TGBlue!10}
\setbeamercolor{block body}{bg=white,fg=black}
\setbeamercolor{alerted text}{fg=TGOrange}
\setbeamercolor{background canvas}{bg=TGSoft}
\setbeamertemplate{itemize item}{\textcolor{TGBlue}{\raisebox{0.15ex}{\tiny$\bullet$}}}
\setbeamertemplate{itemize subitem}{\textcolor{TGBlue}{\scriptsize—}}
\setbeamertemplate{itemize subsubitem}{\textcolor{TGBlue}{\tiny$\bullet$}}
\setbeamertemplate{enumerate item}{\textcolor{TGBlue}{\insertenumlabel.}}

\newcolumntype{Y}{>{\raggedright\arraybackslash}X}
\newcommand{\smallnote}[1]{{\scriptsize\color{TGSlate}#1}}
\newcommand{\keybox}[1]{\begin{block}{Ключевая мысль}#1\end{block}}
\newcommand{\fitfig}[1]{\resizebox{\linewidth}{!}{#1}}
\newcommand{\deckchip}[2][TGBlue]{%
  \tikz[baseline=(chip.base)]\node[
    anchor=base,
    draw=#1!26,
    fill=#1!9,
    text=#1!82!black,
    rounded corners=8pt,
    inner xsep=7pt,
    inner ysep=3pt,
    font=\scriptsize\bfseries
  ] (chip) {#2};%
}
\newcommand{\deckcard}[4]{%
  \begin{minipage}[t]{#1}
    \begin{tikzpicture}
      \node[
        draw=#2!30,
        fill=#2!7,
        rounded corners=7pt,
        inner sep=6pt,
        text width=0.90\linewidth,
        align=left
      ] {
        {\scriptsize\bfseries\color{#2!82!black}#3\par}
        \vspace{0.18em}
        {\scriptsize\color{TGSlate}#4}
      };
    \end{tikzpicture}
  \end{minipage}%
}
\newcommand{\molasset}[2][]{\includegraphics[#1]{#2}}
\newcommand{\plotasset}[2][]{%
  \IfFileExists{figures/generated/#2}{%
    \includegraphics[#1]{figures/generated/#2}%
  }{%
    \IfFileExists{figures/generated/#2.pdf}{%
      \includegraphics[#1]{figures/generated/#2.pdf}%
    }{%
      \IfFileExists{figures/generated/#2.png}{%
        \includegraphics[#1]{figures/generated/#2.png}%
      }{%
        \begin{tikzpicture}
          \node[draw=TGLight,fill=white,rounded corners=6pt,minimum width=0.92\linewidth,minimum height=3.2cm,align=center,text=TGSlate]
            {generated figure missing\\\scriptsize\texttt{\detokenize{figures/generated/#2}}};
        \end{tikzpicture}%
      }%
    }%
  }%
}
\newcommand{\tightitems}{\setlength{\itemsep}{0.25em}}
\newcommand{\barrow}[4]{%
  \node[anchor=east,font=\scriptsize] at (0,#1) {#2};
  \fill[#4!70] (0.15,#1-0.12) rectangle (0.15+#3,#1+0.12);
  \draw[#4] (0.15,#1-0.12) rectangle (0.15+#3,#1+0.12);
}
\newcommand{\roadmapbar}[1]{%
  \begin{tikzpicture}
    \foreach \i/\lab/\name in {
      1/I/Данные,
      2/II/Физика,
      3/III/Архитектура,
      4/IV/Обучение,
      5/V/Диагностика,
      6/VI/Обновления
    }{
      \pgfmathsetmacro{\x}{(\i-1)*1.82}
      \ifnum\i>#1
        \node[draw=TGLight,fill=white,text=TGSlate!58,rounded corners=4pt,
          minimum width=1.58cm,minimum height=0.76cm,align=center,font=\scriptsize]
          at (\x,0) {\lab\\\name};
      \else
        \node[draw=TGBlue,fill=TGBlue!13,text=TGBlueDark,rounded corners=4pt,
          minimum width=1.58cm,minimum height=0.76cm,align=center,font=\scriptsize\bfseries]
          at (\x,0) {\lab\\\name};
      \fi
    }
  \end{tikzpicture}%
}

\title[TGNN-Solv / TIMP]{Физически-информированная графовая нейронная сеть\\для предсказания растворимости твёрдых веществ в жидкостях}
\subtitle{От мастер-уравнения SLE до термодинамически-информированной передачи сообщений (TIMP)}
\author{Никита Поломошнов}
\institute{Лабораторный семинар}
\date{Апрель 2026}

\begin{document}

% 1
\begin{frame}[plain]
  \begin{tikzpicture}[remember picture,overlay]
    \fill[TGSoft] (current page.south west) rectangle (current page.north east);
    \fill[TGBlue!10] (current page.north west) rectangle ([yshift=-0.14cm]current page.north east);
    \fill[TGBlue] ([xshift=0.72cm,yshift=-2.62cm]current page.north west)
      rectangle ([xshift=0.86cm,yshift=-6.55cm]current page.north west);
    \node[anchor=north east,text=TGBlue!62,font=\Large\bfseries] at ([xshift=-0.82cm,yshift=-0.56cm]current page.north east)
      {$\ln x_2=-\Phi-\ln\gamma_2$};
    \node[anchor=north west,fill=TGBlue!9,draw=TGBlue!24,rounded corners=12pt,inner xsep=18pt,inner ysep=9pt,text=TGBlueDark,font=\small\bfseries]
      at ([xshift=1.05cm,yshift=-0.50cm]current page.north west)
      {TGNN-Solv · seminar};
    \node[anchor=north west,text width=0.84\paperwidth,align=left]
      at ([xshift=1.05cm,yshift=-2.78cm]current page.north west)
      {{\color{TGBlueDark}\fontsize{19}{22}\selectfont\bfseries
      Физически-информированная графовая\\
      нейронная сеть для предсказания\\
      растворимости твёрдых веществ в жидкостях\par}
      \vspace{0.45cm}
      {\color{TGSlate}\large От мастер-уравнения SLE до термодинамически-информированной передачи сообщений (TIMP)\par}};
    \node[anchor=south west,text=TGSlate,font=\normalsize]
      at ([xshift=1.05cm,yshift=0.55cm]current page.south west)
      {\insertauthor\quad|\quad\insertdate};
  \end{tikzpicture}
  \note{Открыть доклад как историю инженерного исследования: зачем физика, что именно кодируется в модели, почему потребовался TIMP.}
\end{frame}

\section{Постановка задачи}

% 2
\begin{frame}{Зачем нужно предсказывать растворимость}
  \begin{columns}[c,onlytextwidth]
    \begin{column}{0.50\textwidth}
      \begin{block}{Практический вопрос}
        \large
        В каком растворителе и при какой температуре вещество даст нужную растворимость?
      \end{block}
      \vspace{0.45em}
      \begin{itemize}\tightitems
        \item Кристаллизация, экстракция, очистка и выделение продукта требуют кривую $x_2(T)$.
        \item Планирование синтеза зависит от загрузки, среды реакции и антирастворителя.
        \item Полный экспериментальный перебор занимает часы или сутки на точку.
      \end{itemize}
    \end{column}
    \begin{column}{0.46\textwidth}
      \centering
      \fitfig{
      \begin{tikzpicture}
        \node[draw=TGBlue!22,fill=white,rounded corners=14pt,minimum width=4.7cm,minimum height=1.15cm,align=center] at (0,1.45)
          {\Large $N_{\text{веществ}}\times N_{\text{растворителей}}\times N_T\sim10^6$};
        \node[draw=TGGreen!24,fill=TGGreen!8,rounded corners=14pt,minimum width=4.7cm,minimum height=1.15cm,align=center] at (0,0)
          {\Large $y=\ln x_2\in[-25,0]$};
        \node[draw=TGOrange!24,fill=TGOrange!9,rounded corners=14pt,minimum width=4.7cm,minimum height=1.15cm,align=center,text width=4.1cm] at (0,-1.45)
          {нужна быстрая модель для ранжирования кандидатов};
      \end{tikzpicture}}
    \end{column}
  \end{columns}
  \note{Переход: если задача так велика, нужно выбрать между физическими моделями и машинным обучением, либо объединить их.}
\end{frame}

% 3
\begin{frame}{Два мира, которые не встречаются}
  \small
  \begin{columns}[T,onlytextwidth]
    \begin{column}{0.48\textwidth}
      \begin{block}{Физические модели}
        \scriptsize
        \textbf{Примеры:} UNIFAC, COSMO-RS.\\[0.25em]
        \textcolor{TGGreen}{Плюсы:} интерпретируемость, температурная структура, термодинамическая согласованность.\\
        \textcolor{TGRed}{Минусы:} ограниченный словарь групп, трудная адаптация к новому корпусу.
      \end{block}
    \end{column}
    \begin{column}{0.48\textwidth}
      \begin{block}{Модели машинного обучения}
        \scriptsize
        \textbf{Примеры:} табличные модели, графовые нейросети.\\[0.25em]
        \textcolor{TGGreen}{Плюсы:} учатся на данных, быстро адаптируются, используют большие корпуса.\\
        \textcolor{TGRed}{Минусы:} слабая интерпретация, риск слабой экстраполяции по температуре.
      \end{block}
    \end{column}
  \end{columns}
  \vspace{0.5em}
  \begin{center}
    \Large\textcolor{TGOrange}{Идея проекта: нейросеть предсказывает параметры, физика считает ответ}
  \end{center}
  \note{Разделить плюсы и минусы: физика даёт структуру, машинное обучение даёт адаптацию к данным.}
\end{frame}

% Early notation primer for a broad audience
\begin{frame}{Сначала договоримся об обозначениях}
  \small
  \begin{columns}[T,onlytextwidth]
    \begin{column}{0.49\textwidth}
      \begin{block}{Бинарная смесь}
        \scriptsize
        \begin{itemize}\tightitems
          \item \textbf{Компонент 1} — растворитель.
          \item \textbf{Компонент 2} — растворяемое твёрдое вещество.
          \item $n_i$ — число молей компонента $i$.
          \item $x_i=\dfrac{n_i}{n_1+n_2}$ — мольная доля.
          \item Для бинарной смеси: $x_1+x_2=1$.
          \item Цель модели: $y=\ln x_2$.
        \end{itemize}
      \end{block}
    \end{column}
    \begin{column}{0.49\textwidth}
      \begin{block}{Физхимические величины}
        \scriptsize
        \begin{itemize}\tightitems
          \item $\gamma_i$ — коэффициент активности: поправка к идеальной смеси.
          \item $G^E$ — избыточная энергия Гиббса: отличие реальной смеси от идеальной.
          \item $T_m$ — температура плавления вещества.
          \item $\Delta H_{\mathrm{fus}}$ — энтальпия плавления.
          \item $\tau_{12},\tau_{21},\alpha$ — параметры NRTL для неидеальности.
          \item $\Phi$ — кристаллический штраф в SLE.
        \end{itemize}
      \end{block}
    \end{column}
  \end{columns}
  \smallnote{$R$ — газовая постоянная, $T$ — температура в K. Деление энергии на $RT$ делает параметры безразмерными.}
  \note{Зафиксировать обозначения до данных: 1 всегда растворитель, 2 всегда растворяемое вещество, x2 — именно растворимость вещества.}
\end{frame}

\begin{frame}{Как читать индексы $1$ и $2$}
  \small
  \centering
  \begin{tikzpicture}[>=Latex]
    \node[draw=TGBlue!28,fill=TGBlue!8,rounded corners=10pt,minimum width=3.4cm,minimum height=1.20cm,align=center] (slv) at (0,0)
      {\textbf{1: растворитель}\\{\scriptsize например этанол}};
    \node[draw=TGOrange!30,fill=TGOrange!9,rounded corners=10pt,minimum width=3.4cm,minimum height=1.20cm,align=center] (sol) at (5.0,0)
      {\textbf{2: вещество}\\{\scriptsize например парацетамол}};
    \draw[->,line width=1.0pt,TGSlate!65,bend left=18] (slv.north east) to node[above,font=\scriptsize,text=TGSlate] {$\tau_{12}$: 1 вокруг 2} (sol.north west);
    \draw[->,line width=1.0pt,TGSlate!65,bend left=18] (sol.south west) to node[below,font=\scriptsize,text=TGSlate] {$\tau_{21}$: 2 вокруг 1} (slv.south east);
    \node[draw=TGGreen!26,fill=TGGreen!8,rounded corners=10pt,align=left,text width=8.0cm,inner sep=7pt] at (2.5,-2.45)
      {\scriptsize
      $x_2$ — сколько вещества растворилось в смеси.\\
      $\gamma_2(x_2)$ — насколько вещество ведёт себя неидеально при этой концентрации.\\
      $\ln x_2=-\Phi-\ln\gamma_2$ — кристалл задаёт штраф, растворитель добавляет совместимость.};
  \end{tikzpicture}
  \note{Этот слайд специально снимает путаницу: во всей физической части 1 — растворитель, 2 — вещество.}
\end{frame}

% 4
\begin{frame}{Что такое одна строка данных}
  \small
  \begin{columns}[T,onlytextwidth]
    \begin{column}{0.48\textwidth}
      Каждая экспериментальная точка описывает насыщенный раствор:
      \[
        (\text{вещество},\;\text{растворитель},\;T)\;\longrightarrow\;\ln x_2.
      \]
      \begin{itemize}\tightitems
        \item $x_2$ — мольная доля растворённого вещества.
        \item $T$ — температура измерения.
        \item Одна пара обычно имеет несколько температурных точек.
      \end{itemize}
    \end{column}
    \begin{column}{0.48\textwidth}
      \begin{block}{BigSolDB после обработки}
        \scriptsize
        108\,287 размеченных строк\\
        213 растворителей\\
        около 1\,500 растворённых веществ\\
        медиана: 9 температурных точек на пару
      \end{block}
      \begin{block}{Почему это важно}
        \scriptsize
        Модель должна предсказывать не только одну точку, а кривую растворимости $x_2(T)$.
      \end{block}
    \end{column}
  \end{columns}
  \note{Сразу ввести, что именно лежит в CSV: пара, температура, логарифм мольной доли.}
\end{frame}

% 5
\begin{frame}{Данные: метки, вода и разбиения}
  \small
  \begin{columns}[T,onlytextwidth]
    \begin{column}{0.58\textwidth}
      \scriptsize
      \begin{tabularx}{\linewidth}{lY}
        \toprule
        \textbf{Метка} & \textbf{Зачем нужна}\\
        \midrule
        $T_m$ & температура плавления, супервизия кристаллического блока\\
        $\Delta H_{\mathrm{fus}}$ & энтальпия плавления, супервизия кристаллического штрафа\\
        Hansen & дисперсионная/полярная/H-связевая совместимость для физических каналов\\
        IDAC $\gamma^\infty$ & коэффициент активности при бесконечном разбавлении, прямой якорь для NRTL\\
        \bottomrule
      \end{tabularx}
      \vspace{0.35em}
      \smallnote{Основная цель остаётся $\ln x_2$: вспомогательные метки разрежены. IDAC — не растворимость, а предел активности при $x_2\to0$.}
    \end{column}
    \begin{column}{0.38\textwidth}
      \begin{block}{Вода сохранена в корпусе}
        \scriptsize
        SMILES \texttt{O}: один тяжёлый атом. Раньше общий фильтр \texttt{min\_atoms=2} удалял такие строки; теперь вода включена как важный растворитель: 6\,524 размеченных строки.
      \end{block}
      \begin{block}{Протокол оценки}
        \scriptsize
        Основное разбиение: scaffold 80/10/10. Дополнительно проверяются отложенные вещества и отложенные растворители.
      \end{block}
    \end{column}
  \end{columns}
  \note{Здесь нужно закрыть два вопроса заранее: что такое IDAC и почему вода есть в обучении, несмотря на одноатомный граф. Подробности тестов воды можно дать при вопросах.}
\end{frame}

\begin{frame}{Параметры Хансена: три типа совместимости}
  \small
  \begin{columns}[T,onlytextwidth]
    \begin{column}{0.46\textwidth}
      \begin{block}{Идея}
        \scriptsize
        Растворитель и вещество лучше совместимы, если близки по трём координатам:
        \[
          \boldsymbol{\delta}=(\delta_d,\delta_p,\delta_h).
        \]
      \end{block}
      \begin{itemize}\tightitems
        \item $\delta_d$ — дисперсионные, ван-дер-ваальсовы контакты.
        \item $\delta_p$ — полярные и электростатические контакты.
        \item $\delta_h$ — способность к водородным связям.
      \end{itemize}
      \[
        R_a^2=4(\Delta\delta_d)^2+(\Delta\delta_p)^2+(\Delta\delta_h)^2.
      \]
      \smallnote{Малое $R_a$ — похожие взаимодействия; большое $R_a$ — низкая совместимость.}
    \end{column}
    \begin{column}{0.50\textwidth}
      \centering
      \begin{tikzpicture}
        \draw[->,TGSlate!65] (0,0)--(4.2,0) node[right,font=\scriptsize] {$\delta_d$};
        \draw[->,TGSlate!65] (0,0)--(0,3.2) node[above,font=\scriptsize] {$\delta_p$};
        \draw[->,TGSlate!65] (0,0)--(-1.2,-1.1) node[left,font=\scriptsize] {$\delta_h$};
        \node[circle,draw=TGOrange,fill=TGOrange!18,minimum size=0.42cm,label={[font=\scriptsize]right:вещество}] (api) at (2.4,1.7) {};
        \node[circle,draw=TGBlue,fill=TGBlue!18,minimum size=0.42cm,label={[font=\scriptsize]right:этанол}] at (2.0,1.35) {};
        \node[circle,draw=TGGreen,fill=TGGreen!18,minimum size=0.42cm,label={[font=\scriptsize]right:вода}] at (1.1,2.55) {};
        \node[circle,draw=TGPurple,fill=TGPurple!15,minimum size=0.42cm,label={[font=\scriptsize]right:гексан}] at (3.45,0.25) {};
        \draw[TGOrange!60,dashed] (api) circle (0.72);
        \node[draw=TGBlue!30,fill=TGBlue!7,rounded corners=7pt,align=left,text width=5.0cm,font=\scriptsize] at (2.1,-1.25)
          {Hansen-метки разрежены, но полезны: они дают мягкую подсказку, какой скрытый канал должен отвечать за дисперсию, а какой — за полярность и H-связи.};
      \end{tikzpicture}
    \end{column}
  \end{columns}
  \note{Hansen — не декоративный приём и не основная цель. Это геометрия совместимости: близость в пространстве типов взаимодействий.}
\end{frame}

% 6
\begin{frame}{Цель и задачи}
  \small
  \begin{block}{Цель}
    Разработать и верифицировать физически-информированную графовую нейронную сеть TGNN-Solv для предсказания равновесной растворимости твёрдых веществ в жидкостях.
  \end{block}
  \vspace{0.2em}
  \begin{columns}[T,onlytextwidth]
    \begin{column}{0.48\textwidth}
      \begin{enumerate}\tightitems
        \item Собрать и обработать доступные выборки.
        \item Провести термодинамическое моделирование SLE.
        \item Выбрать представление молекулы и архитектуру сети.
        \item Настроить обучение и регуляризацию.
      \end{enumerate}
    \end{column}
    \begin{column}{0.48\textwidth}
      \begin{enumerate}\setcounter{enumi}{4}\tightitems
        \item Сравнить с существующими аналогами.
        \item Выявить слабые места модели.
        \item Обновить модель: предобучение, дескрипторы, физические каналы.
      \end{enumerate}
    \end{column}
  \end{columns}
  \note{Это новая карта доклада. Дальше каждый блок отвечает одной из задач, а не просто пересказывает детали реализации.}
\end{frame}

% A2
\begin{frame}{Молекула — это не вектор чисел, а граф}
  \small
  \begin{columns}[c,onlytextwidth]
    \begin{column}{0.34\textwidth}
      Обычная нейросеть ожидает фиксированный вход:
      \[
        [\text{признак}_1,\text{признак}_2,\ldots,\text{признак}_n].
      \]
      \pause
      Но молекулы имеют разное число атомов, разную топологию связей, а перестановка атомов не меняет вещество.
      \vspace{0.25em}
      \begin{enumerate}\tightitems
        \item RDKit-дескрипторы: фиксированный набор чисел, но часть структуры теряется.
        \item GNN: работаем с молекулой как с графом.
      \end{enumerate}
    \end{column}
    \begin{column}{0.63\textwidth}
      \begin{center}
      \resizebox{\linewidth}{!}{%
      \begin{tikzpicture}
        \node[draw=TGLight,fill=white,rounded corners=8pt,minimum width=2.30cm,minimum height=2.65cm,align=center] (w) at (0,0)
          {\molasset[width=2.05cm]{water_plain.png}\\[-0.10em]{\scriptsize Вода}};
        \node[draw=TGLight,fill=white,rounded corners=8pt,minimum width=2.80cm,minimum height=2.65cm,align=center] (b) at (3.05,0)
          {\molasset[width=2.55cm]{benzene_plain.png}\\[-0.10em]{\scriptsize Бензол}};
        \node[draw=TGLight,fill=white,rounded corners=8pt,minimum width=3.90cm,minimum height=2.65cm,align=center] (p) at (6.95,0)
          {\molasset[width=3.55cm]{paracetamol_plain.png}\\[-0.10em]{\scriptsize Парацетамол}};
      \end{tikzpicture}}
      \end{center}
      \vspace{0.15em}
      \smallnote{$N_{\text{атомов}}$ меняется от 1 тяжёлого атома у воды до десятков атомов у лекарственного кандидата; порядок атомов в SMILES не должен менять ответ.}
    \end{column}
  \end{columns}
  \note{Дескрипторы — как рост и вес человека; граф — ближе к фотографии структуры. GNN пытается читать саму структуру.}
\end{frame}

% A3
\begin{frame}{Молекула = граф: вершины — атомы, рёбра — связи}
  \small
  \begin{columns}[c,onlytextwidth]
    \begin{column}{0.58\textwidth}
      \molasset[width=0.93\linewidth]{paracetamol_graph.png}
    \end{column}
    \begin{column}{0.38\textwidth}
      \centering
      \molasset[width=0.66\linewidth]{paracetamol.png}
      \vspace{0.15em}
      \begin{block}{Признаки атомов: 35 чисел}
        \scriptsize
        атомный номер, гибридизация, формальный заряд, ароматичность, электроотрицательность, радиус Ван-дер-Ваальса, поляризуемость, \ldots
      \end{block}
      \begin{block}{Признаки связей: 8 чисел}
        \scriptsize
        тип связи, сопряжённость, кольцо, стерео и другие признаки химической связи.
      \end{block}
      \smallnote{В датасете используется граф тяжёлых атомов: водороды обычно неявные, но их вклад кодируется в признаках атома.}
    \end{column}
  \end{columns}
  \note{Это паспорт для каждого атома и связи. GNN читает эти паспорта и передаёт информацию между соседями.}
\end{frame}

% A4
\begin{frame}{Передача сообщений: атомы обмениваются информацией}
  \small
  \begin{columns}[T,onlytextwidth]
    \begin{column}{0.48\textwidth}
      \pause
      \[
        \mathbf{m}_{ij}=\mathrm{MLP}(\mathbf{h}_i,\mathbf{h}_j,\mathbf{e}_{ij})
      \]
      \smallnote{Атом $j$ рассказывает атому $i$ о себе с учётом связи.}
      \pause
      \[
        \mathbf{a}_i=\sum_{j\in\mathcal{N}(i)}\mathbf{m}_{ij}
      \]
      \smallnote{Атом $i$ суммирует сообщения соседей.}
      \pause
      \[
        \mathbf{h}_i\leftarrow
        \mathrm{LayerNorm}\!\left(
        \mathbf{h}_i+\mathrm{MLP}(\mathbf{h}_i,\mathbf{a}_i)
        \right).
      \]
      \smallnote{Повторяем $L=6$ раз: атом собирает контекст в радиусе 6 связей.}
    \end{column}
    \begin{column}{0.48\textwidth}
      \fitfig{
      \begin{tikzpicture}[>=Latex]
        \node[circle,draw=TGBlue,fill=TGBlue!10,minimum size=0.62cm] (c1) at (0,0) {C};
        \node[circle,draw=TGOrange,fill=TGOrange!12,minimum size=0.62cm] (o) at (1.4,0.65) {O};
        \node[circle,draw=TGGreen,fill=TGGreen!10,minimum size=0.62cm] (n) at (1.4,-0.65) {N};
        \node[circle,draw=TGBlue,fill=TGBlue!10,minimum size=0.62cm] (c2) at (2.8,0) {C};
        \draw[thick,TGSlate] (c1)--(o)--(c2)--(n)--(c1);
        \pause
        \draw[->,very thick,TGOrange] (o) to[bend left=12] (c2);
        \draw[->,very thick,TGGreen] (n) to[bend right=12] (c2);
        \pause
        \node[circle,draw=TGBlueDark,fill=TGBlue!28,minimum size=0.74cm] at (2.8,0) {C};
        \node[font=\scriptsize,text=TGSlate] at (1.4,-1.35) {после обновления представление меняется};
      \end{tikzpicture}}
    \end{column}
  \end{columns}
  \note{Каждый атом сначала знает только себя, потом соседей, потом соседей соседей. После нескольких раундов появляется представление молекулярного окружения.}
\end{frame}

% A5
\begin{frame}{Как получить один вектор для всей молекулы?}
  \small
  После передачи сообщений у каждого атома есть представление $\mathbf{h}_i\in\mathbb{R}^d$. Нужно получить один вектор молекулы.
  \begin{columns}[T,onlytextwidth]
    \begin{column}{0.48\textwidth}
      Простая агрегация:
      \[
        \mathbf{g}_{mol}=\frac{1}{N}\sum_i\mathbf{h}_i.
      \]
      \pause
      Агрегация с вниманием:
      \[
        \alpha_i=\mathrm{softmax}(\mathbf{w}^{\top}\mathbf{h}_i),
        \qquad
        \mathbf{g}_{att}=\sum_i\alpha_i\mathbf{h}_i.
      \]
      \[
        \mathbf{g}_{mol}=[\mathbf{g}_{att}\,\|\,\mathbf{g}_{s2s}].
      \]
      \begin{block}{Что такое Set2Set}
        \scriptsize
        $\mathbf{g}_{s2s}$ — многошаговая агрегация множества атомов: маленький LSTM несколько раз ``переспрашивает'' атомы через внимание и собирает более богатый вектор молекулы.
      \end{block}
    \end{column}
    \begin{column}{0.48\textwidth}
      \fitfig{
      \begin{tikzpicture}
        \foreach \x/\lab/\col/\r in {0/O/TGOrange/0.36,1.1/N/TGGreen/0.30,2.2/C/TGBlue/0.18,3.3/C/TGBlue/0.08}{
          \node[circle,draw=\col,fill=\col!12,minimum size=0.72cm] at (\x,0) {\lab};
          \draw[\col,very thick] (\x,-0.55)--(\x,-0.55-\r*3);
        }
        \node[draw=TGBlue,fill=TGBlue!10,rounded corners=5pt,align=center] at (1.65,-2.2) {взвешенный\\вектор молекулы};
      \end{tikzpicture}}
      \smallnote{OH/NH-фрагменты могут получить больший вес, чем инертные CH-группы.}
    \end{column}
  \end{columns}
  \note{Агрегация с вниманием — способ дать больший вес атомам, которые важны для ответа.}
\end{frame}

% A6
\begin{frame}{Дескрипторы: 200 чисел, описывающих молекулу}
  \small
  \begin{columns}[T,onlytextwidth]
    \begin{column}{0.50\textwidth}
      \begin{tabular}{ll}
        \toprule
        \textbf{Дескриптор} & \textbf{Пример для парацетамола}\\
        \midrule
        MolLogP & 0.49\\
        TPSA & 49.3 \AA$^2$\\
        MolWt & 151.2\\
        NumHDonors & 2\\
        NumHAcceptors & 2\\
        FractionCSP3 & 0.125\\
        RingCount & 1\\
        \bottomrule
      \end{tabular}
    \end{column}
    \begin{column}{0.44\textwidth}
      \begin{block}{Плюсы}
        \scriptsize
        Быстро вычисляются, химически интерпретируемы, подходят для простых табличных моделей.
      \end{block}
      \begin{block}{Минусы}
        \scriptsize
        Фиксированный набор нельзя дообучить; часть топологии и контекст взаимодействия с растворителем теряются.
      \end{block}
    \end{column}
  \end{columns}
  \note{Дескрипторы — анкета из 200 вопросов. GNN — собеседование, где признаки адаптируются под задачу.}
\end{frame}

\section{Физический фундамент}

% 6
\begin{frame}{Химический потенциал}
  \small
  \begin{block}{Отправная точка}
    \[
      \mu_i\equiv\left(\frac{\partial G}{\partial n_i}\right)_{T,P,n_{j\neq i}}
    \]
  \end{block}
  Для жидкой смеси:
  \[
    \mu_i^{\mathrm{liq}}=\mu_i^{\circ,\mathrm{liq}}(T,P)+RT\ln(\gamma_i x_i).
  \]
  Для чистого кристалла:
  \[
    \mu_2^{\mathrm{cr}}=\mu_2^{\circ,\mathrm{cr}}(T,P).
  \]
  \smallnote{$\gamma_i=1$ — идеальная смесь; $\gamma_i\gg1$ — компоненту ``некомфортно'' в окружении.}
  \note{Объяснить химический потенциал как стремление компонента покинуть фазу.}
\end{frame}

% 7
\begin{frame}{Условие равновесия $\to$ мастер-уравнение}
  \small
  Равновесие кристалл $\leftrightarrow$ раствор:
  \[
    \mu_2^{\mathrm{cr}}=\mu_2^{\mathrm{liq}}.
  \]
  \pause
  Подстановка:
  \[
    \mu_2^{\circ,\mathrm{cr}}=\mu_2^{\circ,\mathrm{liq}}+RT\ln(\gamma_2 x_2).
  \]
  \pause
  Вводим $\Delta G_{\mathrm{fus}}\equiv\mu_2^{\circ,\mathrm{liq}}-\mu_2^{\circ,\mathrm{cr}}$:
  \[
    \ln(\gamma_2 x_2)=-\frac{\Delta G_{\mathrm{fus}}}{RT}.
  \]
  \pause
  \[
    \boxed{\ln x_2=-\frac{\Delta G_{\mathrm{fus}}(T)}{RT}-\ln\gamma_2}
  \]
  \smallnote{Точное уравнение при чистой кристаллической фазе и $P\approx1$ атм.}
  \note{Показать, что physics bottleneck не произволен: он следует из равновесия.}
\end{frame}

% 8
\begin{frame}{Структурная декомпозиция}
  \[
    \ln x_2=
    \underbrace{-\Phi(T;T_m,\Delta H_{\mathrm{fus}})}_{\text{кристаллический штраф}}
    +
    \underbrace{(-\ln\gamma_2)}_{\text{комфорт в растворе}}.
  \]
  \begin{columns}[T,onlytextwidth]
    \begin{column}{0.48\textwidth}
      \begin{block}{Кристалл}
        \begin{itemize}\tightitems
          \item зависит только от растворяемого вещества;
          \item параметры: $T_m$, $\Delta H_{\mathrm{fus}}$;
          \item вопрос: сколько энергии нужно, чтобы разрушить кристалл?
        \end{itemize}
      \end{block}
    \end{column}
    \begin{column}{0.48\textwidth}
      \begin{block}{Раствор}
        \begin{itemize}\tightitems
          \item зависит от пары вещество–растворитель;
          \item параметры: $\tau_{12}$, $\tau_{21}$, $\alpha$;
          \item зависит от состава $x_2$: $\gamma_2$ зависит от $x_2$, а $x_2$ зависит от $\gamma_2$.
        \end{itemize}
      \end{block}
    \end{column}
  \end{columns}
  \note{Переход: слева нужна кристаллическая голова, справа — парное взаимодействие и NRTL-голова.}
\end{frame}

% 9
\begin{frame}{Свободная энергия плавления}
  \small
  В точке плавления:
  \[
    \Delta G_{\mathrm{fus}}(T_m)=0
    \quad\Rightarrow\quad
    \Delta S_{\mathrm{fus}}(T_m)=\frac{\Delta H_{\mathrm{fus}}(T_m)}{T_m}.
  \]
  \pause
  Интегрирование Кирхгофа от $T_m$ до $T$:
  \[
    \Delta G_{\mathrm{fus}}(T)=
    \Delta H_{\mathrm{fus}}\!\left(1-\frac{T}{T_m}\right)
    +\Delta C_p\!\left[(T-T_m)-T\ln\frac{T}{T_m}\right].
  \]
  \pause
  При $\Delta C_p=0$:
  \[
    \boxed{\Phi(T)=\frac{\Delta H_{\mathrm{fus}}}{R}\left(\frac{1}{T}-\frac{1}{T_m}\right)}
  \]
  \pause
  \smallnote{Цена приближения: для парацетамола при 298 K ошибка порядка $0.85$ в $\ln x_2$, сопоставимо с MAE.}
  \note{Подчеркнуть: Hildebrand упрощает решатель, но создаёт систематическую ошибку, которую голова может компенсировать.}
\end{frame}

% Numeric example: ideal solubility
\begin{frame}{Пример: парацетамол}
  \small
  Дано:
  \[
    T_m=442\ \mathrm{K},\quad
    \Delta H_{\mathrm{fus}}=26\,400\ \mathrm{J/mol},\quad
    T=298.15\ \mathrm{K},\quad
    R=8.314.
  \]
  \pause
  \[
    \Phi
    =
    \frac{26\,400}{8.314}
    \left(\frac{1}{298.15}-\frac{1}{442}\right)
  \]
  \pause
  \[
    =
    3174\,(0.003354-0.002262)
    =
    3174\cdot0.001092
    =
    3.46.
  \]
  \pause
  \[
    x_2^{\mathrm{ideal}}=\exp(-3.46)=0.031.
  \]
  \pause
  \keybox{Без учёта взаимодействий: парацетамол растворяется на 3.1 mol\%. Это идеальная растворимость.}
  \note{Этот пример нужен, чтобы аудитория увидела масштаб: один кристаллический штраф уже даёт численную оценку, но она ещё не учитывает комфорт молекулы в растворителе.}
\end{frame}

% 10
\begin{frame}{NRTL: гипотеза локального состава}
  \small
  Renon–Prausnitz:
  \[
    \frac{x_{12}}{x_{22}}=\frac{x_1}{x_2}\exp(-\alpha_{12}\tau_{12}).
  \]
  \pause
  \begin{block}{Как читать эту формулу}
    \scriptsize
    $x_{12}/x_{22}$ — локальное отношение растворитель/вещество вокруг молекулы компонента 2. Оно не обязано совпадать с объёмным $x_1/x_2$, потому что молекулы притягиваются и отталкиваются не одинаково. $\tau_{12}$ и $\tau_{21}$ — две направленные энергии контакта; $\alpha$ задаёт, насколько сильно локальное окружение отклоняется от случайного.
  \end{block}
  \pause
  \[
    \tau_{ij}=\frac{\Delta g_{ij}}{RT},
    \qquad
    G_{ij}=\exp(-\alpha\tau_{ij}).
  \]
  \pause
  Избыточная энергия Гиббса $G^E$ — добавка к энергии идеальной смеси:
  \[
    \frac{G^E}{RT}
    =x_1x_2
    \left(
      \frac{\tau_{21}G_{21}}{x_1+x_2G_{21}}
      +
      \frac{\tau_{12}G_{12}}{x_2+x_1G_{12}}
    \right).
  \]
  \smallnote{Физика: локальное окружение молекулы отличается от объёмного состава из-за энергий контакта.}
  \note{Показать, что NRTL превращает парное взаимодействие в термодинамически согласованную активность.}
\end{frame}

% 11
\begin{frame}{Коэффициент активности}
  \small
  Дифференцируем $G^E$:
  \[
    \boxed{
    \ln\gamma_2=x_1^2\left[
      \tau_{12}\left(\frac{G_{12}}{x_2+x_1G_{12}}\right)^2
      +
      \frac{\tau_{21}G_{21}}{(x_1+x_2G_{21})^2}
    \right]}
  \]
  \pause
  \begin{columns}[T,onlytextwidth]
    \begin{column}{0.49\textwidth}
      \[
        x_2\to1:\quad \ln\gamma_2\to0
      \]
      \smallnote{Раулев предел.}
    \end{column}
    \begin{column}{0.49\textwidth}
      \[
        x_2\to0:\quad
        \ln\gamma_2^\infty=
        \tau_{12}+\tau_{21}e^{-\alpha\tau_{21}}
      \]
      \smallnote{IDAC-супервизия.}
    \end{column}
  \end{columns}
  \pause
  \begin{center}
    \small\color{TGSlate}
    Если $\gamma_2>1$, растворяемому веществу энергетически хуже, чем в идеальной смеси. Поэтому в SLE-уравнении член $-\ln\gamma_2$ уменьшает растворимость.
  \end{center}
  \pause
  \begin{block}{Зачем здесь Гиббс–Дюгем}
    \[
      x_1\,d\ln\gamma_1+x_2\,d\ln\gamma_2=0
      \qquad (T,P=\mathrm{const}).
    \]
    \smallnote{Это проверка термодинамической согласованности: активности компонентов одной смеси не могут меняться независимо. В NRTL условие выполняется автоматически, потому что обе активности получены из одного $G^E$, а не предсказаны разными нейросетевыми головами.}
  \end{block}
  \note{Переход: теперь уравнение SLE стало трансцендентным, потому что $\gamma_2$ зависит от $x_2$.}
\end{frame}

% Numeric example: nonideality
\begin{frame}{Пример: парацетамол в этаноле}
  \small
  Типичные параметры для иллюстрации:
  \[
    \tau_{12}\approx1.2,\qquad
    \tau_{21}\approx0.8,\qquad
    \alpha=0.30.
  \]
  \pause
  В пределе бесконечного разбавления:
  \[
    \ln\gamma_2^\infty
    =
    \tau_{12}+\tau_{21}\exp(-\alpha\tau_{21})
  \]
  \pause
  \[
    =
    1.2+0.8\exp(-0.3\cdot0.8)
    =
    1.2+0.8\cdot0.787
    =
    1.83.
  \]
  \pause
  \[
    \gamma_2^\infty=\exp(1.83)=6.2.
  \]
  \pause
  \[
    \ln x_2\approx-\Phi-\ln\gamma_2^\infty
    =
    -3.46-1.83
    =
    -5.29,
    \qquad
    x_2\approx\exp(-5.29)=0.005.
  \]
  \pause
  \keybox{В этаноле парацетамолу ``некомфортно'': $\gamma\approx6.2$. Растворимость падает с идеальных 3.1 mol\% до примерно 0.5 mol\%.}
  \note{Подчеркнуть: это иллюстративный расчёт с приближёнными NRTL-параметрами; главное — смысл активности как штрафа за неидеальность.}
\end{frame}

% 12
\begin{frame}{Задача о неподвижной точке}
  \small
  \begin{columns}[T,onlytextwidth]
    \begin{column}{0.47\textwidth}
      \[
        x_2=\frac{\exp(-\Phi)}{\gamma_2(x_2)}\equiv g(x_2).
      \]
      \pause
      \[
        x_2^{(k+1)}=\lambda g(x_2^{(k)})+(1-\lambda)x_2^{(k)}.
      \]
      \pause
      \[
        |g'(x_2^*)|=x_2^*|\eta|<1.
      \]
      \vspace{0.15em}
      \begin{block}{Почему обычно сходится}
        \scriptsize
        $g(x_2)$ — один шаг решателя; $\lambda\in(0,1]$ — демпфирование; $\epsilon$ — порог остановки; $\eta=\partial\ln\gamma_2/\partial x_2$.
        При $x_2^*\ll1$ множитель $x_2^*$ подавляет наклон $\eta$.
      \end{block}
    \end{column}
    \begin{column}{0.51\textwidth}
      \fitfig{
      \begin{tikzpicture}[>=Latex]
        \node[draw,fill=TGBlue!10,rounded corners] (start) at (0,1.5) {$x_2^{(k)}$};
        \node[draw,fill=TGRed!10,rounded corners] (nrtl) at (2.4,1.5) {NRTL $\gamma_2$};
        \node[draw,fill=TGGreen!10,rounded corners] (sle) at (2.4,0) {$g(x_2)$};
        \node[draw,fill=TGOrange!10,rounded corners] (stop) at (0,0) {$|x^{k+1}-x^k|<\epsilon$};
        \draw[->,thick] (start)--(nrtl);
        \draw[->,thick] (nrtl)--(sle);
        \draw[->,thick] (sle)--(stop);
        \draw[->,thick] (stop.west) .. controls (-1.2,0.5) and (-1.2,1.1) .. (start.west);
      \end{tikzpicture}}
    \end{column}
  \end{columns}
  \note{Объяснить демпфирование как инженерную защиту: берём смесь старого и нового значения, а не весь шаг сразу.}
\end{frame}

\begin{frame}{Решатель в действии: 5 итераций для парацетамола}
  \scriptsize
  \centering
  \begin{block}{Что делает одна итерация}
    \scriptsize
    Берём текущую оценку $x_2^{(k)}$, считаем по NRTL $\gamma_2(x_2^{(k)})$ и обновляем
    $x_2^{new}=\exp(-\Phi)/\gamma_2$. Старт — идеальная растворимость $x_2^{ideal}=\exp(-\Phi)$.
  \end{block}
  \vspace{-0.2em}
  \setlength{\tabcolsep}{4.5pt}
  \renewcommand{\arraystretch}{1.18}
  \begin{tabular}{r c c c c}
    \toprule
    \textbf{$k$} & \textbf{$x_2^{(k)}$} & \textbf{$\gamma_2$} &
    \textbf{\shortstack{$x_2^{new}$\\$=\exp(-\Phi)/\gamma_2$}} &
    \textbf{\shortstack{невязка\\$|x^{new}-x^{(k)}|$}}\\
    \midrule
    \uncover<1->{0} & \uncover<1->{0.03143} & \uncover<1->{—} &
      \uncover<1->{$x_2^{ideal}$} & \uncover<1->{—}\\
    \uncover<2->{1} & \uncover<2->{0.03143} & \uncover<2->{5.44} &
      \uncover<2->{0.00578} & \uncover<2->{0.02565}\\
    \uncover<3->{2} & \uncover<3->{0.00578} & \uncover<3->{6.07} &
      \uncover<3->{0.00518} & \uncover<3->{0.00060}\\
    \uncover<4->{3} & \uncover<4->{0.00518} & \uncover<4->{6.09} &
      \uncover<4->{0.00516} & \uncover<4->{$1.36\cdot10^{-5}$}\\
    \uncover<5->{4} & \uncover<5->{0.00516} & \uncover<5->{6.09} &
      \uncover<5->{0.00516} & \uncover<5->{$3.08\cdot10^{-7}$}\\
    \uncover<6->{5} & \uncover<6->{0.00516} & \uncover<6->{6.09} &
      \uncover<6->{0.00516} & \uncover<6->{$6.98\cdot10^{-9}$}\\
    \bottomrule
  \end{tabular}
  \renewcommand{\arraystretch}{1.0}
  \uncover<7->{
  \[
    \boxed{x_2^*=0.00516,\qquad \ln x_2^*=-5.27.}
  \]
  \vspace{-0.3em}
  \keybox{Сошлось за 5 итераций. Первое приближение — идеальный раствор — было примерно в 6 раз больше решения SLE/NRTL для выбранных параметров.}
  }
  \note{Числа рассчитаны по формуле NRTL из NRTLLayer: Phi=3.46, tau12=1.2, tau21=0.8, alpha=0.3. Важно показать, что решатель неподвижной точки быстро уходит от идеальной растворимости к неидеальной.}
\end{frame}

% 13
\begin{frame}{Неявное дифференцирование через решатель}
  \small
  Проблема: развёрнутое обратное распространение через $N$ итераций даёт затухание порядка $|g'|^N$.
  \pause
  \[
    F(x_2^*,\theta)=x_2^*-g(x_2^*;\theta)=0.
  \]
  \pause
  По теореме о неявной функции:
  \[
    \frac{dx_2^*}{d\theta_j}
    =
    -\left(\frac{\partial F}{\partial x_2^*}\right)^{-1}
    \frac{\partial F}{\partial\theta_j}
    =
    -\frac{\partial g/\partial\theta_j}{1-g'(x_2^*)}.
  \]
  \pause
  \[
    \boxed{
    \frac{d\ln x_2^*}{d\theta_j}
    =
    -\frac{\partial(\Phi+\ln\gamma_2)/\partial\theta_j}{1+x_2^*\eta}}
  \]
  \smallnote{$O(1)$ памяти, $O(1)$ вычислений NRTL, градиент равновесного уравнения вместо развёрнутой RNN.}
  \note{Сформулировать, что решатель остаётся физическим, но обучаемая верхняя часть модели получает точные градиенты.}
\end{frame}

\begin{frame}{Где мы сейчас?}
  \centering
  \vspace{0.4cm}
  \roadmapbar{2}
  \vspace{1.0cm}

  {\fontsize{20}{24}\selectfont
  Мы вывели мастер-уравнение SLE и показали, что нужно предсказать
  физические параметры.\par}
  \vspace{0.7cm}
  {\fontsize{22}{26}\selectfont\bfseries\color{TGBlueDark}
  Физика говорит, \textbf{что} предсказывать.\\
  Архитектура — \textbf{как}.\par}
  \note{Переход: математически понятно, какие величины нужны; теперь построим сеть, которая выдаёт эти величины из молекулярных графов.}
\end{frame}

\begin{frame}{Контрольное сравнение: где именно проверяется физика}
  \small
  \fitfig{
  \begin{tikzpicture}[>=Latex]
    \tikzstyle{card}=[draw,rounded corners=12pt,minimum height=0.82cm,align=center]
    \node[anchor=west,text=TGSlate] at (0,3.35) {Два способа использовать один и тот же молекулярный энкодер};

    \node[card,fill=TGSlate!7,draw=TGSlate!20,minimum width=2.6cm] (d0) at (1.2,2.25) {графы};
    \node[card,fill=TGBlue!9,draw=TGBlue!22,minimum width=2.8cm] (d1) at (4.3,2.25) {общий GNN};
    \node[card,fill=TGOrange!10,draw=TGOrange!25,minimum width=2.8cm] (d2) at (7.5,2.25) {прямой MLP};
    \node[card,fill=TGGreen!10,draw=TGGreen!25,minimum width=2.0cm] (d3) at (10.4,2.25) {$\ln x_2$};
    \draw[->,thick,TGSlate] (d0)--(d1)--(d2)--(d3);
    \node[anchor=east,text=TGSlate,font=\bfseries] at (-0.15,2.25) {DirectGNN};

    \node[card,fill=TGSlate!7,draw=TGSlate!20,minimum width=2.6cm] (t0) at (1.2,0.85) {графы};
    \node[card,fill=TGBlue!9,draw=TGBlue!22,minimum width=2.8cm] (t1) at (4.3,0.85) {общий GNN};
    \node[card,fill=TGGreen!10,draw=TGGreen!25,minimum width=3.2cm] (t2) at (7.5,0.85) {$T_m,\Delta H,\tau,\alpha$};
    \node[card,fill=TGRed!9,draw=TGRed!25,minimum width=2.2cm] (t3) at (10.4,0.85) {SLE};
    \node[card,fill=TGGreen!10,draw=TGGreen!25,minimum width=2.0cm] (t4) at (12.8,0.85) {$\ln x_2$};
    \draw[->,thick,TGSlate] (t0)--(t1)--(t2)--(t3)--(t4);
    \node[anchor=east,text=TGBlueDark,font=\bfseries] at (-0.15,0.85) {TGNN-Solv};

    \node[draw=TGBlue!20,fill=white,rounded corners=10pt,align=left,text width=11.4cm,anchor=west] at (0,-0.55)
      {\textbf{Проверка честная:} общая графовая основа, данные и бюджет совпадают. Бюджет здесь — число эпох и вычислительных ресурсов обучения. Разница — только после энкодера.};
  \end{tikzpicture}}
  \note{Этот слайд переносит контрольный вопрос в архитектурный блок: физика проверяется сравнением с прямой моделью при одинаковом backbone.}
\end{frame}

\section{Построение архитектуры}

% 14
\begin{frame}{Что нужно предсказать}
  \small
  \begin{tabularx}{\textwidth}{lYY}
    \toprule
    \textbf{Параметр} & \textbf{Смысл} & \textbf{Зависит от}\\
    \midrule
    $T_m$ & температура плавления & только вещество\\
    $\Delta H_{\mathrm{fus}}$ & энтальпия плавления & только вещество\\
    $\tau_{12}(T)$ & энергия контакта 1–2 & пара + $T$\\
    $\tau_{21}(T)$ & энергия контакта 2–1 & пара + $T$\\
    $\alpha_{12}$ & неслучайность локального состава & пара\\
    \bottomrule
  \end{tabularx}
  \vspace{0.4em}
  \keybox{Архитектура должна отделить чистые свойства вещества от взаимодействия пары и температурной NRTL-головы.}
  \note{Переход: строим архитектуру снизу вверх — сначала граф, затем кристаллические головы, затем взаимодействие.}
\end{frame}

% 15
\begin{frame}{Шаг 1: графовый энкодер}
  \small
  \begin{columns}[c,onlytextwidth]
    \begin{column}{0.64\textwidth}
      \centering
      \begin{minipage}[c]{0.56\linewidth}
        \centering
        \molasset[width=\linewidth]{paracetamol_graph.png}\\[-0.15em]
        {\scriptsize\color{TGSlate}вещество: парацетамол}
      \end{minipage}\hfill
      \begin{minipage}[c]{0.38\linewidth}
        \centering
        \molasset[width=\linewidth]{dmso_graph.png}\\[-0.15em]
        {\scriptsize\color{TGSlate}растворитель: ДМСО}
      \end{minipage}
      \vspace{0.35em}
      \[
        \mathbf{m}_{ij}^{(\ell)}=\phi(\mathbf{h}_i,\mathbf{h}_j,\mathbf{e}_{ij})
      \]
      \smallnote{Один энкодер читает графы атомов и связей; роль вещества и растворителя задаётся адаптерами после общей основы.}
    \end{column}
    \begin{column}{0.32\textwidth}
      \begin{itemize}\tightitems
        \item \textbf{энкодер} — блок, который превращает граф в атомные/молекулярные векторы;
        \item атомы: 35 признаков;
        \item связи: 8 признаков;
        \item общая основа + адаптеры ролей;
        \item варианты: MPNN (локальные сообщения), GPS/GraphGPS (локальные сообщения + глобальное внимание), физические каналы.
      \end{itemize}
    \end{column}
  \end{columns}
  \smallnote{\textbf{Дизайн-ограничение:} температура не подаётся в энкодер. Так кристаллические головы учат свойства вещества, а не рабочие условия эксперимента.}
  \note{Сказать: изоляция температуры — одно из физических ограничений архитектуры.}
\end{frame}

% 16
\begin{frame}{Шаг 2: кристаллические головы}
  \small
  \smallnote{\textbf{Голова} — небольшой MLP после энкодера: он переводит вектор молекулы в конкретную физическую величину.}
  \vspace{0.2em}
  \[
    \mathbf{g}_{2}^{\mathrm{pre}}
    =
    \mathrm{MeanPool}(\{\mathbf{h}_i^{(2)}\}).
  \]
  \[
    T_m=100+600\cdot\sigma(\mathrm{MLP}(\mathbf{g}_{2}^{\mathrm{pre}}))\in[100,700]~\mathrm{K}.
  \]
  \[
    \Delta H_{\mathrm{fus}}
    =
    S_H\cdot\mathrm{softplus}(\mathrm{MLP}(\mathbf{g}_{2}^{\mathrm{pre}}))>0.
  \]
  \begin{block}{Опциональный GC-приор}
    \[
      T_m^{\mathrm{pred}}=T_m^{\mathrm{GC,cal}}+50\cdot\tanh(h).
    \]
    \scriptsize GC-приор — групповая оценка свойства по функциональным фрагментам молекулы; сеть учит ограниченную поправку к этой оценке.
  \end{block}
  \note{Показать, что ограничения встроены в голову, а не штрафуются постфактум.}
\end{frame}

% B1
\begin{frame}{Joback: оценка $T_m$ по кусочкам молекулы}
  \small
  \begin{columns}[T,onlytextwidth]
    \begin{column}{0.42\textwidth}
      \molasset[width=\linewidth]{paracetamol.png}
      \[
        \boxed{T_m^{GC}=122.5+\sum_k n_k\,\Delta T_{m,k}}
      \]
      \smallnote{Типичная точность GC-груба: MAE порядка 30–50 K, но это лучше случайного старта.}
    \end{column}
    \begin{column}{0.54\textwidth}
      \scriptsize
      \begin{tabular}{lcc}
        \toprule
        \textbf{Группа} & \textbf{Кол-во} & \textbf{Вклад}\\
        \midrule
        ароматич. CH & 4 & $4\times20.0$ K\\
        ароматич. C & 2 & $2\times17.2$ K\\
        фенольный OH & 1 & $54.0$ K\\
        амидный NH & 1 & $32.8$ K\\
        амидный C=O & 1 & $33.4$ K\\
        CH$_3$ & 1 & $23.6$ K\\
        \bottomrule
      \end{tabular}
      \vspace{0.35em}
      \keybox{$T_m^{GC}\approx460$ K; эксперимент для парацетамола около 442 K.}
    \end{column}
  \end{columns}
  \note{Joback похож на оценку дома по комнатам и гаражу: грубо, но даёт разумный ориентир до обучения.}
\end{frame}

% B2
\begin{frame}{Нейросеть исправляет ошибку GC, а не учит с нуля}
  \small
  \begin{columns}[T,onlytextwidth]
    \begin{column}{0.48\textwidth}
      Без GC-приора:
      \[
        T_m=100+600\sigma(h)\in[100,700]\;\mathrm{K}.
      \]
      \pause
      С GC-приором:
      \[
        \boxed{T_m^{pred}=T_m^{GC,cal}+50\tanh(h)}
      \]
      \[
        \Delta H_{fus}^{pred}=\Delta H_{fus}^{GC}\exp(r),
        \qquad r\in[\ln0.3,\ln3.0].
      \]
    \end{column}
    \begin{column}{0.48\textwidth}
      \fitfig{
      \begin{tikzpicture}
        \draw[->] (0,0)--(5.6,0) node[right] {$T_m$};
        \draw[TGSlate!50] (0.4,0.65)--(5.2,0.65);
        \fill[TGOrange!35] (0.4,0.48) rectangle (5.2,0.82);
        \node[font=\scriptsize] at (2.8,1.05) {без приора: 600 K};
        \draw[TGSlate!50] (2.45,-0.65)--(3.25,-0.65);
        \fill[TGBlue!35] (2.45,-0.82) rectangle (3.25,-0.48);
        \draw[very thick,TGBlue] (2.85,-1.02)--(2.85,-0.28);
        \node[font=\scriptsize] at (2.85,-1.25) {GC $\pm50$ K};
      \end{tikzpicture}}
      \smallnote{Пространство поиска для $T_m$ сужается примерно в 12 раз: сеть учит ограниченную поправку.}
    \end{column}
  \end{columns}
  \note{С подсказкой сеть оценивает не абсолютный рост человека, а поправку к разумному диапазону.}
\end{frame}

% 17
\begin{frame}{Шаг 3: взаимодействие вещества и растворителя}
  \small
  \begin{columns}[T,onlytextwidth]
    \begin{column}{0.58\textwidth}
      \molasset[width=\linewidth]{cross_attention_pair.png}
    \end{column}
    \begin{column}{0.38\textwidth}
      \scriptsize
      \begin{block}{Смысл}
        \scriptsize
        До этого момента две молекулы кодировались по отдельности.
        Теперь нужно построить \textbf{представление пары}, а не
        просто положить рядом два независимых вектора.
      \end{block}
      \begin{block}{Что делает перекрёстное внимание}
        \scriptsize
        Каждый атом вещества смотрит на все атомы растворителя
        и выбирает, от кого взять больше информации.
        После этого атомные представления уже учитывают партнёра.
      \end{block}
      \[
        \tilde{\mathbf{h}}_{\text{вещ}}
        =
        \mathrm{CrossAttn}(\mathbf{h}_{\text{вещ}},\mathbf{h}_{\text{раств}})
      \]
    \end{column}
  \end{columns}
  \smallnote{Именно здесь из двух отдельных молекул появляется признак взаимодействия. Дальше из него строится $\mathbf g_{\mathrm{pair}}$, который идёт в NRTL-голову.}
  \note{Этот слайд должен ответить на простой вопрос: зачем вообще нужен отдельный блок взаимодействия. Ответ: растворимость определяется не только свойствами молекул по отдельности, но и тем, как они контактируют друг с другом.}
\end{frame}

% C1
\begin{frame}{Перекрёстное внимание: как оно работает}
  \footnotesize
  \begin{columns}[T,onlytextwidth]
    \begin{column}{0.50\textwidth}
      \footnotesize
      Для атома вещества $i$ и атома растворителя $j$:
      \[
        \mathbf{q}_i=W_Q\mathbf{h}_i,\quad
        \mathbf{k}_j=W_K\mathbf{h}_j,\quad
        \mathbf{v}_j=W_V\mathbf{h}_j.
      \]
      \pause
      \[
        s_{ij}=\frac{\mathbf{q}_i^{\top}\mathbf{k}_j}{\sqrt{d}},
        \qquad
        \alpha_{ij}=\mathrm{softmax}_j(s_{ij}).
      \]
      \pause
      \[
        \tilde{\mathbf{h}}_i=\sum_j\alpha_{ij}\mathbf{v}_j.
      \]
      \vspace{-0.2em}
      {\scriptsize
      $i$ — атом вещества; $j$ — атом растворителя; $d$ — размерность внимания; $\alpha_{ij}$ — вес контакта; $\sum_j\alpha_{ij}=1$.}
      \smallnote{\textbf{Запрос} — что ищет атом; \textbf{ключ} — как другой атом предъявляет себя для сравнения; \textbf{значение} — что передаётся.}
    \end{column}
    \begin{column}{0.46\textwidth}
      \fitfig{
      \begin{tikzpicture}[>=Latex]
        \foreach \y/\lab/\col in {1.5/O/TGOrange,0/N/TGGreen,-1.5/C/TGBlue}
          \node[circle,draw=\col,fill=\col!12,minimum size=0.58cm] (s\lab) at (0,\y) {\lab};
        \node[circle,draw=TGOrange,fill=TGOrange!12,minimum size=0.58cm] (vO) at (4.2,1.8) {O};
        \node[circle,draw=TGGreen,fill=TGGreen!12,minimum size=0.58cm] (vS) at (4.2,0.6) {S};
        \node[circle,draw=TGBlue,fill=TGBlue!12,minimum size=0.58cm] (vC1) at (4.2,-0.6) {C};
        \node[circle,draw=TGBlue,fill=TGBlue!12,minimum size=0.58cm] (vC2) at (4.2,-1.8) {C};
        \draw[->,line width=2.8pt,TGOrange] (sO)--(vO);
        \draw[->,line width=1.8pt,TGGreen] (sN)--(vS);
        \draw[->,line width=0.8pt,TGLight] (sC)--(vC2);
        \draw[->,line width=0.7pt,TGLight] (sO)--(vC1);
        \node[font=\scriptsize,text=TGSlate] at (2.1,-2.45) {толщина стрелки $\sim\alpha_{ij}$};
      \end{tikzpicture}}
    \end{column}
  \end{columns}
  \note{Этот слайд должен закрыть и формальный, и интуитивный вопрос сразу: как из атомов двух молекул получается таблица контактов и почему веса $\alpha_{ij}$ можно интерпретировать как важность контакта. Здесь же стоит проговорить, что внимание двунаправленное: вещество смотрит на растворитель, и растворитель смотрит на вещество.}
\end{frame}

\begin{frame}{Как читать карту внимания}
  \small
  \begin{columns}[T,onlytextwidth]
    \begin{column}{0.52\textwidth}
      \molasset[width=\linewidth]{cross_attention_pair.png}
    \end{column}
    \begin{column}{0.44\textwidth}
      \begin{block}{Строка и столбец}
        \scriptsize
        Строка — атом вещества, столбец — атом растворителя. Яркая ячейка означает большой вес $\alpha_{ij}$: атом $i$ берёт больше информации от атома $j$.
      \end{block}
      \begin{block}{Что проверять}
        \scriptsize
        Для полярных систем должны выделяться OH...O, NH...O, C=O...H и похожие контакты. Для неполярных систем карта обычно менее локализована. Это не доказательство причинности, а диагностический способ понять, куда модель смотрит.
      \end{block}
    \end{column}
  \end{columns}
  \note{Здесь цель простая: научить аудиторию читать карту внимания. Не выдавать её за прямой физический эксперимент, но и не принижать её диагностическую ценность.}
\end{frame}

% 18
\begin{frame}{Шаг 4: NRTL-голова и SLE-решатель}
  \small
  Температура входит только в NRTL-голову:
  \[
    \tau_{12}(T)
    =
    \tau_{12}^{\mathrm{ref}}
    +
    \beta_{12}\left(\frac{1}{T}-\frac{1}{T_{\mathrm{ref}}}\right).
  \]
  Решатель:
  \[
    \Phi \to \gamma_2(x_2) \to
    x_2^{(k+1)}=\lambda\frac{\exp(-\Phi)}{\gamma_2^{(k)}}+(1-\lambda)x_2^{(k)}
    \to \ln x_2^*.
  \]
  Ограниченная коррекция:
  \[
    \ln x_2^{\text{итог}}
    =
    \ln x_2^{\text{физ}}
    +(1-w)\cdot\mathrm{clip}(\delta,\pm2.0).
  \]
  \note{Пояснить коррекцию как ограниченную поправку, а не обход физики.}
\end{frame}

% D1
\begin{frame}{Физическая модель неидеальна — нужен ``клапан''}
  \small
  Даже при хороших параметрах SLE+NRTL имеет встроенные приближения:
  \begin{itemize}\tightitems
    \item $\Delta C_p=0$ может давать ошибку порядка $0.85$ в $\ln x_2$;
    \item NRTL — компактная 3-параметрическая модель;
    \item полиморфизм, сольваты, ассоциация и специфическая химия вне решателя.
  \end{itemize}
  \vspace{0.25em}
  \begin{columns}[T,onlytextwidth]
    \begin{column}{0.48\textwidth}
      \begin{block}{Свободная поправка}
        \scriptsize
        MLP $\to \delta$ любого размера: модель быстро обходит физику и теряет интерпретируемость.
      \end{block}
    \end{column}
    \begin{column}{0.48\textwidth}
      \begin{block}{Ограниченная коррекция}
        \scriptsize
        Ограниченные поправки к физическим параметрам и повторное решение SLE: структура сохраняется.
      \end{block}
    \end{column}
  \end{columns}
  \note{Коррекция нужна не как лазейка вокруг физики, а как клапан для известных ограничений модели.}
\end{frame}

% D2
\begin{frame}{Коррекция через повторное решение SLE}
  \small
  \begin{block}{Идея}
    \scriptsize
    Не добавляем свободную поправку к ответу. Сначала ограниченно меняем физические параметры, затем \textbf{повторно решаем SLE}; так коррекция не обходит физику.
  \end{block}
  \vspace{-0.2em}
  \fitfig{
  \begin{tikzpicture}[>=Latex]
    \tikzstyle{b}=[draw,rounded corners=5pt,fill=white,minimum height=0.78cm,minimum width=1.75cm,align=center,font=\scriptsize,inner xsep=5pt]
    \node[b,draw=TGRed] (s1) at (0,0) {SLE$_1$\\физика};
    \node[b,draw=TGBlue] (d) at (2.10,0) {ограниченные\\поправки};
    \node[b,draw=TGBlue] (p) at (4.20,0) {скоррект.\\параметры};
    \node[b,draw=TGRed] (s2) at (6.30,0) {SLE$_2$\\предложение};
    \node[b,draw=TGOrange,minimum width=1.95cm] (g) at (8.55,0) {шлюз $w$\\SLE$_1$/SLE$_2$};
    \node[b,draw=TGGreen,minimum width=1.35cm] (o) at (10.75,0) {итог\\$\ln x_2$};
    \draw[->,thick,TGSlate!72] (s1)--(d);
    \draw[->,thick,TGSlate!72] (d)--(p);
    \draw[->,thick,TGSlate!72] (p)--(s2);
    \draw[->,thick,TGSlate!72] (s2)--(g);
    \draw[->,thick,TGSlate!72] (g)--(o);
    \node[font=\tiny,align=center,text=TGSlate,anchor=north] at (d.south)
      {$\delta T_m\in[-50,50]$ K\\$\delta\tau$ ограничена};
    \node[font=\tiny,align=center,text=TGSlate,anchor=north] at (s2.south)
      {предложение\\физически согласовано};
    \node[font=\tiny,align=center,text=TGSlate,anchor=north] at (g.south)
      {$w=1$: исходная физика\\$w=0$: предложение};
  \end{tikzpicture}}
  \vspace{-0.35em}
  \[
    \boxed{
    \ln x_2^{\text{итог}}
    =
    \ln x_2^{\text{физ}}
    +(1-w)\,\mathrm{clip}\!\left(
    \ln x_2^{\text{предл}}-\ln x_2^{\text{физ}},\pm2.0
    \right)}
  \]
  \vspace{-0.5em}
  \smallnote{$\pm2$ в $\ln x_2$ соответствует множителю $e^{\pm2}\approx[0.14,7.4]$ в $x_2$: это клапан, а не свободный обход решателя.}
  \note{Коррекция работает через физику: сначала меняем параметры, затем снова решаем SLE, затем шлюз решает, насколько доверять предложению.}
\end{frame}

\begin{frame}{Прямой проход: что сеть отдаёт физике}
  \vspace{-0.2em}
  \centering
  \resizebox{0.98\textwidth}{!}{%
  \begin{tikzpicture}[>=Latex]
    \tikzstyle{fbox}=[
      draw,
      rounded corners=8pt,
      fill=white,
      align=center,
      font=\scriptsize,
      inner sep=6pt,
      minimum height=0.76cm
    ]
    \node[fbox,draw=TGGreen!45,fill=TGGreen!7,minimum width=2.20cm] (inp) at (0,0)
      {граф вещества\\граф растворителя};
    \node[fbox,draw=TGBlue!45,fill=TGBlue!8,minimum width=2.05cm] (enc) at (2.75,0)
      {общий\\энкодер};
    \node[fbox,draw=TGBlue!45,fill=white,minimum width=2.85cm] (crys) at (5.85,1.05)
      {голова кристалла\\$T_m,\Delta H_{\mathrm{fus}},\Delta C_p$};
    \node[fbox,draw=TGOrange!50,fill=white,minimum width=2.85cm] (pair) at (5.85,-1.05)
      {NRTL-голова\\$\tau_{12}(T),\tau_{21}(T),\alpha$};
    \node[fbox,draw=TGRed!45,fill=TGRed!7,minimum width=3.10cm,minimum height=1.35cm] (solver) at (9.65,0)
      {жёсткая физика\\$\Phi(T)$ и $\ln\gamma_2(x_2)$\\$\Rightarrow\ln x_2^{\mathrm{физ}}$};
    \node[fbox,draw=TGBlue!45,fill=TGBlue!8,minimum width=2.75cm] (corr) at (13.15,0)
      {коррекция\\$\delta T_m,\delta\tau,w$};
    \node[fbox,draw=TGGreen!45,fill=TGGreen!8,minimum width=1.55cm] (out) at (16.05,0)
      {$\ln x_2^{\mathrm{итог}}$};

    \draw[->,line width=1pt,TGSlate!75] (inp) -- (enc);
    \draw[->,line width=1pt,TGSlate!75] (enc.east) -- (crys.west);
    \draw[->,line width=1pt,TGSlate!75] (enc.east) -- (pair.west);
    \draw[->,line width=1pt,TGSlate!75] (crys.east) -- ($(solver.west)+(0,0.32)$);
    \draw[->,line width=1pt,TGSlate!75] (pair.east) -- ($(solver.west)+(0,-0.32)$);
    \draw[->,line width=1pt,TGSlate!75] (solver) -- (corr);
    \draw[->,line width=1pt,TGSlate!75] (corr) -- (out);

    \node[fbox,draw=TGBlue!25,fill=TGBlue!5,minimum width=4.1cm,minimum height=0.85cm] at (6.15,-2.45)
      {$\Phi=\dfrac{\Delta H_{\mathrm{fus}}}{R}\!\left(\dfrac1T-\dfrac1{T_m}\right)$};
    \node[fbox,draw=TGOrange!30,fill=TGOrange!6,minimum width=4.2cm,minimum height=0.85cm] at (10.45,-2.45)
      {$\ln\gamma_2=\mathrm{NRTL}(x_2,\tau_{12},\tau_{21},\alpha)$};
    \node[fbox,draw=TGGreen!30,fill=TGGreen!6,minimum width=4.4cm,minimum height=0.85cm] at (14.90,-2.45)
      {$\ln x_2^{\mathrm{итог}}=\ln x_2^{\mathrm{физ}}+\text{шлюз + ограничение}$};
  \end{tikzpicture}}
  \vspace{0.35em}
  \smallnote{Синие/оранжевые блоки обучаемые; красный SLE/NRTL-решатель имеет 0 обучаемых параметров. Нейросеть выдаёт физические параметры, а растворимость “доделывает” решатель.}
  \note{На этом слайде важно проговорить трассу: энкодер выдаёт представления, головы переводят их в физические параметры, формулы считают растворимость, коррекция ограниченно правит результат.}
\end{frame}

\begin{frame}{Один пример через весь прямой проход}
  \small
  \begin{columns}[T,onlytextwidth]
    \begin{column}{0.34\textwidth}
      \centering
      \molasset[width=0.95\linewidth]{pair_paracetamol_ethanol.png}
      \smallnote{Парацетамол + этанол, $T=298.15$ K. Числа иллюстративные, но физически согласованные с SLE+NRTL.}
    \end{column}
    \begin{column}{0.62\textwidth}
      \scriptsize
      \begin{tabularx}{\linewidth}{lYY}
        \toprule
        \textbf{Шаг} & \textbf{Что выдаёт} & \textbf{Что делает физика}\\
        \midrule
        Кристаллическая голова & $T_m=442$ K, $\Delta H_{\mathrm{fus}}=26.4$ кДж/моль & $\Phi=3.46$\\
        NRTL-голова & $\tau_{12}=1.2$, $\tau_{21}=0.8$, $\alpha=0.30$ & $\ln\gamma_2^\infty=1.83$\\
        SLE-решатель & 0 обучаемых параметров & $\ln x_2\approx-3.46-1.83=-5.29$\\
        Интерпретация & $x_2\approx0.005$ & примерно 0.5 мол.\%\\
        \bottomrule
      \end{tabularx}
      \vspace{0.55em}
      \begin{block}{Что важно}
        \scriptsize
        Сеть не предсказывает растворимость напрямую. Она отдаёт физические параметры; формулы $\Phi$ и NRTL доделывают расчёт и дают разложение ошибки по причинам.
      \end{block}
    \end{column}
  \end{columns}
  \note{Это слайд для широкой аудитории: показать один числовой след от молекул до ответа. Он связывает ранее выведенную физику с архитектурной схемой.}
\end{frame}

\begin{frame}{Где подключаются функции потерь}
  \small
  \begin{columns}[c,onlytextwidth]
    \begin{column}{0.53\textwidth}
      \centering
      \vspace{0.15cm}
      \resizebox{\linewidth}{!}{%
      \begin{tikzpicture}[>=Latex]
        \tikzstyle{nodebox}=[
          draw,
          rounded corners=7pt,
          fill=white,
          align=center,
          font=\scriptsize,
          inner sep=5pt,
          minimum height=0.62cm
        ]
        \node[nodebox,draw=TGBlue!40,fill=TGBlue!7,minimum width=3.0cm] (model) at (0,0)
          {один прямой\\проход};
        \node[nodebox,draw=TGGreen!45,fill=TGGreen!8,minimum width=2.55cm] (sol) at (3.65,1.45)
          {$\ln x_2^{\mathrm{итог}}$};
        \node[nodebox,draw=TGBlue!45,fill=white,minimum width=2.55cm] (tm) at (3.65,0.45)
          {$T_m,\Delta H_{\mathrm{fus}}$};
        \node[nodebox,draw=TGOrange!45,fill=white,minimum width=2.55cm] (nrtl) at (3.65,-0.55)
          {$\tau_{12},\tau_{21},\alpha$};
        \node[nodebox,draw=TGPurple!45,fill=TGPurple!7,minimum width=2.55cm] (curve) at (3.65,-1.55)
          {кривая по $T$};

        \node[nodebox,draw=TGGreen!45,fill=TGGreen!6,minimum width=2.35cm] (lsol) at (7.15,1.45)
          {$\mathcal L_{\mathrm{sol}}$\\Huber};
        \node[nodebox,draw=TGBlue!45,fill=TGBlue!6,minimum width=2.35cm] (lcrys) at (7.15,0.45)
          {$\mathcal L_{T_m}$, $\mathcal L_{\Delta H}$};
        \node[nodebox,draw=TGOrange!45,fill=TGOrange!6,minimum width=2.35cm] (lgam) at (7.15,-0.55)
          {$\mathcal L_{\gamma^\infty}$\\IDAC};
        \node[nodebox,draw=TGPurple!45,fill=TGPurple!6,minimum width=2.35cm] (lshape) at (7.15,-1.55)
          {$\mathcal L_{\mathrm{mono}}$, $\mathcal L_{\mathrm{rank}}$, $\mathcal L_{\mathrm{vH}}$};

        \draw[->,line width=0.9pt,TGSlate!70] (model) -- (sol);
        \draw[->,line width=0.9pt,TGSlate!70] (model) -- (tm);
        \draw[->,line width=0.9pt,TGSlate!70] (model) -- (nrtl);
        \draw[->,line width=0.9pt,TGSlate!70] (model) -- (curve);
        \draw[->,line width=0.9pt,TGSlate!70] (sol) -- (lsol);
        \draw[->,line width=0.9pt,TGSlate!70] (tm) -- (lcrys);
        \draw[->,line width=0.9pt,TGSlate!70] (nrtl) -- (lgam);
        \draw[->,line width=0.9pt,TGSlate!70] (curve) -- (lshape);
      \end{tikzpicture}}
    \end{column}
    \begin{column}{0.43\textwidth}
      \begin{block}{Главная потеря}
        \scriptsize
        $\mathcal L_{\mathrm{sol}}$ сравнивает итоговое $\ln x_2$ с экспериментом. Это основной сигнал обучения.
      \end{block}
      \begin{block}{Вспомогательные якоря}
        \scriptsize
        $T_m$, $\Delta H_{\mathrm{fus}}$ и $\gamma^\infty$ учат промежуточные величины, чтобы NRTL не компенсировал ошибки кристаллической ветки.
      \end{block}
      \begin{block}{Форма температурной кривой}
        \scriptsize
        Монотонность, rank и Вант-Гофф — мягкие ограничения: они штрафуют нефизичную зависимость от температуры, но не заменяют экспериментальную потерю.
      \end{block}
    \end{column}
  \end{columns}
  \note{Этот слайд связывает архитектуру с обучением: каждая потеря цепляется не к абстрактной сети, а к конкретному физическому выходу.}
\end{frame}

\begin{frame}{Смесь экспертов по типам растворителей}
  \small
  \begin{columns}[T,onlytextwidth]
    \begin{column}{0.31\textwidth}
      \begin{block}{Зачем}
        \scriptsize
        Вода, спирты и неполярные растворители дают разные режимы ошибок. Одна общая поправка вынуждена усреднять несовместимые случаи.
      \end{block}
      \vspace{0.25em}
      \begin{block}{Примеры маршрутизации}
        \scriptsize
        \setlength{\tabcolsep}{2pt}
        \begin{tabular}{lcc}
          \toprule
          Растворитель & главный эксперт & вес\\
          \midrule
          вода & водный & 0.62\\
          этанол & спиртовой & 0.55\\
          гексан & неполярный & 0.81\\
          ДМСО & полярн. апрот. & 0.58\\
          \bottomrule
        \end{tabular}
      \end{block}
    \end{column}
    \begin{column}{0.67\textwidth}
      \centering
      \resizebox{\linewidth}{!}{%
      \begin{tikzpicture}[>=Latex]
        \tikzstyle{mcard}=[draw,rounded corners=8pt,fill=white,align=center,font=\scriptsize,inner sep=6pt]
        \node[mcard,draw=TGBlue!40,fill=TGBlue!7,minimum width=2.55cm,minimum height=0.72cm] (pair) at (0,0.90)
          {$\mathbf g_{\mathrm{pair}}$\\представление пары};
        \node[mcard,draw=TGGreen!40,fill=TGGreen!7,minimum width=2.55cm,minimum height=0.72cm] (type) at (0,-0.30)
          {тип растворителя\\категория + вектор};
        \node[mcard,draw=TGOrange!45,fill=TGOrange!8,minimum width=2.25cm,minimum height=1.05cm] (gate) at (3.35,0.30)
          {шлюз\\softmax-веса};

        \node[draw=TGLight,fill=white,rounded corners=11pt,minimum width=5.35cm,minimum height=3.25cm] (experts) at (7.20,0.30) {};
        \node[font=\scriptsize\bfseries,text=TGSlate] at (7.20,1.72) {экспертные остаточные поправки};
        \node[mcard,draw=TGBlue!35,fill=TGBlue!7,minimum width=1.55cm] (aq) at (5.55,0.93) {водный};
        \node[mcard,draw=TGOrange!35,fill=TGOrange!8,minimum width=1.55cm] (alc) at (7.20,0.93) {спиртовой};
        \node[mcard,draw=TGGreen!35,fill=TGGreen!8,minimum width=1.55cm] (pa) at (8.85,0.93) {полярный\\апротонный};
        \node[mcard,draw=TGPurple!35,fill=TGPurple!7,minimum width=1.55cm] (np) at (5.55,-0.45) {неполярный};
        \node[mcard,draw=TGRed!35,fill=TGRed!7,minimum width=1.55cm] (ar) at (7.20,-0.45) {ароматич.};
        \node[mcard,draw=TGSlate!30,fill=TGSlate!6,minimum width=1.55cm] (fb) at (8.85,-0.45) {общий};

        \node[mcard,draw=TGOrange!42,fill=TGOrange!7,minimum width=2.05cm,minimum height=0.86cm] (mix) at (11.30,0.30)
          {взвешенная\\сумма};
        \node[mcard,draw=TGGreen!40,fill=TGGreen!8,minimum width=2.25cm,minimum height=0.86cm] (out) at (14.20,0.30)
          {$\mathbf g'_{\mathrm{pair}}$\\для голов физики};

        \draw[->,line width=1.0pt,TGSlate!70] (pair.east) -- ($(gate.west)+(0,0.22)$);
        \draw[->,line width=1.0pt,TGSlate!70] (type.east) -- ($(gate.west)+(0,-0.22)$);
        \draw[->,line width=1.0pt,TGSlate!70] (gate.east) -- (experts.west);
        \draw[->,line width=1.0pt,TGSlate!70] (experts.east) -- (mix.west);
        \draw[->,line width=1.0pt,TGSlate!70] (mix.east) -- (out.west);
      \end{tikzpicture}}
      \smallnote{$\pi_1,\ldots,\pi_6$ — веса экспертов из шлюза. Итоговая поправка масштабируется малым обучаемым коэффициентом, поэтому смесь экспертов стартует почти как тождественное отображение.}
      \vspace{0.25em}
      \begin{block}{Математика коротко}
        \scriptsize
        \[
        \begin{aligned}
          \pi &=
          \mathrm{softmax}(G([\mathbf g\|\mathbf e_t])),\\
          \mathbf g' &=
          \mathbf g+\tanh(s)\sum_{e=1}^{6}\pi_e\,E_e(\mathbf g).
        \end{aligned}
        \]
        \smallnote{$\mathbf g=\mathbf g_{\mathrm{pair}}$; $s$ стартует около нуля. \texttt{moe\_balance} не даёт всем примерам уйти в одного эксперта.}
      \end{block}
    \end{column}
  \end{columns}
  \note{Важно сказать, что смесь экспертов есть в текущем репозитории и включена в основных конфигах, но это не центральная научная ставка, а инженерная адаптация под неоднородность растворителей.}
\end{frame}

\begin{frame}{Как кодируется температура}
  \small
  \begin{columns}[T,onlytextwidth]
    \begin{column}{0.48\textwidth}
      Одно сырое число $T$ в Кельвинах неудобно: масштаб большой, единица измерения произвольна, а физика SLE/NRTL часто зависит от $1/T$ и $\ln T$.
      \vspace{0.35em}
      \[
        \mathbf t(T)=
        \left[
        \frac{T-300}{50},\;
        \underbrace{300\left(\frac{1}{T}-\frac{1}{300}\right)}_{=\,300/T-1},\;
        \ln\frac{T}{300}
        \right].
      \]
      \smallnote{Это ровно компактное температурное состояние из \texttt{make\_temperature\_features(T)}: три безразмерных признака вместо одного сырого числа.}
      \vspace{0.35em}
      \begin{block}{Зачем не одним числом}
        \scriptsize
        Три безразмерных признака дают модели линейный, обратный и логарифмический масштаб температуры. Это устойчивее, чем подавать сырые Кельвины в MLP.
      \end{block}
    \end{column}
    \begin{column}{0.48\textwidth}
      \fitfig{
      \begin{tikzpicture}
        \begin{axis}[
          width=6.6cm,height=4.4cm,
          xmin=240,xmax=430,
          xlabel={$T$, K},
          ylabel={значение признака},
          grid=both,
          legend style={font=\tiny,draw=none,fill=white,fill opacity=0.72,text opacity=1},
          legend pos=north east
        ]
          \addplot[domain=240:430,samples=120,TGBlue,thick] {(x-300)/50};
          \addlegendentry{$(T-300)/50$}
          \addplot[domain=240:430,samples=120,TGOrange,thick] {(1/x - 1/300)*300};
          \addlegendentry{$300/T-1$}
          \addplot[domain=240:430,samples=120,TGGreen,thick] {ln(x/300)};
          \addlegendentry{$\ln(T/300)$}
        \end{axis}
      \end{tikzpicture}}
      \smallnote{$T$ не подаётся в кристаллические головы: иначе модель могла бы сделать $T_m$ зависимым от рабочей температуры эксперимента.}
    \end{column}
  \end{columns}
  \note{Объяснить: температурный код нужен не для красоты, а чтобы дать модели физически удобные координаты T, 1/T и log T без утечки T в свойства кристалла.}
\end{frame}

% 20
\begin{frame}{DirectGNN: контрольная базовая модель}
  \begin{columns}[T,onlytextwidth]
    \begin{column}{0.44\textwidth}
      \begin{itemize}\tightitems
        \item тот же графовый энкодер;
        \item то же перекрёстное внимание;
        \item тот же агрегатор чтения графа;
        \item вместо физического пути: кодирование $T$ + прямой MLP.
      \end{itemize}
    \end{column}
    \begin{column}{0.52\textwidth}
      \fitfig{
      \begin{tikzpicture}[>=Latex]
        \tikzstyle{box}=[draw,rounded corners=3pt,minimum height=0.65cm,align=center,font=\scriptsize]
        \node[box,fill=TGBlue!10,minimum width=1.8cm] (g) at (0,1.0) {GNN\\пара};
        \node[box,fill=TGOrange!10,minimum width=1.9cm] (t) at (0,-0.6) {термо\\код $T$};
        \node[box,fill=TGBlue!10,minimum width=1.8cm] (cat) at (2.7,0.25) {объединение};
        \node[box,fill=TGOrange!12,minimum width=1.7cm] (m) at (5.1,0.25) {прямой\\MLP};
        \node[box,fill=TGGreen!12,minimum width=1.4cm] (y) at (7.2,0.25) {$\ln x_2$};
        \draw[->,thick] (g)--(cat);
        \draw[->,thick] (t)--(cat);
        \draw[->,thick] (cat)--(m)--(y);
        \draw[TGOrange,thick,rounded corners=2pt] (-1.15,-1.45) rectangle (-0.85,0.30);
        \foreach \yy/\op in {-1.25/15,-0.95/22,-0.65/35,-0.35/55,-0.05/80}
          \fill[TGOrange!\op] (-1.13,\yy) rectangle (-0.87,\yy+0.22);
        \node[font=\tiny,text=TGSlate,align=center] at (-1.0,-1.75) {термометрическое\\кодирование};
      \end{tikzpicture}}
    \end{column}
  \end{columns}
  \keybox{Разница TGNN и DirectGNN измеряет чистый эффект физического узкого места.}
  \note{Переход: нужна ещё одна базовая модель — фиксированные дескрипторы без GNN.}
\end{frame}

% 21
\begin{frame}{RF-базовая модель: необучаемый потолок}
  \[
    \hat{y}=f_{\mathrm{RF}}([\mathbf{d}_{\text{вещ}}\|\mathbf{d}_{\text{раств}}\|T]).
  \]
  \begin{itemize}\tightitems
    \item RDKit-дескрипторы: MolLogP, TPSA, MolWt, HBD/HBA, отпечаток Morgan.
    \item Не требует видеокарты и обучается за минуты.
    \item Не имеет физической структуры и не предсказывает $T_m,\Delta H,\tau$.
  \end{itemize}
  \keybox{Роль RF: оценить, сколько даёт фиксированная экспертная химия без обучения нейросетевых представлений.}
  \note{Объяснить, почему RF может быть сильным и почему это не провал идеи TGNN.}
\end{frame}

% 22
\begin{frame}{Три уровня сравнения}
  \small
  \begin{tabularx}{\textwidth}{lYY}
    \toprule
    \textbf{Модель} & \textbf{Что сравниваем} & \textbf{Роль}\\
    \midrule
    RF-дескрипторы & фиксированные признаки + деревья & потолок дескрипторов\\
    DirectGNN & обучаемые графовые представления & чистое машинное обучение\\
    TGNN-Solv & GNN + физический решатель & физика + машинное обучение\\
    \bottomrule
  \end{tabularx}
  \vspace{0.6em}
  \begin{center}
    \Large\textcolor{TGOrange}{RF $>$ DirectGNN $>$ TGNN? Или наоборот?}
  \end{center}
  \note{Переход к обучению: прежде чем сравнивать, нужно заставить физический путь стабильно учиться.}
\end{frame}

\begin{frame}{FastSolv и SolProp: внешние базовые модели}
  \small
  \begin{columns}[T,onlytextwidth]
    \begin{column}{0.48\textwidth}
      \begin{block}{FastSolv}
        \scriptsize
        Внешняя дескрипторная базовая модель. В репозитории есть обёртка для предсказания, сравнения и обучения на наших разбиениях по химическим каркасам.
        \vspace{0.25em}

        Сильная сторона: быстрый внешний ориентир без нашей физической архитектуры.\\
        Ограничение для нашего вопроса: не даёт $T_m$, $\Delta H$, NRTL-параметры и SLE-декомпозицию.
      \end{block}
    \end{column}
    \begin{column}{0.48\textwidth}
      \begin{block}{SolProp}
        \scriptsize
        Внешняя архитектура для растворимости, обычно в отдельном окружении. Обёртка поддерживает режим без дообучения, калибровку на обучающем разбиении и полное переобучение на нашей цели $\ln x_2$.
        \vspace{0.25em}

        Сильная сторона: внешний опубликованный ориентир.\\
        Ограничение: режим без дообучения не является строго контролируемой абляцией; для статьи честнее полное переобучение на нашей цели.
      \end{block}
    \end{column}
  \end{columns}
  \vspace{0.35em}
  \keybox{FastSolv/SolProp отвечают на вопрос “как мы выглядим против внешних моделей”, но не заменяют главное сравнение TGNN-Solv vs DirectGNN.}
  \note{Подчеркнуть: внешние модели полезны как конкуренты, но не дают чистого ответа на physics bottleneck, потому что меняется не только физика, но и вся архитектура/обучающий стек.}
\end{frame}

\begin{frame}{Как сравниваем внешние базовые модели}
  \small
  \begin{columns}[T,onlytextwidth]
    \begin{column}{0.50\textwidth}
      \begin{block}{Протокол}
        \scriptsize
        \begin{itemize}\tightitems
          \item те же обучающее, валидационное и тестовое разбиения;
          \item та же целевая величина $\ln x_2$;
          \item те же метрики: MAE, RMSE, $R^2$;
          \item тот же набор артефактов: предсказания, отчёт, манифест, карточка сравнения;
          \item отдельно отмечаем: без дообучения, с калибровкой или с полным переобучением.
        \end{itemize}
      \end{block}
    \end{column}
    \begin{column}{0.46\textwidth}
      \begin{block}{Что контролируем}
        \scriptsize
        \begin{itemize}\tightitems
          \item пересчёт единиц: `logS` $\leftrightarrow$ $\ln x_2$;
          \item температуру: 298.15 K или явный температурно-зависимый режим;
          \item утечки: сравнение только на тестовом разбиении;
          \item окружение: внешние зависимости не должны менять наш путь обучения и оценки.
        \end{itemize}
      \end{block}
      \begin{block}{Чего они не объясняют}
        \scriptsize
        Почему ответ такой: кристалл, активность, NRTL, коррекция. Это остаётся преимуществом TGNN-Solv как интерпретируемой модели.
      \end{block}
    \end{column}
  \end{columns}
  \note{Этот слайд нужен, чтобы аудитория не смешивала внешние базовые модели и строго контролируемую абляцию DirectGNN.}
\end{frame}

\begin{frame}{Где мы сейчас?}
  \centering
  \vspace{0.2cm}
  \roadmapbar{3}
  \vspace{0.75cm}

  {\fontsize{19}{23}\selectfont
  Архитектура готова: графовый энкодер, перекрёстное внимание, NRTL-голова и SLE-решатель.\par}
  \vspace{0.55cm}
  {\fontsize{21}{25}\selectfont\bfseries\color{TGBlueDark}
  Но обучить её наивно нельзя:\\
  решатель получает нефизичные параметры,\\
  и итерации становятся нестабильными.\par}
  \note{Переход к поэтапному обучению: сложная физическая архитектура требует поэтапного обучения, иначе случайные головы делают решатель нестабильным.}
\end{frame}

\section{Обучение}

\begin{frame}{Этап 0: зачем предобучать энкодер}
  \small
  \begin{center}
  \resizebox{0.96\textwidth}{!}{%
  \begin{tikzpicture}[>=Latex]
    \tikzstyle{pcard}=[draw,rounded corners=9pt,fill=white,align=center,font=\footnotesize,inner sep=7pt,minimum height=0.88cm]
    \node[pcard,draw=TGGreen!40,fill=TGGreen!7,minimum width=2.85cm] (smi) at (0,0)
      {много SMILES\\без $\ln x_2$};
    \node[pcard,draw=TGBlue!40,fill=TGBlue!7,minimum width=2.95cm] (graph) at (3.7,0)
      {графы RDKit\\атомы + связи};
    \node[pcard,draw=TGBlue!50,fill=TGBlue!9,minimum width=3.20cm] (tasks) at (7.65,0)
      {учебные задачи\\маски и дескрипторы};
    \node[pcard,draw=TGOrange!45,fill=TGOrange!7,minimum width=2.95cm] (ckpt) at (11.75,0)
      {тёплый старт\\чекпойнт};
    \node[pcard,draw=TGGreen!45,fill=TGGreen!8,minimum width=3.25cm] (down) at (15.85,0)
      {основное\\дообучение};
    \draw[->,line width=1.1pt,TGSlate!70] (smi) -- (graph);
    \draw[->,line width=1.1pt,TGSlate!70] (graph) -- (tasks);
    \draw[->,line width=1.1pt,TGSlate!70] (tasks) -- (ckpt);
    \draw[->,line width=1.1pt,TGSlate!70] (ckpt) -- (down);
  \end{tikzpicture}}
  \end{center}
  \vspace{0.25em}
  \begin{columns}[c,onlytextwidth]
    \begin{column}{0.34\textwidth}
      \centering
      \molasset[width=0.88\linewidth]{paracetamol_plain.png}\\[-0.15em]
      {\scriptsize\color{TGSlate}структура из SMILES}
    \end{column}
    \begin{column}{0.34\textwidth}
      \centering
      \molasset[width=0.90\linewidth]{paracetamol_graph.png}\\[-0.15em]
      {\scriptsize\color{TGSlate}тот же объект как граф}
    \end{column}
    \begin{column}{0.28\textwidth}
      \begin{itemize}\tightitems
        \item источник: ZINC250k или SMILES из BigSolDB;
        \item учатся GNN и агрегатор;
        \item временные головы затем удаляются;
        \item чекпойнт используется как старт.
      \end{itemize}
    \end{column}
  \end{columns}
  \smallnote{Это не фаза 1: этап 0 учит модель читать молекулы до обучения на растворимости.}
  \note{Для широкой аудитории: это как предварительно научить модель химическому алфавиту до задачи растворимости.}
\end{frame}

\begin{frame}{Легенда: как читать графовые слайды}
  \small
  \begin{columns}[T,onlytextwidth]
    \begin{column}{0.40\textwidth}
      \centering
      \molasset[width=0.86\linewidth]{paracetamol_graph.png}\\[-0.15em]
      {\scriptsize\color{TGSlate}парацетамол как граф из RDKit}
      \vspace{0.35em}

      \resizebox{0.90\linewidth}{!}{%
      \begin{tikzpicture}
        \tikzstyle{legatom}=[circle,minimum size=0.48cm,font=\scriptsize\bfseries]
        \node[legatom,draw=TGBlue!55,fill=TGBlue!10] at (0,0) {C};
        \node[anchor=west,font=\scriptsize] at (0.38,0) {углерод};
        \node[legatom,draw=TGGreen!65,fill=TGGreen!12] at (2.10,0) {N};
        \node[anchor=west,font=\scriptsize] at (2.48,0) {азот};
        \node[legatom,draw=TGOrange!65,fill=TGOrange!12] at (4.0,0) {O};
        \node[anchor=west,font=\scriptsize] at (4.38,0) {кислород};
        \node[legatom,draw=TGSlate,fill=TGSlate!84,text=white] at (6.25,0) {?};
        \node[anchor=west,font=\scriptsize] at (6.63,0) {скрытый атом};
      \end{tikzpicture}}
    \end{column}
    \begin{column}{0.56\textwidth}
      \begin{block}{Что такое граф тяжёлых атомов}
        \scriptsize
        Узлы \textemdash{} атомы без явных водородов. Водороды обычно хранятся как признаки атома: число H, донорность, валентность. Это стандартный формат RDKit/PyG и он не теряет ключевую информацию для органических молекул.
      \end{block}
      \begin{block}{Что получает сеть}
        \scriptsize
        Для каждого атома: тип элемента, заряд, гибридизация, ароматичность, электроотрицательность, радиус, поляризуемость и другие признаки. Для каждой связи: тип, кольцо, сопряжение, стерео и дополнительные физические признаки для варианта с физическими каналами.
      \end{block}
      \begin{block}{Что означает маска}
        \scriptsize
        Маска \textemdash{} искусственно удалённая информация. Если модель восстанавливает её из соседей, значит она выучила химический контекст, а не просто запомнила список атомов.
      \end{block}
    \end{column}
  \end{columns}
  \note{Этот слайд нужен перед задачами предобучения: аудитория должна понимать, что на следующих рисунках скрывается и что модель пытается восстановить.}
\end{frame}

\begin{frame}{Этап 0: четыре учебные задачи}
  \small
  \begin{columns}[T,onlytextwidth]
    \begin{column}{0.49\textwidth}
      \begin{block}{1. Маскированный двухшаговый подграф}
        \scriptsize
        Зануляем связную окрестность атома и просим восстановить признаки атомов. Модель учится локальной химической топологии.
      \end{block}
      \begin{block}{2. Предсказание типа связи}
        \scriptsize
        По состояниям двух атомов восстановить тип связи: одинарная, двойная, тройная или ароматическая.
      \end{block}
    \end{column}
    \begin{column}{0.49\textwidth}
      \begin{block}{3. RDKit-свойства}
        \scriptsize
        Предсказать 12 вычисляемых дескрипторов: MolLogP, MolWt, TPSA, доноры/акцепторы H-связей, FractionCSP3 и другие.
      \end{block}
      \begin{block}{4. Контрастивное обучение}
        \scriptsize
        Две слегка испорченные версии одной молекулы должны иметь близкие векторы; разные молекулы — далёкие.
      \end{block}
    \end{column}
  \end{columns}
  \vspace{0.25em}
  \[
    \mathcal L_{0}
    =
    1.0\,\mathcal L_{\mathrm{atom}}
    +0.5\,\mathcal L_{\mathrm{bond}}
    +1.0\,\mathcal L_{\mathrm{prop}}
    +0.5\,\mathcal L_{\mathrm{contrastive}}.
  \]
  \smallnote{Параметры по умолчанию: 30 эпох, пакет 128, AdamW, косинусное расписание скорости обучения, обрезка градиента 1.0.}
  \note{Проговорить, что каждая задача дешёвая: не нужны лабораторные измерения растворимости, достаточно SMILES и RDKit.}
\end{frame}

\begin{frame}{Этап 0a: локальная химия в графе}
  \small
  \vspace{-0.25em}
  \centering
  \resizebox{0.97\textwidth}{!}{%
  \begin{tikzpicture}[>=Latex]
    \tikzstyle{catom}=[circle,draw=TGBlue!55,fill=TGBlue!10,minimum size=0.62cm,font=\scriptsize\bfseries]
    \tikzstyle{oatom}=[circle,draw=TGOrange!65,fill=TGOrange!12,minimum size=0.62cm,font=\scriptsize\bfseries]
    \tikzstyle{natom}=[circle,draw=TGGreen!65,fill=TGGreen!12,minimum size=0.62cm,font=\scriptsize\bfseries]
    \tikzstyle{maskatom}=[circle,draw=TGSlate,fill=TGSlate!84,minimum size=0.62cm,font=\scriptsize\bfseries,text=white]
    \tikzstyle{card}=[draw,rounded corners=8pt,fill=white,align=center,font=\scriptsize,inner sep=6pt]

    \node[catom] (c1) at (0.0,0.0) {C};
    \node[catom] (c2) at (0.85,1.10) {C};
    \node[catom] (c3) at (2.20,1.05) {C};
    \node[maskatom] (c4) at (3.05,0.00) {?};
    \node[catom] (c5) at (2.20,-1.05) {C};
    \node[catom] (c6) at (0.85,-1.10) {C};
    \foreach \a/\b in {c1/c2,c2/c3,c3/c4,c4/c5,c5/c6,c6/c1}
      \draw[line width=1.45pt,TGSlate!62] (\a) -- (\b);
    \draw[line width=0.85pt,TGSlate!36] (1.52,0) circle (0.72cm);

    \node[oatom] (oh) at (-1.08,-0.18) {O};
    \draw[line width=1.45pt,TGSlate!62] (c1) -- (oh);
    \node[font=\tiny,text=TGSlate] at (-1.35,0.42) {OH};

    \node[maskatom] (n) at (4.15,0.20) {?};
    \node[maskatom] (co) at (5.25,0.18) {?};
    \node[oatom] (oxo) at (6.02,1.02) {O};
    \node[maskatom] (me) at (6.30,-0.62) {?};
    \draw[line width=1.45pt,TGSlate!62] (c4) -- (n);
    \draw[line width=1.45pt,TGSlate!62] (n) -- (co);
    \draw[line width=3.0pt,TGOrange!86] (co) -- (oxo);
    \draw[line width=1.45pt,TGSlate!62] (co) -- (me);

    \node[card,draw=TGSlate!35,fill=white] at (3.55,-1.75)
      {парацетамол: граф тяжёлых атомов\\маска вокруг амидной группы};
    \draw[->,line width=1.0pt,TGSlate!60] (3.55,-1.35) -- (n);
    \draw[->,line width=1.0pt,TGSlate!60] (3.55,-1.35) -- (co);
    \draw[->,line width=1.0pt,TGSlate!60] (3.55,-1.35) -- (me);

    \node[card,draw=TGBlue!35,fill=TGBlue!7,minimum width=2.15cm,minimum height=0.95cm] (gnn) at (8.35,0.20)
      {GNN\\контекст};
    \node[card,draw=TGGreen!40,fill=TGGreen!8,minimum width=2.30cm,minimum height=0.78cm] (atomhead) at (11.30,0.85)
      {восстановить\\признаки атомов};
    \node[card,draw=TGOrange!45,fill=TGOrange!8,minimum width=2.30cm,minimum height=0.78cm] (bondhead) at (11.30,-0.55)
      {предсказать\\тип связи};

    \draw[->,line width=1.15pt,TGSlate!70] (6.70,0.18) -- (gnn.west);
    \draw[->,line width=1.15pt,TGSlate!70] (gnn.east) -- (atomhead.west);
    \draw[->,line width=1.15pt,TGSlate!70] (gnn.east) -- (bondhead.west);
    \node[font=\scriptsize,text=TGSlate,align=center] at (11.30,1.70)
      {$\hat{\mathbf x}_{mask}\approx\mathbf x_{mask}$};
    \node[font=\scriptsize,text=TGSlate,align=center] at (11.30,-1.32)
      {$\hat y_{\mathrm{bond}}\in\{\text{один.},\text{двойн.},\text{тройн.},\text{аром.}\}$};
  \end{tikzpicture}}
  \vspace{0.25em}
  \begin{columns}[T,onlytextwidth]
    \begin{column}{0.49\textwidth}
      \begin{block}{Маскированный подграф}
        \scriptsize
        Зануляем признаки связной окрестности и просим сеть восстановить атомные векторы из химического контекста:
        \[
          \mathcal L_{\mathrm{atom}}=\|\hat{\mathbf x}_{mask}-\mathbf x_{mask}\|_2^2.
        \]
      \end{block}
    \end{column}
    \begin{column}{0.49\textwidth}
      \begin{block}{Тип связи}
        \scriptsize
        Для части рёбер зануляем признаки связи и классифицируем её по состояниям двух конечных атомов:
        \[
          \mathcal L_{\mathrm{bond}}=\mathrm{CE}(\hat y_{\mathrm{bond}},y_{\mathrm{bond}}).
        \]
      \end{block}
    \end{column}
  \end{columns}
  \smallnote{Чёрные узлы \textemdash{} скрытые атомы; оранжевая связь \textemdash{} целевая связь. Так сеть учит химический синтаксис до обучения на растворимости.}
  \note{Смысл: сеть учится химическому синтаксису — какие атомы и связи правдоподобны в данном окружении.}
\end{frame}

\begin{frame}{Этап 0b/c: глобальный смысл молекулы}
  \small
  \begin{columns}[T,onlytextwidth]
    \begin{column}{0.49\textwidth}
      \begin{block}{RDKit-свойства}
        \scriptsize
        После агрегации графа получаем $\mathbf g_{\mathrm{mol}}$ и просим предсказать 12 бесплатных дескрипторов.
      \end{block}
      \centering
      \resizebox{\linewidth}{!}{%
      \begin{tikzpicture}[>=Latex]
        \tikzstyle{card}=[draw,rounded corners=7pt,fill=white,align=center,font=\scriptsize,inner sep=5pt,minimum height=0.68cm]
        \node[card,draw=TGBlue!35,fill=TGBlue!7,minimum width=2.0cm] (g) at (0,0) {$\mathbf g_{\mathrm{mol}}$};
        \node[card,draw=TGOrange!35,fill=TGOrange!7,minimum width=2.35cm] (h) at (2.9,0) {голова\\свойств};
        \node[card,draw=TGGreen!35,fill=TGGreen!7,minimum width=1.35cm] (p1) at (5.7,0.85) {MolLogP};
        \node[card,draw=TGGreen!35,fill=TGGreen!7,minimum width=1.35cm] (p2) at (7.25,0.85) {TPSA};
        \node[card,draw=TGGreen!35,fill=TGGreen!7,minimum width=1.35cm] (p3) at (5.7,0.0) {MolWt};
        \node[card,draw=TGGreen!35,fill=TGGreen!7,minimum width=1.35cm] (p4) at (7.25,0.0) {HBD/HBA};
        \node[card,draw=TGGreen!35,fill=TGGreen!7,minimum width=2.9cm] (p5) at (6.48,-0.85) {ещё 8 дескрипторов};
        \draw[->,line width=0.9pt,TGSlate!70] (g) -- (h);
        \foreach \p in {p1,p2,p3,p4,p5}
          \draw[->,line width=0.7pt,TGSlate!55] (h.east) -- (\p.west);
      \end{tikzpicture}}
      \[
        \mathcal L_{\mathrm{prop}}=\|\hat{\mathbf p}-\mathbf p_{\mathrm{RDKit}}\|_2^2.
      \]
    \end{column}
    \begin{column}{0.49\textwidth}
      \begin{block}{Контрастивное обучение}
        \scriptsize
        Делаем две слегка испорченные версии одной молекулы. Их векторы должны быть близко; векторы разных молекул — далеко.
      \end{block}
      \centering
      \resizebox{\linewidth}{!}{%
      \begin{tikzpicture}[>=Latex]
        \tikzstyle{view}=[draw,rounded corners=7pt,fill=white,align=center,font=\scriptsize,inner sep=5pt,minimum width=2.2cm,minimum height=0.85cm]
        \node[view,draw=TGBlue!35,fill=TGBlue!7] (v1) at (0,0.75) {вид A\\маска узлов};
        \node[view,draw=TGBlue!35,fill=TGBlue!7] (v2) at (0,-0.75) {вид B\\маска рёбер};
        \node[view,draw=TGOrange!35,fill=TGOrange!7] (enc1) at (3.0,0.75) {GNN\\$z_A$};
        \node[view,draw=TGOrange!35,fill=TGOrange!7] (enc2) at (3.0,-0.75) {GNN\\$z_B$};
        \node[draw=TGGreen!40,fill=TGGreen!8,rounded corners=7pt,align=center,font=\scriptsize,minimum width=2.4cm,minimum height=0.8cm] (pos) at (6.0,0)
          {та же молекула\\притянуть};
        \node[draw=TGRed!40,fill=TGRed!7,rounded corners=7pt,align=center,font=\scriptsize,minimum width=2.4cm,minimum height=0.8cm] (neg) at (6.0,-1.55)
          {другая молекула\\оттолкнуть};
        \draw[->,line width=0.9pt,TGSlate!70] (v1) -- (enc1);
        \draw[->,line width=0.9pt,TGSlate!70] (v2) -- (enc2);
        \draw[->,line width=0.9pt,TGGreen!75] (enc1.east) -- (pos.west);
        \draw[->,line width=0.9pt,TGGreen!75] (enc2.east) -- (pos.west);
        \draw[->,line width=0.8pt,TGRed!70,dashed] (enc2.east) -- (neg.west);
      \end{tikzpicture}}
      \[
        \mathcal L_{\mathrm{contrastive}}
        =\mathrm{CE}(z_Az_B^\top/\tau,\; \text{та же молекула}).
      \]
    \end{column}
  \end{columns}
  \note{Смысл: потеря по свойствам учит глобальные химические свойства, контрастивная потеря учит устойчивость к небольшим искажениям графа.}
\end{frame}

\begin{frame}{Что меняется после предобучения}
  \small
  \centering
  \resizebox{0.92\textwidth}{!}{%
  \begin{tikzpicture}[>=Latex]
    \tikzstyle{panel}=[draw=TGLight,fill=white,rounded corners=10pt,minimum width=4.65cm,minimum height=3.35cm]
    \node[panel] (before) at (0,0) {};
    \node[panel] (after) at (8.3,0) {};
    \node[font=\scriptsize\bfseries,text=TGSlate] at (0,1.80) {до этапа 0: случайная геометрия};
    \node[font=\scriptsize\bfseries,text=TGSlate] at (8.3,1.80) {после этапа 0: химически осмысленные группы};

    \foreach \x/\y/\c in {
      -1.65/0.85/TGBlue,-1.05/-0.80/TGOrange,-0.42/0.25/TGGreen,
      0.18/-1.00/TGBlue,0.72/0.75/TGOrange,1.25/-0.20/TGGreen,
      1.55/1.10/TGRed,-1.35/-0.20/TGRed,-0.85/1.15/TGPurple,
      0.52/0.05/TGPurple,1.65/-0.85/TGSlate,-1.75/-1.20/TGSlate}
      \fill[\c!72] (\x,\y) circle (0.08);

    \foreach \x/\y in {6.65/0.95,6.95/0.72,7.18/1.05,6.78/1.25}
      \fill[TGBlue!72] (\x,\y) circle (0.08);
    \foreach \x/\y in {8.00/-0.90,8.28/-1.10,8.55/-0.78,8.12/-0.58}
      \fill[TGOrange!78] (\x,\y) circle (0.08);
    \foreach \x/\y in {9.40/0.88,9.68/0.65,9.82/1.12,9.25/1.22}
      \fill[TGGreen!78] (\x,\y) circle (0.08);
    \foreach \x/\y in {9.18/-0.48,9.52/-0.28,9.75/-0.62,9.35/-0.88}
      \fill[TGRed!72] (\x,\y) circle (0.08);
    \draw[TGBlue!35,rounded corners=14pt,line width=0.9pt] (6.42,0.55) rectangle (7.42,1.45);
    \draw[TGOrange!35,rounded corners=14pt,line width=0.9pt] (7.78,-1.35) rectangle (8.78,-0.42);
    \draw[TGGreen!35,rounded corners=14pt,line width=0.9pt] (9.04,0.45) rectangle (10.04,1.45);
    \draw[TGRed!35,rounded corners=14pt,line width=0.9pt] (8.98,-1.08) rectangle (9.92,-0.10);

    \node[font=\tiny,text=TGBlueDark] at (6.92,0.34) {ароматика};
    \node[font=\tiny,text=TGOrange!80!black] at (8.28,-1.58) {спирты/OH};
    \node[font=\tiny,text=TGGreen!70!black] at (9.55,0.32) {амиды};
    \node[font=\tiny,text=TGRed!80!black] at (9.45,-1.30) {карбонилы};

    \node[draw=TGOrange!35,fill=TGOrange!8,rounded corners=8pt,align=center,font=\scriptsize,minimum width=2.6cm,minimum height=0.78cm] (stage0) at (4.15,0)
      {4 задачи\\без ручных меток};
    \draw[->,line width=1.1pt,TGSlate!70] (2.45,0) -- (stage0.west);
    \draw[->,line width=1.1pt,TGSlate!70] (stage0.east) -- (5.86,0);
  \end{tikzpicture}}
  \vspace{0.35em}
  \begin{block}{Интерпретация}
    \scriptsize
    Этап 0 не видит $\ln x_2$ и не учит растворимость напрямую. Он делает стартовую точку энкодера менее случайной: похожие химические мотивы оказываются ближе в пространстве представлений, а дальнейшему обучению остаётся связать эти мотивы с физическими параметрами.
  \end{block}
  \note{Это визуальный ответ на вопрос, зачем предобучение: не ради прямой метрики растворимости, а ради геометрии представлений.}
\end{frame}

% 23
\begin{frame}{Трёхфазное обучение}
  \fitfig{
  \begin{tikzpicture}
    \draw[thick,TGSlate] (0,0) -- (12,0);
    \foreach \x/\lab in {0/0,2/50,10/250,12/300}
      \draw (\x,0.08)--(\x,-0.08) node[below,font=\scriptsize] {\lab};
    \fill[TGBlue!18] (0,0.35) rectangle (2,1.25);
    \fill[TGOrange!20] (2,0.35) rectangle (10,1.25);
    \fill[TGGreen!18] (10,0.35) rectangle (12,1.25);
    \node[align=center,font=\scriptsize] at (1,0.80) {Фаза 1\\только головы\\$L_{sol}=0$};
    \node[align=center,font=\scriptsize] at (6,0.80) {Фаза 2\\полный SLE\\12 потерь};
    \node[align=center,font=\scriptsize] at (11,0.80) {Фаза 3\\fine tune\\LR $\downarrow$};
  \end{tikzpicture}}
  \begin{itemize}\tightitems
    \item Фаза 1: разумные $T_m$, $\Delta H$, Hansen, $\gamma^\infty$.
    \item Фаза 2: включаем SLE, потеря по растворимости доминирует.
    \item Фаза 3: сильнее монотонность, слабее коррекция.
  \end{itemize}
  \note{Сказать: если включить всё сразу, случайные головы ломают решатель и градиенты.}
\end{frame}

% 24
\begin{frame}{12 компонент функции потерь}
  \small
  \begin{columns}[T,onlytextwidth]
    \begin{column}{0.50\textwidth}
      \[
        \mathcal{L}
        =
        \lambda_{\mathrm{sol}}\mathcal{L}_{\mathrm{sol}}
        +\lambda_{T_m}\mathcal{L}_{T_m}
        +\lambda_{\Delta H}\mathcal{L}_{\Delta H}
        +\lambda_{\gamma^\infty}\mathcal{L}_{\gamma^\infty}
        +\ldots
      \]
      \[
        \mathcal{L}_{\mathrm{sol}}
        =
        \mathrm{Huber}_{\delta=1}(\Delta\ln x_2).
      \]
      \[
        f_{\mathrm{sol}}
        =
        \frac{\lambda_{\mathrm{sol}}\mathcal{L}_{\mathrm{sol}}}{\mathcal{L}_{\mathrm{total}}}
        >50\%.
      \]
      \vspace{0.2em}
      {\scriptsize
      \centering
      \begin{tabular}{lccc}
        \toprule
        & \textbf{Фаза 1} & \textbf{Фаза 2} & \textbf{Фаза 3}\\
        \midrule
        $\mathcal{L}_{sol}$ & 0 & \textbf{1.0} & \textbf{1.0}\\
        $\mathcal{L}_{T_m}$ & \textbf{1.0} & 0.1 & 0.05\\
        $\mathcal{L}_{\Delta H}$ & \textbf{1.0} & 0.1 & 0.05\\
        $\mathcal{L}_{mono}$ & 0 & 0.1 & \textbf{0.5}\\
        $\mathcal{L}_{\gamma^\infty}$ & \textbf{1.0} & 0.1 & 0.05\\
        \bottomrule
      \end{tabular}
      \par}
      \smallnote{Полная таблица 12 весов вынесена в дополнительные материалы.}
    \end{column}
    \begin{column}{0.46\textwidth}
	      \begin{block}{Почему понадобилась стабилизация vH}
        \raggedright
        \scriptsize
	        $L_{\mathrm{vH}}$ — согласованность Вант-Гоффа: штраф за неправильный локальный наклон $\ln x_2$ по $1/T$ внутри одной пары. На шумных или почти плоских температурных сериях он доминировал над основной задачей:
        \[
          L_{\mathrm{vH}}:\;3811\to240148,\qquad f_{\mathrm{sol}}<1\%.
        \]
        Конкретные исправления:
        \begin{itemize}\tightitems
	          \item считать только пары с $\ge2$ конечными температурными точками;
	          \item маска некорректных численных значений и \texttt{nan\_to\_num} перед суммой;
	          \item ограничить остаток наклона и вклад $L_{\mathrm{vH}}$ в суммарную потерю;
	          \item меньший фазовый вес, чтобы $f_{\mathrm{sol}}>50\%$ оставался инвариантом.
        \end{itemize}
      \end{block}
    \end{column}
  \end{columns}
  \note{Подчеркнуть sol_fraction как инженерный инвариант обучения, а не просто метрику в логах.}
\end{frame}

% E1
\begin{frame}{Робастная регрессия: не реагировать на выбросы}
  \small
  \begin{columns}[T,onlytextwidth]
    \begin{column}{0.48\textwidth}
      MSE штрафует выбросы квадратично, а в корпусе бывают ошибки единиц и экспериментальные выбросы.
      \[
        L_{\delta}(r)=
        \begin{cases}
          \frac{1}{2}r^2, & |r|\le\delta,\\
          \delta\left(|r|-\frac{1}{2}\delta\right), & |r|>\delta.
        \end{cases}
      \]
      \smallnote{В проекте $\delta=1$: малые ошибки получают сильный квадратичный градиент, крупные не взрывают обучение.}
    \end{column}
    \begin{column}{0.48\textwidth}
      \fitfig{
      \begin{tikzpicture}
        \begin{axis}[
          width=7cm,height=4.5cm,
          xlabel={$r$},ylabel={$L(r)$},
          xmin=-3,xmax=3,ymin=0,ymax=5,
          legend style={font=\tiny,draw=none,fill=white,fill opacity=0.72,text opacity=1},
          legend cell align=left,
          legend pos=north west,
          grid=both
        ]
          \addplot[domain=-3:3,samples=120,TGBlue,thick] {0.5*x*x};
          \addlegendentry{MSE}
          \addplot[domain=-3:3,samples=120,TGOrange,thick] {abs(x)};
          \addlegendentry{MAE}
          \addplot[domain=-3:3,samples=120,TGGreen,thick] {abs(x) <= 1 ? 0.5*x*x : abs(x)-0.5};
          \addlegendentry{Huber}
        \end{axis}
      \end{tikzpicture}}
    \end{column}
  \end{columns}
  \note{Huber — плавный компромисс между MSE около нуля и MAE на выбросах.}
\end{frame}

% E2
\begin{frame}{Супервизия свойств чистого вещества}
  \small
  \begin{columns}[T,onlytextwidth]
    \begin{column}{0.52\textwidth}
      \[
        \mathcal{L}_{T_m}=\frac{1}{N_{T_m}}\sum (T_m^{\text{пред}}-T_m^{\text{ист}})^2,
      \]
      \[
        \mathcal{L}_{\Delta H}=\frac{1}{N_{\Delta H}}\sum(\Delta H^{\text{пред}}-\Delta H^{\text{ист}})^2.
      \]
      \smallnote{Сумма берётся только по строкам, где есть метка: супервизия с маской.}
    \end{column}
    \begin{column}{0.44\textwidth}
      \fitfig{
      \begin{tikzpicture}
        \draw[->] (0,0)--(5.5,0) node[right] {покрытие};
        \fill[TGBlue!65] (0,1.0) rectangle (4.1,1.35);
        \node[anchor=east,font=\scriptsize] at (-0.1,1.17) {$T_m$};
        \node[anchor=west,font=\scriptsize] at (4.2,1.17) {51\%};
        \fill[TGOrange!75] (0,0.25) rectangle (0.08,0.60);
        \node[anchor=east,font=\scriptsize] at (-0.1,0.42) {$\Delta H$};
        \node[anchor=west,font=\scriptsize] at (0.2,0.42) {1\%};
      \end{tikzpicture}}
      \smallnote{Супервизия $T_m$ плотная; $\Delta H_{fus}$ почти всегда учится косвенно через растворимость.}
    \end{column}
  \end{columns}
  \note{Учитель видит ответы для половины примеров по Tm и почти не видит по энтальпии; остальное приходит через SLE-потерю.}
\end{frame}

% E3
\begin{frame}{Единственная прямая супервизия на NRTL}
  \small
  \begin{columns}[T,onlytextwidth]
    \begin{column}{0.52\textwidth}
      \[
        \mathcal{L}_{\gamma^\infty}
        =
        \frac{1}{N_{\gamma^\infty}}\sum
        \left(\ln\gamma_2^{\infty,pred}-\ln\gamma_2^{\infty,true}\right)^2.
      \]
      \[
        \boxed{\ln\gamma_2^\infty=\tau_{12}+\tau_{21}\exp(-\alpha\tau_{21})}
      \]
      \smallnote{IDAC: 404 строки, 138 пар из Zenodo/NIST ThermoML.}
      \vspace{0.15em}
      \begin{block}{Зачем нужен этот якорь}
        \scriptsize
        Loss по растворимости видит только сумму $-\Phi-\ln\gamma_2$. IDAC отдельно привязывает NRTL-параметры к измеряемому пределу активности и уменьшает компенсаторные решения.
      \end{block}
    \end{column}
    \begin{column}{0.44\textwidth}
      \begin{block}{$\gamma_2^\infty$ vs $\gamma_2(x_2)$}
        \scriptsize
        $\gamma_2^\infty$ — вспомогательная целевая величина при бесконечном разбавлении.\\
        $\gamma_2(x_2)$ — коэффициент активности на стороне решателя при конечной концентрации.\\
        При $x_2\to0$ предел решателя должен совпасть с IDAC.
      \end{block}
      \plotasset[width=\linewidth]{nrtl_gamma_example.pdf}
    \end{column}
  \end{columns}
  \note{IDAC — якорь для NRTL-параметров: данных мало, но это прямой сигнал именно на активность.}
\end{frame}

% E4
\begin{frame}{Физические ограничения на форму кривой}
  \small
  \begin{columns}[T,onlytextwidth]
    \begin{column}{0.54\textwidth}
      Монотонность:
      \[
        \mathcal{L}_{mono}=\frac{1}{N}\sum \max\!\left(0,-\frac{\partial\ln x_2}{\partial T}\right)^2.
      \]
      Ранговая потеря:
      \[
        \mathcal{L}_{rank}=\sum\max(0,\ln x_2(T_i)-\ln x_2(T_j)+\epsilon),
        \quad T_i<T_j.
      \]
      Van't Hoff consistency:
      \[
        \mathcal{L}_{vH}=\sum\left(\frac{\Delta \ln x_2}{\Delta(1/T)}-s_{\text{целевой}}\right)^2.
      \]
    \end{column}
    \begin{column}{0.42\textwidth}
      \plotasset[width=\linewidth]{ideal_sle_example.pdf}
      \smallnote{Парно-ориентированная сборка батчей нужна, чтобы в одном батче были несколько температур одной пары.}
    \end{column}
  \end{columns}
  \note{Это мягкие физические штрафы: модель может нарушать их, но платит потерю.}
\end{frame}

% E5
\begin{frame}{Regular Solution bridge: полезен, но опасен}
  \small
  \begin{columns}[T,onlytextwidth]
    \begin{column}{0.52\textwidth}
      Bridge-потеря связывает расстояние Hansen с энергией NRTL:
      \[
        \Delta g_{12}^{Hansen}=V_1[(\delta_{d1}-\delta_{d2})^2+
        \kappa_p(\delta_{p1}-\delta_{p2})^2+
        \kappa_h(\delta_{h1}-\delta_{h2})^2].
      \]
      Проблема:
      \[
        \Delta g_{12}^{Hansen}\ge0 \quad \text{всегда}.
      \]
      Но реальные выгодные контакты могут иметь $\Delta g<0$ и $\gamma_2<1$.
    \end{column}
    \begin{column}{0.44\textwidth}
      \begin{block}{Текущий вывод}
        \scriptsize
        Для систем с H-связями bridge-потеря может тянуть NRTL в неправильную сторону. Поэтому \texttt{bridge\_loss\_weight=0} по умолчанию.
      \end{block}
      \begin{block}{Нейтральная альтернатива}
        \scriptsize
        Walden check:
        \[
          \mathcal{L}_{Walden}=\max(0,|\Delta H/T_m-56.5|-30)^2.
        \]
      \end{block}
    \end{column}
  \end{columns}
  \note{Bridge-потеря — пример физического регуляризатора, который может навредить, если его область применимости уже, чем корпус.}
\end{frame}

% 25
\begin{frame}{Проблема идентифицируемости}
  \small
  Для одной пары при одной температуре:
  \[
    \ln x_2=-\Phi(T_m,\Delta H)-\ln\gamma_2(\tau_{12},\tau_{21},\alpha).
  \]
  \[
    \mathcal{M}=\{\theta\in\mathbb{R}^7:F(\theta;T,y)=0\},
    \qquad
    \dim\mathcal{M}=6.
  \]
  \begin{columns}[T,onlytextwidth]
    \begin{column}{0.56\textwidth}
      \begin{itemize}\tightitems
        \item NRTL может абсорбировать ошибки кристаллических голов.
        \item На обучении это выглядит нормально.
        \item На тестовом разбиении факторизация разваливается.
      \end{itemize}
    \end{column}
    \begin{column}{0.40\textwidth}
      \fitfig{
      \begin{tikzpicture}
        \draw[->] (-2,0)--(2.2,0) node[right] {$\theta_1$};
        \draw[->] (0,-1.4)--(0,1.6) node[above] {$\theta_2$};
        \draw[thick,TGOrange] plot[smooth] coordinates {(-1.7,-1.0)(-0.7,-0.4)(0.2,0.15)(1.3,1.1)};
        \node[font=\scriptsize] at (1.1,-0.8) {$F(\theta)=0$};
      \end{tikzpicture}}
    \end{column}
  \end{columns}
  \note{Переход: нужна подсказка параметрам, которые идут в решатель, особенно когда метки редкие.}
\end{frame}

% 26
\begin{frame}{Подсказка истинного $T_m$ при обучении}
  \small
  \[
    T_m^{\mathrm{solver}}(i)=
    \begin{cases}
      T_m^{\mathrm{true}}, & \text{если метка есть и }u<p(t),\\
      T_m^{\mathrm{pred}}, & \text{иначе}.
    \end{cases}
  \]
  \[
    p(t)=
    \max\!\left(0,\,
    1-\frac{t-t_{\mathrm{start}}}{t_{\mathrm{end}}-t_{\mathrm{start}}}
    \right).
  \]
  \begin{itemize}\tightitems
    \item Фаза 1: $p=1$ — решатель получает правильные кристаллические свойства.
    \item Фаза 2: $p$ убывает — модель учится самостоятельно.
    \item Фаза 3: $p=0$ — чистое предсказание.
  \end{itemize}
  \note{Сравнить с teacher forcing в последовательных моделях, но подчеркнуть, что здесь подсказывается физический параметр.}
\end{frame}

% 27
\begin{frame}{Чувствительность: инверсия иерархии}
  \small
  \begin{tabularx}{\textwidth}{lYY}
    \toprule
    \textbf{Параметр} & \textbf{Чувствительность $\ln x_2$} & \textbf{Прямая супервизия}\\
    \midrule
    $\Delta g_{12}$ & $\pm0.8$ на 2000 Дж/моль & почти нет: 138 IDAC-пар\\
    $\Delta g_{21}$ & $\pm0.8$ & почти нет\\
    $\Delta H_{\mathrm{fus}}$ & $\pm0.25$ & около 1\% строк\\
    $T_m$ & $\pm0.19$ & около 51\% строк\\
    \bottomrule
  \end{tabularx}
  \vspace{0.7em}
  \keybox{Самые важные параметры — наименее супервизированы. Это центральная проблема обучения физического узкого места.}
  \note{Переход к диагностике: теперь объясним, что произошло в результатах.}
\end{frame}

\begin{frame}{Где мы сейчас?}
  \centering
  \vspace{0.15cm}
  \roadmapbar{4}
  \vspace{0.65cm}

  {\fontsize{19}{22}\selectfont
  Мы стабилизировали обучение: поэтапный режим, маскированные вспомогательные потери,
  подстановки истинных параметров и физические регуляризаторы.\par}
  \vspace{0.45cm}
  {\fontsize{21}{24}\selectfont\bfseries\color{TGBlueDark}
  Теперь переходим к результатам:\\
  что работает, а что нет.\par}
  \note{Переход к диагностике: после инженерии обучения нужно не рекламировать модель, а разобраться в фактических ошибках.}
\end{frame}

\section{Диагностика и решения}

\begin{frame}{Как оценивать успех модели?}
  \small
  Ошибка измеряется в логарифме мольной доли:
  \[
    \mathrm{MAE}=\frac1N\sum_i\left|\ln x_{2,i}^{\mathrm{pred}}-\ln x_{2,i}^{\mathrm{true}}\right|.
  \]
  \vspace{0.25em}
  \begin{columns}[T,onlytextwidth]
    \begin{column}{0.32\textwidth}
      \begin{block}{MAE = 0.5}
        \centering
        \Large $\exp(0.5)=1.65$\\[-0.1em]
        \scriptsize ошибка примерно в 1.7 раза
      \end{block}
    \end{column}
    \begin{column}{0.32\textwidth}
      \begin{block}{MAE = 1.0}
        \centering
        \Large $\exp(1)=2.72$\\[-0.1em]
        \scriptsize ошибка примерно в 2.7 раза
      \end{block}
    \end{column}
    \begin{column}{0.32\textwidth}
      \begin{block}{MAE = 2.0}
        \centering
        \Large $\exp(2)=7.39$\\[-0.1em]
        \scriptsize ошибка примерно в 7.4 раза
      \end{block}
    \end{column}
  \end{columns}
  \vspace{0.35em}
  \keybox{\scriptsize Поэтому разрыв $0.09$ MAE — это около 9\% в $x_2$, а разрыв $0.7$ MAE — уже почти двукратный множитель.}
  \note{Перед таблицей результатов объяснить шкалу. Иначе числа MAE 1.65, 1.74 и 1.82 не имеют интуитивного смысла для аудитории.}
\end{frame}

% 28
\begin{frame}{Результаты базовых моделей: короткий бюджет}
  \centering
  \scriptsize
  \setlength{\tabcolsep}{4pt}
  \begin{tabular}{rlccc}
    \toprule
    \textbf{\#} & \textbf{Модель} & \textbf{MAE} & \textbf{RMSE} & \textbf{$R^2$}\\
    \midrule
    1 & \textbf{DirectGNN} & \textbf{1.652} & 2.254 & 0.478\\
    2 & DirectGNN + дескрипторы & 1.668 & 2.242 & 0.483\\
    3 & RF, смешанные признаки & 1.722 & 2.315 & 0.449\\
    4 & TGNN + дескрипторы & 1.729 & 2.288 & 0.461\\
    5 & TGNN MPNN & 1.741 & 2.338 & 0.438\\
    6 & RF, RDKit-дескрипторы & 1.764 & 2.365 & 0.425\\
    7 & Идеальный SLE без обучения & 6.841 & 9.215 & $-7.733$\\
    \bottomrule
  \end{tabular}

  \vspace{0.6em}
  \begin{columns}[T,onlytextwidth]
    \begin{column}{0.48\textwidth}
      \keybox{\scriptsize DirectGNN лидирует с MAE 1.652. TGNN MPNN находится близко к лучшему RF, но жёсткий физический путь пока не даёт выигрыша.}
    \end{column}
    \begin{column}{0.48\textwidth}
      \smallnote{Короткий бюджет: 1/4/1 эпохи, один seed $=42$. RF обучены полностью. Идеальный SLE — чистая физика без обучения.}
    \end{column}
  \end{columns}
  \note{До введения TIMP показываем только базовые модели: прямую GNN, физическую TGNN и RF. Главный вывод: графовая модель уже обошла RF, но физический путь пока стоит небольшую цену.}
\end{frame}

% 29
\begin{frame}{Линейная проба: что знает базовый энкодер}
  \small
  \begin{columns}[T,onlytextwidth]
    \begin{column}{0.46\textwidth}
      \centering
      \scriptsize
      \setlength{\tabcolsep}{5pt}
      \begin{tabular}{lc}
        \toprule
        \textbf{RDKit-дескриптор} & \textbf{$R^2$}\\
        \midrule
        FractionCSP3 & 0.989\\
        NumHDonors & 0.798\\
        TPSA & 0.719\\
        MolLogP & 0.681\\
        MolWt & 0.406\\
        \midrule
        \textbf{Медиана по всем} & \textbf{0.565}\\
        \bottomrule
      \end{tabular}
    \end{column}
    \begin{column}{0.50\textwidth}
      \begin{block}{Что измеряем}
        \scriptsize
        Замораживаем энкодер и обучаем простую линейную регрессию:
        можно ли из скрытого вектора восстановить известные химические
        дескрипторы?

        \vspace{0.35em}
        Высокий $R^2$ означает, что нужная информация присутствует
        в представлении. Низкий $R^2$ означает, что голова модели
        должна угадывать физику из неполного описания.
      \end{block}
      \smallnote{Базовый MPNN хорошо кодирует некоторые локальные признаки, но молекулярная масса и липофильность восстанавливаются хуже.}
    \end{column}
  \end{columns}
  \note{Этот слайд вводит саму методику линейной пробы до обсуждения новых энкодеров. Сравнение с TIMP появится позже, после введения TIMP.}
\end{frame}

% 30
\begin{frame}{Декомпозиция ошибки: обновлённая картина}
  \small
  \[
    \Delta_{\text{физ}}
    =
    \mathrm{MAE}_{\mathrm{TGNN}}-\mathrm{MAE}_{\mathrm{DirectGNN}}
    =
    1.741-1.652=+0.089
  \]
  \[
    \Delta_{\text{энкодер}}
    =
    \mathrm{MAE}_{\mathrm{DirectGNN}}-\mathrm{MAE}_{\mathrm{RF}_{\text{лучший}}}
    =
    1.652-1.722=-0.070
  \]

  \vspace{0.15em}
  \fitfig{
  \begin{tikzpicture}[>=Latex]
    \tikzstyle{model}=[draw,rounded corners=7pt,minimum width=2.4cm,minimum height=1.3cm,align=center,font=\scriptsize]

    \node[model,draw=TGSlate!55,fill=TGSlate!8] (ideal) at (0,0)
      {Идеальный SLE\\[0.1em]{\normalsize 6.84}};
    \node[model,draw=TGGreen!55,fill=TGGreen!8] (rf) at (3.8,0)
      {RF смешанный\\[0.1em]{\normalsize 1.72}};
    \node[model,draw=TGBlue!55,fill=TGBlue!8] (tgnn) at (7.6,0)
      {TGNN MPNN\\[0.1em]{\normalsize 1.74}};
    \node[model,draw=TGOrange!60,fill=TGOrange!8] (direct) at (11.4,0)
      {DirectGNN\\[0.1em]{\normalsize\textbf{1.65}}};

    \draw[->,line width=0.7pt,TGSlate] (ideal)--(rf)
      node[midway,above,font=\tiny] {обучение: $-5.12$};
    \draw[->,line width=0.7pt,TGGreen] (rf)--(tgnn)
      node[midway,above,font=\tiny] {$+0.02$};
    \draw[->,line width=0.7pt,TGOrange] (tgnn)--(direct)
      node[midway,above,font=\tiny,text=TGOrange]
      {физика: $+0.09$};
  \end{tikzpicture}}

  \vspace{0.3em}
  \begin{columns}[T,onlytextwidth]
    \begin{column}{0.48\textwidth}
      \keybox{\scriptsize DirectGNN обошёл лучший RF: $\Delta_{\text{энкодер}}=-0.07$. Цена физического пути: $+0.09$ MAE.}
    \end{column}
    \begin{column}{0.48\textwidth}
      \smallnote{По сравнению с предыдущими прогонами: раньше RF доминировал. Включение водных данных и новое разбиение изменили картину.}
    \end{column}
  \end{columns}
  \note{Нарратив изменился: теперь DirectGNN лучше RF, и основной вопрос — цена физического пути около 0.09 MAE.}
\end{frame}

\begin{frame}{KNN и modelability: прямой ответ на критику}
  \small
  \centering
  \plotasset[height=0.48\textheight]{knn_summary_table}

  \vspace{0.05em}
  \keybox{\scriptsize Full split: 1-NN по pair Tanimoto даёт MAE 2.530, то есть заметно хуже RF 1.722 и DirectGNN 1.652. Subset diagnostic повторяет вывод: даже честный sklearn KNN не превосходит pair-Tanimoto 1-NN.}
  \note{Этот слайд отвечает на риторический тезис про KNN конкретными числами. Важно подчеркнуть, что main evidence здесь -- full-split 1-NN; две нижние строки нужны только как честная проверка, что sklearn KNN не скрывает более сильный результат.}
\end{frame}

\begin{frame}{Почему тезис ``плохой KNN $\Rightarrow$ плохая выборка'' не проходит}
  \small
  \centering
  \plotasset[height=0.52\textheight]{knn_vs_adaptive_benchmarks}

  \vspace{0.05em}
  \keybox{\scriptsize Простая neighbor-based модель сильно слабее RF и DirectGNN. Значит, representation-space содержит полезный сигнал, а наивный nearest-neighbor lookup его не извлекает. Это не буквальный MODI, а эмпирический контраргумент к тезису ``KNN плох, значит dataset плох''.}
  \note{Смысл слайда: RF и DirectGNN обе существенно лучше простого соседа. Если бы выборка была полностью немоделируема в практическом смысле, такого разрыва не было бы.}
\end{frame}

\begin{frame}{Что именно говорят данные о modelability}
  \small
  \centering
  \plotasset[height=0.47\textheight]{knn_modelability_diagnostics}

  \vspace{0.05em}
  \keybox{\scriptsize Локальных cliffs много, но выборка не выглядит хаотичной. При высокой структурной близости согласованность растёт: для pair Tanimoto $\ge0.8$ медианный $|\Delta \ln x_2|$ падает до 0.69 на $n=57$.}
  \smallnote{Оригинальный MODI у Golbraikh et al. определён для бинарной классификации. Здесь показан регрессионный nearest-neighbor аналог: cliffs как функция pair Tanimoto, то есть локальная диагностика modelability, а не буквальный MODI из статьи.}
  \note{Этот слайд нужен, чтобы не скатиться в крайности. Корректная формулировка не ``выборка идеальна'' и не ``выборка бесполезна'', а ``простых соседей недостаточно, при этом у реально близких пар локальная согласованность появляется''.}
\end{frame}

\begin{frame}{Подстановка истинных кристаллических параметров}
  \small
  \centering
  \setlength{\tabcolsep}{5pt}
  \begin{tabular}{lccc}
    \toprule
    \textbf{Модель} & \textbf{MAE обычно} & \textbf{MAE с подсказкой} & \textbf{$\Delta$}\\
    \midrule
    TGNN MPNN & 1.741 & 1.746 & $+0.005$\\
    TGNN + дескрипторы & 1.729 & 1.734 & $+0.005$\\
    \bottomrule
  \end{tabular}

  \vspace{0.5em}
  \keybox{Подстановка истинных $T_m$ и $\Delta H_{\mathrm{fus}}$ не меняет MAE. Узкое место — не в кристаллической голове.}

  \smallnote{Два объяснения: (1) кристаллическая голова уже достаточно хороша; (2) компенсаторная дегенерация — NRTL скомпенсировала ошибки кристаллических параметров. Для различения нужен граф ``предсказано против эксперимента'' для $T_m$.}
  \note{Это диагностический результат: кристаллическую голову можно оставить как есть; фокус переносится на NRTL, решатель и обучение физических голов.}
\end{frame}

% 31
\begin{frame}{Обновлённый диагноз}
  \small
  \begin{columns}[T,onlytextwidth]
    \begin{column}{0.48\textwidth}
      \begin{block}{Что уже ясно}
        \scriptsize
        \begin{itemize}\tightitems
          \item DirectGNN уже обошёл лучший RF: bottleneck базового энкодера больше не доминирует.
          \item Подстановка истинных $T_m$ и $\Delta H_{\mathrm{fus}}$ почти не меняет MAE.
          \item Значит, главный bottleneck сейчас находится не в кристаллической голове.
        \end{itemize}
      \end{block}
      \begin{block}{Что всё ещё мешает}
        \scriptsize
        \begin{itemize}\tightitems
          \item Physics bottleneck даёт измеримую цену: $+0.09$ MAE.
          \item Идентифицируемость параметров и bridge-регуляризация остаются открытыми вопросами.
        \end{itemize}
      \end{block}
    \end{column}
    \begin{column}{0.48\textwidth}
      \begin{block}{Новая рабочая гипотеза}
        \scriptsize
        Проблема сместилась из ветки чистого вещества в ветку
        взаимодействия: физическим головам не хватает
        структурированного описания контактов между молекулами.
      \end{block}
      \begin{block}{Ключевой диагностический сигнал}
        \scriptsize
        Попарное внимание в TGNN получает градиент примерно
        в 33 раза слабее, чем в DirectGNN. Значит, попарное
        представление обучается хуже, чем остальная модель.
      \end{block}
      \keybox{\scriptsize Следующий архитектурный шаг должен менять не solver, а само представление взаимодействий: разделять дисперсионные и полярные контакты ещё на уровне message passing.}
    \end{column}
  \end{columns}
  \note{Этот слайд теперь объединяет и резюме, и диагноз. Он должен считываться за несколько секунд: что уже исключено, что осталось проблемой и почему следующий шаг касается именно ветки взаимодействия.}
\end{frame}

\begin{frame}{Поток градиентов: где теряется сигнал}
  \small
  \begin{columns}[T,onlytextwidth]
    \begin{column}{0.55\textwidth}
      \scriptsize
      \setlength{\tabcolsep}{4pt}
      \begin{tabular}{lccc}
        \toprule
        \textbf{Блок} & \textbf{TGNN} & \textbf{DirectGNN} & \textbf{Отношение}\\
        \midrule
        Энкодер, слой 0 & 0.037 & 0.008 & \textbf{4.9$\times$}\\
        Энкодер, слой 5 & 0.042 & — & —\\
        Попарное внимание & \textbf{0.0007} & \textbf{0.023} & \textbf{0.03$\times$}\\
        Голова кристалла / прямая голова & 1.25 & 0.78 & 1.6$\times$\\
        Адаптер растворителя & $6\cdot10^{-6}$ & $2\cdot10^{-4}$ & 0.03$\times$\\
        \bottomrule
      \end{tabular}
      \vspace{0.45em}
      \smallnote{Норма градиента усреднена по короткому диагностическому прогону.}
    \end{column}
    \begin{column}{0.41\textwidth}
      \begin{block}{Ключевой вывод}
        \scriptsize
        Энкодер TGNN получает сильный сигнал, но в основном через
        кристаллическую голову. Слой взаимодействия, который должен
        учить совместимость вещества и растворителя, получает слишком
        слабый градиент.
      \end{block}
      \keybox{\scriptsize H10: физический путь сохраняет градиент до NRTL-головы,
      но сигнал затухает на участке ``пара молекул $\to$ попарное внимание''.}
    \end{column}
  \end{columns}
  \note{Это уточняет диагноз. Проблема не в том, что градиенты вообще не доходят до энкодера. Они доходят, но по неправильному маршруту: через кристаллическую голову. А взаимодействие вещества и растворителя остаётся недообученным.}
\end{frame}

\begin{frame}{Как оживить слой взаимодействия}
  \small
  \begin{columns}[T,onlytextwidth]
    \begin{column}{0.51\textwidth}
      \begin{block}{1. Вспомогательная прямая ветка}
        \scriptsize
        \[
          \mathbf g_{\mathrm{pair}}
          \to f_{\mathrm{всп}}
          \to \ln x_2^{\mathrm{всп}}
        \]
        \[
          \mathcal L_{\mathrm{всп}}
          =
          \mathrm{Huber}(\ln x_2^{\mathrm{всп}}, \ln x_2^{\mathrm{эксп}})
        \]
        Вес по фазам: $0 \to 0.1 \to 0.01$.
      \end{block}
      \smallnote{Эта ветка нужна только для обучения. В финальном прогнозе остаётся физический путь SLE.}
    \end{column}
    \begin{column}{0.45\textwidth}
      \begin{block}{2. Не дать кристаллу захватить энкодер}
        \scriptsize
        Во второй фазе можно подать в кристаллическую голову
        отсоединённый от градиента вектор:
        \[
          \mathbf g_{\mathrm{sol}}^{crystal}
          =
          \operatorname{sg}(\mathbf g_{\mathrm{sol}})
        \]
        где $\operatorname{sg}$ — оператор отсечения градиента.
        Голова $T_m,\Delta H$ продолжает учиться, но энкодер
        больше не переписывается её сильным градиентом.
      \end{block}
      \keybox{\scriptsize Это проверка H10: если слой взаимодействия оживает,
      должны вырасти градиенты попарного внимания и улучшиться NRTL-параметры.}
    \end{column}
  \end{columns}
  \note{Важная оговорка: это не замена физической модели прямой нейросетью. Прямая ветка не используется как ответ; она только даёт обучающий сигнал слоям взаимодействия. Если она начнёт доминировать, мы потеряем смысл физического bottleneck, поэтому вес малый и фазовый.}
\end{frame}

\begin{frame}{Проверка: вспомогательная ветка оживляет внимание}
  \small
  \begin{columns}[T,onlytextwidth]
    \begin{column}{0.52\textwidth}
      \plotasset[width=\linewidth]{gradient_flow_fix.pdf}
    \end{column}
    \begin{column}{0.44\textwidth}
      \scriptsize
      \setlength{\tabcolsep}{4pt}
      \begin{tabular}{lcc}
        \toprule
        \textbf{Блок} & \textbf{TGNN + ветка} & \textbf{DirectGNN}\\
        \midrule
        Попарное внимание & \textbf{0.0377} & 0.0381\\
        Энкодер, слой 0 & 0.0338 & 0.0099\\
        NRTL-голова & 0.137 & —\\
        Кристаллическая голова & 1.09 & —\\
        \bottomrule
      \end{tabular}
      \vspace{0.4em}
      \keybox{\scriptsize Отношение по попарному вниманию:
      $0.0377/0.0381=0.99$. Слой взаимодействия больше не недополучает
      градиент относительно DirectGNN.}
      \smallnote{Это проверка только потока градиентов на 50 шагах. Она не доказывает,
      что MAE улучшится; это должен проверить отдельный proxy/medium запуск.}
    \end{column}
  \end{columns}
  \note{Показать осторожно: вспомогательная ветка решает именно механическую проблему голодания градиента. Но это ещё не результат по качеству. Следующий эксперимент — сравнить MAE с тем же бюджетом.}
\end{frame}

\begin{frame}{Где мы сейчас?}
  \centering
  \vspace{0.15cm}
  \roadmapbar{5}
  \vspace{0.65cm}

  {\fontsize{18}{22}\selectfont
  Диагностика изменила картину: физический путь стоит мало,
  но ему нужен более информативный вход.\par}
  \vspace{0.45cm}
  {\fontsize{20}{24}\selectfont\bfseries\color{TGBlueDark}
  Цена физического пути сейчас около 0.09 MAE.\\
  Следующий вопрос — как разделить типы атомных взаимодействий.\par}
  \note{Переход к TIMP: теперь вводим физически разделённую передачу сообщений как ответ на диагностический вывод.}
\end{frame}

% 32
\begin{frame}{TIMP: решение в передаче сообщений}
  \small
  \begin{columns}[T,onlytextwidth]
    \begin{column}{0.54\textwidth}
      \molasset[width=\linewidth]{timp_paracetamol_channels.png}
    \end{column}
    \begin{column}{0.42\textwidth}
      \[
        \mathbf{m}_{ij}=
        \underbrace{\phi_{\mathrm{disp}}(...)\sqrt{\alpha_i\alpha_j}}_{\text{\textcolor{TGBlue}{Ван-дер-Ваальс}}}
        +
        \underbrace{\phi_{\mathrm{polar}}(...)\operatorname{sp}(\delta q_i\delta q_j)}_{\text{\textcolor{TGOrange}{полярность/H-связь}}}.
      \]
      \smallnote{Аналогии: Hansen $\delta_t^2=\delta_d^2+\delta_p^2+\delta_h^2$; COSMO-RS разделяет Ван-дер-Ваальс, misfit и H-связи.}
    \end{column}
  \end{columns}
  \note{Показать, что TIMP не заменяет решатель, а улучшает представления, которые решатель получает.}
\end{frame}

\begin{frame}{TIMP на уровне одной связи}
  \small
  \begin{columns}[T,onlytextwidth]
    \begin{column}{0.48\textwidth}
      \centering
      \begin{tikzpicture}[>=Latex]
        \tikzstyle{atom}=[circle,draw,minimum size=0.62cm,font=\scriptsize,fill=white]
        \node[atom,draw=TGBlue] (c1) at (0,2.0) {C};
        \node[atom,draw=TGBlue] (c2) at (2.2,2.0) {C};
        \draw[line width=2.6pt,TGBlue] (c1)--(c2);
        \node[anchor=west,font=\scriptsize,text=TGSlate] at (2.65,2.0) {C--C: дисперсия доминирует};

        \node[atom,draw=TGOrange] (o) at (0,0.75) {O};
        \node[atom,draw=TGOrange] (h) at (2.2,0.75) {H};
        \draw[line width=2.6pt,TGOrange] (o)--(h);
        \node[anchor=west,font=\scriptsize,text=TGSlate] at (2.65,0.75) {O--H: полярность и H-связь};

        \node[atom,draw=TGBlue] (c3) at (0,-0.50) {C};
        \node[atom,draw=TGOrange] (o2) at (2.2,-0.50) {O};
        \draw[line width=1.7pt,TGBlue] (c3) to[bend left=8] (o2);
        \draw[line width=1.7pt,TGOrange] (c3) to[bend right=8] (o2);
        \node[anchor=west,font=\scriptsize,text=TGSlate] at (2.65,-0.50) {C=O: оба канала};
      \end{tikzpicture}
    \end{column}
    \begin{column}{0.48\textwidth}
      \begin{block}{Что меняется относительно MPNN}
        \scriptsize
        Обычный MPNN применяет один тип сообщения ко всем связям. TIMP считает два сообщения:
        \[
          m_{ij}^{disp}\quad\text{и}\quad m_{ij}^{polar}.
        \]
        Затем масштабирует их физическими множителями: поляризуемость усиливает дисперсионный канал, заряды и H-связи усиливают полярный.
      \end{block}
      \begin{block}{Зачем это нужно NRTL}
        \scriptsize
        NRTL-параметры — это эффективные энергии контактов. TIMP заранее разделяет контакты по типам, чтобы NRTL-голове не приходилось извлекать всё из одного смешанного вектора.
      \end{block}
    \end{column}
  \end{columns}
  \note{Это визуальный мост: показать, что TIMP не просто новая формула, а разделение разных химических контактов до предсказания NRTL-параметров.}
\end{frame}

% 33
\begin{frame}{Физические множители TIMP}
  \scriptsize
  \begin{columns}[T,onlytextwidth]
    \begin{column}{0.35\textwidth}
      \molasset[width=0.95\linewidth]{paracetamol.png}
    \end{column}
    \begin{column}{0.31\textwidth}
      \textbf{\textcolor{TGBlue}{Дисперсионный канал}}
      \begin{align*}
        E_{\mathrm{London}}&\propto-\alpha_i\alpha_j/r^6\\
        s_{ij}^{\mathrm{disp}}&\sim\sqrt{\alpha_i\alpha_j}
      \end{align*}
    \end{column}
    \begin{column}{0.31\textwidth}
      \textbf{\textcolor{TGOrange}{Полярный канал}}
      \begin{align*}
        E_{\mathrm{Coulomb}}&\propto\delta q_i\delta q_j/r\\
        s_{ij}^{\mathrm{polar}}&\sim\operatorname{sp}(\delta q_i\delta q_j)
      \end{align*}
    \end{column}
  \end{columns}
  \vspace{0.15em}
  \scriptsize
  \begin{tabularx}{\textwidth}{lccY}
    \toprule
    \textbf{Атом} & $\alpha$ & $\delta q$ & \textbf{Роль}\\
    \midrule
    O гидроксильный & 0.80 & $-0.65$ & полярный акцептор\\
    H гидроксильный & 0.67 & $+0.42$ & полярный донор\\
    C ароматич. & 1.76 & $-0.03$ & дисперсионный\\
    C амидный & 1.76 & $+0.31$ & смешанный\\
    \bottomrule
  \end{tabularx}
  \note{Снять возражение: это не классическая энергия, а физически мотивированное масштабирование обучаемых сообщений.}
\end{frame}

% 34
\begin{frame}{Структурированное представление пары}
  \small
  \[
    \mathbf{g}_{\mathrm{pair}}
    =
    \mathrm{Proj}\!\left(
      [\underbrace{\mathbf{g}^{\mathrm{comb}}}_{\text{общий}}
      \|
      \underbrace{\mathbf{g}^{\mathrm{disp}}}_{\sim\delta_d}
      \|
      \underbrace{\mathbf{g}^{\mathrm{polar}}}_{\sim\delta_p+\delta_h}]
    \right).
  \]
  \fitfig{
  \begin{tikzpicture}
    \node[anchor=east] at (0,1.0) {MPNN $\to$ NRTL};
    \fill[TGSlate!45] (0.2,0.8) rectangle (4.8,1.2);
    \node[anchor=west] at (5.0,1.0) {обобщённый $\to \tau_{12},\tau_{21},\alpha$};
    \node[anchor=east] at (0,0.0) {TIMP $\to$ NRTL};
    \fill[TGBlue!60] (0.2,-0.2) rectangle (1.8,0.2);
    \fill[TGOrange!65] (1.8,-0.2) rectangle (3.4,0.2);
    \fill[TGGreen!55] (3.4,-0.2) rectangle (4.8,0.2);
    \node[anchor=west] at (5.0,0.0) {структурированный $\to \tau_{12},\tau_{21},\alpha$};
  \end{tikzpicture}}
  \smallnote{Дополнительная супервизия каналов: $\mathbf{g}^{disp}\to\delta_d$, $\mathbf{g}^{polar}\to\delta_p,\delta_h$.}
  \note{Переход: TIMP — основной фикс, добавление дескрипторов — страховка, если GNN всё ещё теряет информацию.}
\end{frame}

\begin{frame}{Доступная поверхность: какие атомы видит растворитель?}
  \small
  \begin{columns}[c,onlytextwidth]
    \begin{column}{0.45\textwidth}
      \centering
      \begin{tikzpicture}[scale=1.00,>=Latex]
        \tikzstyle{atom}=[circle,draw,fill=white,minimum size=0.48cm,font=\scriptsize\bfseries]
        \draw[TGBlue!20,line width=13pt,line cap=round]
          (-1.4,0.55) -- (-0.45,1.05) -- (0.55,0.55) -- (0.58,-0.55)
          -- (-0.32,-1.10) -- (-1.28,-0.55) -- cycle;
        \draw[TGOrange!24,line width=11pt,line cap=round] (0.58,-0.55) -- (1.52,-1.15);
        \draw[TGGreen!20,line width=10pt,line cap=round] (-0.45,1.05) -- (-0.55,1.95);

        \node[atom,draw=TGBlue] (c1) at (-1.4,0.55) {C};
        \node[atom,draw=TGBlue] (c2) at (-0.45,1.05) {C};
        \node[atom,draw=TGBlue] (c3) at (0.55,0.55) {C};
        \node[atom,draw=TGBlue] (c4) at (0.58,-0.55) {C};
        \node[atom,draw=TGBlue] (c5) at (-0.32,-1.10) {C};
        \node[atom,draw=TGBlue] (c6) at (-1.28,-0.55) {C};
        \node[atom,draw=TGOrange,fill=TGOrange!10] (o) at (1.52,-1.15) {O};
        \node[atom,draw=TGGreen,fill=TGGreen!10] (n) at (-0.55,1.95) {N};
        \draw[line width=1.2pt] (c1)--(c2)--(c3)--(c4)--(c5)--(c6)--(c1);
        \draw[line width=1.2pt] (c4)--(o);
        \draw[line width=1.2pt] (c2)--(n);

        \draw[->,line width=1.2pt,TGOrange] (2.35,1.25) -- (o);
        \draw[->,line width=1.2pt,TGGreen!70!black] (2.35,1.25) -- (n);
        \node[draw=TGOrange!35,fill=TGOrange!8,rounded corners=5pt,
          align=center,font=\scriptsize,text width=2.35cm] at (2.55,1.65)
          {большой вес:\\доступные O/N};
        \node[draw=TGBlue!30,fill=TGBlue!7,rounded corners=5pt,
          align=center,font=\scriptsize,text width=2.55cm] at (-1.45,-2.0)
          {меньший вес:\\экранированные атомы};
      \end{tikzpicture}
    \end{column}
    \begin{column}{0.51\textwidth}
      \textbf{Идея:} растворитель взаимодействует не со всеми атомами одинаково,
      а прежде всего с доступной поверхностью молекулы.
      \vspace{0.5em}
      \[
        w_i^{\mathrm{SASA}}
        =
        \frac{S_i}{\sum_j S_j+\varepsilon},
        \qquad
        S_i\ge0
      \]
      \vspace{-0.5em}
      \[
        \mathbf g_{\mathrm{disp}}
        =
        \sum_i w_i^{\mathrm{SASA}}\mathbf a_i^{\mathrm{disp}},
        \qquad
        \mathbf g_{\mathrm{polar}}
        =
        \sum_i w_i^{\mathrm{SASA}}\mathbf a_i^{\mathrm{polar}}.
      \]
      \begin{block}{Почему это важно}
        \scriptsize
        Каналы TIMP говорят, \emph{какой тип взаимодействия} несёт атом.
        SASA-взвешивание говорит, \emph{насколько этот атом доступен} для контакта
        с растворителем.
      \end{block}
      \smallnote{$S_i$ — вклад атома в доступную растворителю площадь поверхности (SASA).
      Это опциональное расширение выходного блока: базовая сборка использует
      канально-раздельную агрегацию вниманием, а SASA-вариант требует устойчивой
      оценки $S_i$.}
    \end{column}
  \end{columns}
  \note{Здесь важно не переобещать: явная агрегация по доступной поверхности — опциональный вариант выходного блока TIMP. Смысл в том, что энергии контакта должны сильнее зависеть от поверхностно доступных атомов, а не от глубоко экранированных фрагментов.}
\end{frame}

% I1
\begin{frame}{Три силы, два канала: почему так?}
  \small
  \begin{tabularx}{\textwidth}{lYY}
    \toprule
    \textbf{Сила} & \textbf{Физика} & \textbf{Отображение в TIMP}\\
    \midrule
    London / Ван-дер-Ваальс & между всеми атомами; $E\sim-\alpha_i\alpha_j/r^6$; всегда притяжение & дисперсионный канал, Hansen $\delta_d$\\
    полярные & заряды и диполи; $E\sim q_iq_j/r$; притяжение или отталкивание & полярный канал, Hansen $\delta_p$\\
    H-связи & донор–акцептор, локально и направленно, $\sim20$ кДж/моль & полярный канал, Hansen $\delta_h$\\
    \bottomrule
  \end{tabularx}
  \vspace{0.4em}
  \keybox{Два канала вместо трёх: полярность и H-связи часто коррелируют, а три раздельных канала сложнее идентифицировать на разреженных вспомогательных метках.}
  \note{Разделение на два канала — инженерный компромисс между физической декомпозицией и идентифицируемостью.}
\end{frame}

% I2
\begin{frame}{4 новых числа для каждой химической связи}
  \small
  \begin{columns}[T,onlytextwidth]
    \begin{column}{0.52\textwidth}
      Стандартные признаки связей: 8 чисел. TIMP может добавить:
      \[
        \Delta\chi=|\chi_i-\chi_j|,\qquad
        \Delta r_{vdW}=|r_i-r_j|,
      \]
      \[
        \mu_{bond}\approx\Delta\chi\,(r_i+r_j),\qquad
        h_{cap}\in\{0,1\}.
      \]
      Итого $e_{ij}\in\mathbb{R}^{12}$.
    \end{column}
    \begin{column}{0.44\textwidth}
      \scriptsize
      \begin{tabular}{lccc}
        \toprule
        \textbf{Связь} & $\Delta\chi$ & $\mu_{bond}$ & $h_{cap}$\\
        \midrule
        C–C & 0.0 & 0.0 & 0\\
        C=O & $\sim0.9$ & $\sim2.4$ & 1\\
        O–H & $\sim1.4$ & $\sim2.8$ & 1\\
        C–Br & $\sim0.4$ & асимм. & 0\\
        \bottomrule
      \end{tabular}
    \end{column}
  \end{columns}
  \par\vspace{0.45em}
  \begin{minipage}{0.92\textwidth}
    \scriptsize\color{TGSlate}
    $h_{cap}=1$, если связь указывает на способность быть донором или акцептором H-связи.
  \end{minipage}
  \note{Эти признаки не заменяют обучаемые сообщения, а дают им физически осмысленные скаляры на ребре.}
\end{frame}

% I3
\begin{frame}{Парциальные заряды: как распределена электронная плотность}
  \small
  \begin{columns}[T,onlytextwidth]
    \begin{column}{0.42\textwidth}
      \plotasset[width=\linewidth]{timp_dual_channel_molecule.pdf}
    \end{column}
    \begin{column}{0.54\textwidth}
      Заряды Gasteiger–Marsili:
      \begin{itemize}\tightitems
        \item итеративная электроотрицательностная схема;
        \item не требует квантовой химии;
        \item обычно $<1$ мс на молекулу;
        \item сумма зарядов нейтральной молекулы $\approx0$.
      \end{itemize}
      \scriptsize
      \begin{tabular}{lcc}
        \toprule
        Атом & заряд & роль\\
        \midrule
        O гидроксильный & $-0.65$ & акцептор\\
        H гидроксильный & $+0.42$ & донор\\
        N амидный & $-0.36$ & полярный\\
        C карбонильный & $+0.31$ & смешанный\\
        \bottomrule
      \end{tabular}
    \end{column}
  \end{columns}
  \note{Заряды приближённые, но достаточны как индуктивная подсказка; сама модель может доучить эффективные взаимодействия.}
\end{frame}

% I4
\begin{frame}{Подсказка для внимания: ``смотри на выгодные контакты''}
  \small
  Стандартное внимание:
  \[
    \alpha_{ij}=\mathrm{softmax}\left(\frac{\mathbf{q}_i^\top\mathbf{k}_j}{\sqrt{d}}\right).
  \]
  Термодинамическое смещение:
  \[
    \alpha_{ij}=\mathrm{softmax}\left(\frac{\mathbf{q}_i^\top\mathbf{k}_j}{\sqrt{d}}+\beta E_{ij}^{contact}\right),
  \]
  \[
    E_{ij}^{contact}=\mathrm{MLP}_{energy}([\mathbf{h}_i\|\mathbf{h}_j])\,a_i\,a_j,
    \qquad
    a_i=\sigma(\mathrm{MLP}_{acc}(\mathbf{h}_i)).
  \]
  \begin{columns}[T,onlytextwidth]
    \begin{column}{0.48\textwidth}\smallnote{$\beta\to0$: стандартное внимание.}\end{column}
    \begin{column}{0.48\textwidth}\smallnote{$\beta$ растёт: смещение выделяет физически выгодные контакты.}\end{column}
  \end{columns}
  \note{Термодинамическое перекрёстное внимание — опциональный режим: оно не меняет базовую модель, а добавляет физическое смещение в логиты внимания.}
\end{frame}

% I5
\begin{frame}{Каналы не должны ``схлопнуться''}
  \small
  Без вспомогательного сигнала оба TIMP-канала могут выучить одно и то же. Поэтому добавлены мягкие Hansen-пробы:
  \[
    \mathcal{L}_{disp}=\mathrm{MSE}(\mathrm{MLP}(\mathbf{g}_{sol}^{disp}),\delta_d^{true}),
  \]
  \[
    \mathcal{L}_{polar}=
    \mathrm{MSE}(\mathrm{MLP}(\mathbf{g}_{sol}^{polar}),\delta_p^{true})
    +
    \mathrm{MSE}(\mathrm{MLP}(\mathbf{g}_{sol}^{polar}),\delta_h^{true}).
  \]
  \begin{block}{Роль}
    Малые веса 0.02–0.05 и разреженное покрытие 0.85\%: это структурная регуляризация, а не основная целевая супервизия.
  \end{block}
  \note{Супервизия каналов нужна не для точного предсказания Hansen, а чтобы сохранить физический смысл декомпозиции каналов.}
\end{frame}

% I6
\begin{frame}{Hansen-contrastive: сохраняем геометрию совместимости}
  \small
  \begin{columns}[T,onlytextwidth]
    \begin{column}{0.58\textwidth}
      \centering
      \begin{tikzpicture}[scale=0.93, every node/.style={font=\scriptsize}]
        \node[
          draw=TGBlue!28,
          fill=TGBlue!6,
          rounded corners=8pt,
          minimum width=3.95cm,
          minimum height=3.15cm,
          label={[font=\scriptsize\bfseries,text=TGBlueDark]above:Hansen-пространство}
        ] (hbox) at (0,0) {};
        \node[
          draw=TGOrange!35,
          fill=TGOrange!8,
          rounded corners=8pt,
          minimum width=3.95cm,
          minimum height=3.15cm,
          label={[font=\scriptsize\bfseries,text=TGOrange!85!black]above:latent-пространство GNN}
        ] (zbox) at (5.2,0) {};

        \coordinate (hA) at (-1.25,0.85);
        \coordinate (hB) at (0.55,0.45);
        \coordinate (hC) at (1.05,-0.95);
        \coordinate (zA) at (4.15,0.95);
        \coordinate (zB) at (5.55,0.55);
        \coordinate (zC) at (6.25,-0.80);

        \draw[TGBlue!45,dashed,line width=0.55pt] (hA) -- node[above,font=\tiny,text=TGSlate] {$R_a$ мал.} (hB);
        \draw[TGBlue!35,dashed,line width=0.55pt] (hA) -- node[left,font=\tiny,text=TGSlate] {$R_a$ бол.} (hC);
        \draw[TGOrange!55,dashed,line width=0.55pt] (zA) -- node[above,font=\tiny,text=TGSlate] {$\|\mathbf g_i-\mathbf g_j\|$ мал.} (zB);
        \draw[TGOrange!45,dashed,line width=0.55pt] (zA) -- node[right,font=\tiny,text=TGSlate] {далеко} (zC);

        \foreach \p/\lab/\col in {hA/A/TGBlue,hB/B/TGBlue,hC/C/TGBlue,zA/A/TGOrange,zB/B/TGOrange,zC/C/TGOrange}{
          \fill[\col!70] (\p) circle (3.7pt);
          \node[white,font=\tiny\bfseries] at (\p) {\lab};
        }

        \draw[-{Latex[length=3mm]},line width=0.9pt,TGSlate!70]
          (2.15,0.25) -- node[above,font=\tiny,text=TGSlate,align=center]
          {учим сохранять\\расстояния} (3.05,0.25);

        \node[align=center,text=TGSlate,font=\tiny] at (0,-1.92)
          {$R_a^2=4(\Delta\delta_d)^2+(\Delta\delta_p)^2+(\Delta\delta_h)^2$};
        \node[align=center,text=TGSlate,font=\tiny] at (5.2,-1.92)
          {$\|\mathbf g_i-\mathbf g_j\|_2\approx\alpha R_a+\beta$};
      \end{tikzpicture}
    \end{column}
    \begin{column}{0.39\textwidth}
      \begin{block}{Три уровня}
        \scriptsize
        \begin{itemize}\tightitems
          \item \textbf{Молекулярный}: близкие по Hansen молекулы должны иметь близкие эмбеддинги.
          \item \textbf{Канальный TIMP}: $\mathbf g^{disp}\leftrightarrow\delta_d$, $\mathbf g^{polar}\leftrightarrow(\delta_p,\delta_h)$.
          \item \textbf{Парный}: пары с похожей совместимостью $R_a$ должны быть близки в $\mathbf g_{pair}$.
        \end{itemize}
      \end{block}
      \vspace{0.25em}
      \[
        \mathcal L_{HC}=
        \mathrm{MSE}\!\left(
        \|\mathbf g_i-\mathbf g_j\|_2,\;
        \mathrm{sg}(\alpha R_{ij}+\beta)
        \right)
      \]
      \smallnote{Экспериментальные Hansen-метки имеют вес 1.0; pseudo-Hansen из RDKit-дескрипторов — вес 0.3.}
    \end{column}
  \end{columns}
  \note{Объяснить идею без деталей оптимизации: мы не заставляем сеть точно предсказывать Hansen, а просим сохранить геометрию совместимости. Если две молекулы близки по дисперсии, полярности и H-связям, то их скрытые представления тоже должны быть близки. Для TIMP это дополнительно удерживает дисперсионный и полярный каналы от схлопывания.}
\end{frame}

% I7
\begin{frame}{Линейная проба после TIMP}
  \small
  \begin{columns}[T,onlytextwidth]
    \begin{column}{0.54\textwidth}
      \centering
      \scriptsize
      \setlength{\tabcolsep}{3pt}
      \begin{tabular}{lccc}
        \toprule
        \textbf{Дескриптор} & \textbf{MPNN} & \textbf{TIMP} & \textbf{TIMP+HC}\\
        \midrule
        FractionCSP3 & 0.989 & 0.979 & \textbf{0.992}\\
        NumHDonors & 0.798 & 0.790 & \textbf{0.803}\\
        TPSA & 0.719 & 0.718 & \textbf{0.730}\\
        MolLogP & 0.681 & 0.722 & \textbf{0.751}\\
        MolWt & 0.406 & 0.466 & \textbf{0.482}\\
        \midrule
        \textbf{Медиана $R^2$} & 0.565 & 0.575 & \textbf{0.601}\\
        $R^2 \ge 0.8$ & 10 & 17 & \textbf{19}\\
        $R^2 < 0.5$ & 84 & 85 & \textbf{76}\\
        \bottomrule
      \end{tabular}
    \end{column}
    \begin{column}{0.42\textwidth}
      \begin{block}{Что изменилось}
        \scriptsize
        TIMP и Hansen-contrastive улучшают скрытое представление:
        больше RDKit-дескрипторов восстанавливаются линейной моделью.

        \vspace{0.35em}
        Особенно заметны MolLogP и MolWt — признаки, связанные с
        дисперсионным вкладом и размером молекулы.
      \end{block}
      \smallnote{Это диагностика представлений, а не финальная ошибка растворимости.}
    \end{column}
  \end{columns}
  \note{Теперь TIMP уже введён, поэтому можно показать сравнительную линейную пробу. Смысл: представления стали богаче, но это ещё нужно превратить в выигрыш по MAE.}
\end{frame}

% I8
\begin{frame}{TIMP + HC: каналы разделились}
  \small
  \begin{columns}[T,onlytextwidth]
    \begin{column}{0.48\textwidth}
      \centering
      \setlength{\tabcolsep}{4pt}
      \scriptsize
      \begin{tabular}{lcc}
        \toprule
        \textbf{Метрика} & \textbf{TIMP} & \textbf{TIMP+HC}\\
        \midrule
        Косинусная близость & 0.058 & \textbf{0.006}\\
        Норма дисперс. канала & 24.5 & 15.3\\
        Норма полярн. канала & 23.9 & \textbf{38.2}\\
        Доля дисперс. канала & 0.51 & \textbf{0.29}\\
        \bottomrule
      \end{tabular}

      \vspace{0.45em}
      \begin{block}{Интерпретация}
        \scriptsize
        Hansen-contrastive сделал каналы практически независимыми:
        cosine $0.006$ близок к нулю. Полярный канал стал доминирующим,
        что согласуется с корпусом, где среди частых растворителей много
        спиртов и воды.
      \end{block}
    \end{column}
    \begin{column}{0.48\textwidth}
      \fitfig{
      \begin{tikzpicture}
        \begin{axis}[
          ybar,
          width=6.7cm,
          height=5.0cm,
          xtick={1,2},
          xticklabels={TIMP,TIMP+HC},
          xmin=0.45,
          xmax=2.55,
          ylabel={норма канала},
          legend style={font=\tiny,at={(0.02,0.98)},anchor=north west,draw=none,fill=white,fill opacity=0.82,text opacity=1},
          bar width=12pt,
          ymin=0,
          ymax=45,
          nodes near coords,
          nodes near coords align={vertical},
          every node near coord/.append style={font=\tiny},
          grid=major,
          grid style={TGSlate!10},
        ]
        \addplot[fill=TGBlue!70,draw=TGBlue!80] coordinates {(1,24.5)(2,15.3)};
        \addlegendentry{дисперс.}
        \addplot[fill=TGOrange!70,draw=TGOrange!80] coordinates {(1,23.9)(2,38.2)};
        \addlegendentry{полярн.}
        \end{axis}
      \end{tikzpicture}}
      \smallnote{Порог схлопывания каналов: cosine $>0.7$.
      Оба варианта далеко от схлопывания.}
    \end{column}
  \end{columns}
  \note{Показать главный факт: HC не просто добавил ещё одну потерю, а изменил геометрию каналов. Каналы стали почти ортогональными, значит они несут разные типы информации.}
\end{frame}

% I8
% I9
\begin{frame}{Парадокс: лучший энкодер, худшая MAE}
  \small
  \begin{columns}[T,onlytextwidth]
    \begin{column}{0.48\textwidth}
      \begin{block}{Что улучшилось в представлении}
        \scriptsize
        \begin{itemize}\tightitems
          \item Медиана линейной пробы $R^2$: $0.565 \to 0.601$.
          \item Линейная проба MolLogP: $0.681 \to 0.751$.
          \item Дескрипторов с $R^2\ge 0.8$: $10 \to 19$.
          \item Косинус каналов: $0.058 \to 0.006$.
        \end{itemize}
      \end{block}
      \begin{block}{Что ухудшилось в ответе}
        \scriptsize
        \begin{itemize}\tightitems
          \item MAE: $1.741 \to 1.820$.
          \item $R^2$: $0.438 \to 0.384$.
        \end{itemize}
      \end{block}
    \end{column}
    \begin{column}{0.48\textwidth}
      \begin{block}{Гипотеза H9}
        \scriptsize
        TIMP содержит больше химической информации, но голова NRTL
        не успевает адаптироваться к структурированному входу
        за 4 эпохи основной фазы.

        \vspace{0.35em}
        MPNN даёт ``серое'' представление. TIMP даёт ``раскрашенное''
        представление: дисперсия + полярность. Голова NRTL должна
        научиться использовать эту структуру.
      \end{block}
      \vspace{0.25em}
      \keybox{\scriptsize Проверка: средний бюджет 10/40/10.
      Если MAE TIMP снижается при стабильном $R^2$ линейной пробы,
      значит ограничение было в обучении после энкодера.}
    \end{column}
  \end{columns}
  \note{Здесь важно не продать TIMP как готовую победу. Корректный вывод тоньше: представление стало лучше, но конечная MAE пока хуже. Это полезная диагностика, потому что показывает следующий bottleneck.}
\end{frame}

% 35
\begin{frame}{Добавление дескрипторов как страховка}
  \small
  \[
    \mathbf{g}_{\mathrm{pair}}
    =
    \mathrm{Proj}([\mathbf{g}_{\mathrm{pair}}^{\mathrm{GNN}}\|
    \mathbf{g}_{\mathrm{pair}}^{\mathrm{descr}}]).
  \]
  \begin{itemize}\tightitems
    \item RDKit: MolLogP, TPSA, MolWt, HBD/HBA, отпечаток Morgan, топологические дескрипторы.
    \item Дескрипторы нормируются по статистикам обучающей выборки; некорректные численные значения заменяются безопасными.
    \item Физический путь не меняется: головы $\to$ SLE $\to \ln x_2$.
  \end{itemize}
  \keybox{Дескрипторы обогащают вход для голов, но не обходят физический путь.}
  \note{Переход: после обновлений представления возвращаемся к протоколу сравнения и результатам.}
\end{frame}

\section{Результаты и приложения}

% 37
\begin{frame}{Как вспомогательные метки входят в обучение}
  \small
  \begin{tabularx}{\textwidth}{lYY}
    \toprule
    \textbf{Источник} & \textbf{Покрытие} & \textbf{Роль}\\
    \midrule
    $T_m$ Bradley & около 51\% строк & супервизия кристаллической головы\\
    $\Delta H_{\mathrm{fus}}$ NIST & около 1\% & разреженная супервизия кристалла\\
    Hansen & около 0.85\% & супервизия каналов TIMP\\
    $\gamma^\infty$ IDAC, Zenodo & 404 строки, 138 пар & супервизия NRTL\\
    \bottomrule
  \end{tabularx}
  \vspace{0.5em}
  \begin{block}{Важно}
    $\gamma_2^\infty$ — IDAC-метка при бесконечном разбавлении; $\gamma_2(x_2)$ — коэффициент активности, который решатель вычисляет при конечном составе.
  \end{block}
  \note{Это не первое введение IDAC, а возврат к вопросу обучения: какие метки за какие головы отвечают. IDAC — единственная прямая супервизия на NRTL-параметры, поэтому она важна непропорционально размеру.}
\end{frame}

% 38
\begin{frame}{Разбиения и протокол}
  \begin{columns}[T,onlytextwidth]
    \begin{column}{0.52\textwidth}
      \begin{itemize}\tightitems
        \item Каноническое: scaffold-разбиение 80/10/10.
        \item Дополнительные: отложенная проверка по веществам и растворителям.
        \item Ни один Murcko scaffold из теста не встречается в обучении.
        \item Seeds: 42, 123, 456.
      \end{itemize}
    \end{column}
    \begin{column}{0.44\textwidth}
      \begin{block}{Честное сравнение}
        Настроенный TGNN, настроенный DirectGNN и RF: одинаковые данные и бюджет.
      \end{block}
    \end{column}
  \end{columns}
  \note{Переход: результаты пока предварительные; ключевое — запланирован полный бюджет.}
\end{frame}

% 39
\begin{frame}{Текущие результаты и ожидания}
  \small
  \begin{tabularx}{\textwidth}{lYY}
    \toprule
    \textbf{Семейство} & \textbf{Текущее наблюдение} & \textbf{Ожидание полного бюджета}\\
    \midrule
    RF-дескрипторы & сильная одиночная базовая модель & останется жёстким ориентиром\\
    DirectGNN, настроенная & лучше TGNN на коротком/среднем бюджете & должен показать потолок энкодера без физики\\
    TGNN-Solv, настроенная & после фикса потерь близок к DirectGNN & проверить, исчезает ли разрыв при 50/200/50\\
    TIMP & реализован как фикс из гипотез & должен улучшить пробы и вход NRTL-головы\\
    \bottomrule
  \end{tabularx}
  \vspace{0.5em}
  \smallnote{Числа на слайдах результатов помечаются как предварительные, пока не закрыты 3 случайные инициализации полного бюджета.}
  \note{Переход к плану экспериментов: как именно проверять гипотезы.}
\end{frame}

% H1
\begin{frame}{Не все предсказания одинаково надёжны}
  \small
  \begin{columns}[T,onlytextwidth]
    \begin{column}{0.46\textwidth}
      Модель всегда выдаёт число, но для одних систем она интерполирует, а для других угадывает.
      \begin{block}{Два вопроса}
        \scriptsize
        1. Насколько надёжен конкретный прогноз? $\to$ неопределённость.\\
        2. Система похожа на обучающую область? $\to$ OOD / область применимости.
      \end{block}
    \end{column}
    \begin{column}{0.50\textwidth}
      \fitfig{
      \begin{tikzpicture}
        \begin{axis}[width=7cm,height=4.6cm,xlabel={true},ylabel={pred},xmin=-12,xmax=0,ymin=-12,ymax=0,grid=both]
          \addplot[only marks,mark=*,mark size=1.5pt,TGBlue!70] coordinates {(-10,-9.8)(-8,-8.2)(-6,-6.1)(-4,-3.8)(-2,-2.2)(-7,-5.2)(-3,-5.7)};
          \addplot[TGSlate,domain=-12:0] {x};
        \end{axis}
      \end{tikzpicture}}
      \smallnote{Идея: точки с высокой неопределённостью должны чаще лежать дальше от диагонали.}
    \end{column}
  \end{columns}
  \note{Важно объяснить, что неопределённость — часть поддержки принятия решений: не только ответ, но и доверие к ответу.}
\end{frame}

% H2
\begin{frame}{MC-Dropout: 30 стохастических предсказаний}
  \small
  Dropout — это случайное выключение части нейронов при обучении.
  Обычно на предсказании его отключают. MC-Dropout оставляет его включённым
  и делает $N=30$ проходов:
  \[
    \hat{y}_{mean}=\frac{1}{N}\sum_n \hat{y}^{(n)},\qquad
    \sigma=\sqrt{\frac{1}{N-1}\sum_n(\hat{y}^{(n)}-\hat{y}_{mean})^2}.
  \]
  \begin{columns}[T,onlytextwidth]
    \begin{column}{0.48\textwidth}
      \begin{block}{Интерпретация}
        \scriptsize
        Малое $\sigma$: прогноз стабилен.\\
        Большое $\sigma$: модель чувствительна к случайной маске, значит менее уверена.
      \end{block}
    \end{column}
    \begin{column}{0.48\textwidth}
      \begin{block}{Калибровка}
        \scriptsize
        PICP$_{90}$ — доля точек внутри 90\% интервала. Идеально: PICP$_{90}\approx0.90$; ниже — интервалы переуверенные.
      \end{block}
    \end{column}
  \end{columns}
  \note{MC-Dropout дешёвый и работает поверх существующей модели: меняется режим предсказания, а не архитектура.}
\end{frame}

% H3
\begin{frame}{Внутри области применимости или вне её}
  \scriptsize
  \begin{columns}[T,onlytextwidth]
    \begin{column}{0.52\textwidth}
      Расстояние в пространстве представлений:
      \[
        D_M=\sqrt{(\mathbf{g}-\mu)^\top\Sigma^{-1}(\mathbf{g}-\mu)}.
      \]
      Молекулярное сходство:
      \[
        T_{\text{вещ}}=\max_j Tanimoto(fp_{\text{вещ}},fp_j^{\text{обуч}}),
      \]
      \[
        T_{\text{раств}}=\max_j Tanimoto(fp_{\text{раств}},fp_j^{\text{обуч}}).
      \]
      Внутри области применимости: $D_M$ ниже обучающего перцентиля, а обе оценки Tanimoto выше порога.
    \end{column}
    \begin{column}{0.44\textwidth}
      \fitfig{
      \begin{tikzpicture}
        \begin{axis}[width=6.5cm,height=4.5cm,xtick=\empty,ytick=\empty,xlabel={координата 1},ylabel={координата 2},grid=both]
          \addplot[only marks,TGGreen,mark=*,mark size=2pt] coordinates {(0,0)(0.4,0.2)(-0.3,0.1)(0.1,-0.4)(0.5,-0.3)(-0.5,-0.2)};
          \addplot[only marks,TGRed,mark=*,mark size=2.4pt] coordinates {(2.0,1.8)(-2.0,1.5)(1.9,-1.4)};
        \end{axis}
      \end{tikzpicture}}
      \smallnote{Зелёные — похожи на обучение; красные — OOD, требуются осторожность и эксперименты.}
    \end{column}
  \end{columns}
  \note{Область применимости нужна, чтобы не применять модель за пределами химического пространства, где она обучалась.}
\end{frame}

% 43
\begin{frame}{Что прикладной слой не делает}
  \small
  \begin{columns}[T,onlytextwidth]
    \begin{column}{0.49\textwidth}
      \begin{block}{Не заявляем}
        \scriptsize
        \begin{itemize}\tightitems
          \item кинетику кристаллизации и рост кристаллов;
          \item полиморфизм, сольваты и механизмы зародышеобразования;
          \item полноценную PBPK/PK/PD-модель;
          \item генерацию маршрута синтеза.
        \end{itemize}
      \end{block}
    \end{column}
    \begin{column}{0.49\textwidth}
      \begin{block}{Что действительно делаем}
        \scriptsize
        \begin{itemize}\tightitems
          \item считаем равновесную растворимость $x_2(T)$;
          \item ранжируем растворители и температурные окна;
          \item строим эвристики поверх $x_2(T)$;
          \item явно показываем, где нужна экспериментальная проверка.
        \end{itemize}
      \end{block}
    \end{column}
  \end{columns}
  \vspace{0.35em}
  \keybox{\scriptsize Прикладные модули — decision-support поверх растворимости, а не замена лабораторного процесса или регуляторного доказательства.}
  \note{Этот слайд нужен как защита от чрезмерных claims: приложение полезно, но не должно звучать как ретросинтез, PBPK или кинетическая модель кристаллизации.}
\end{frame}

\begin{frame}{Прикладной слой}
  \begin{columns}[T,onlytextwidth]
    \begin{column}{0.50\textwidth}
      \begin{enumerate}\tightitems
        \item Отбор растворителей.
        \item Оптимизация процесса.
        \item BCS-подобная оценка пригодности кандидата.
        \item PK-профиль растворимости для сред ЖКТ.
      \end{enumerate}
    \end{column}
    \begin{column}{0.46\textwidth}
      \begin{block}{Принцип}
        Это надстройка над $x_2(T)$, а не отдельные модели.
      \end{block}
      \smallnote{Поддерживается для поверхностей предсказания TGNN-Solv и DirectGNN.}
    \end{column}
  \end{columns}
  \note{Показать связь Streamlit-приложений с основной исследовательской моделью.}
\end{frame}

% G1
\begin{frame}{4 приложения, 1 базовый сервис}
  \small
  \begin{center}
  \begin{tikzpicture}[>=Latex]
    \tikzstyle{box}=[draw,rounded corners=5pt,fill=white,minimum height=0.55cm,align=center,font=\scriptsize]
    \node[box,draw=TGOrange,minimum width=3.0cm] at (0,2.4) {BCS-оценка};
    \node[box,draw=TGBlue,minimum width=3.6cm] at (-2.2,1.55) {отбор растворителей};
    \node[box,draw=TGGreen,minimum width=3.6cm] at (2.2,1.55) {PK-профиль};
    \node[box,draw=TGPurple,minimum width=4.2cm] at (0,0.75) {оптимизация процесса};
    \node[box,draw=TGRed,fill=TGRed!8,minimum width=8.5cm] at (0,-0.15) {предсказание растворимости $x_2(T)$};
    \node[box,draw=TGSlate,fill=TGSlate!5,minimum width=8.5cm] at (0,-1.05) {TGNN-Solv или DirectGNN};
  \end{tikzpicture}
  \end{center}
  \keybox{Верхние рабочие сценарии — эвристики поверх одного прогноза растворимости, а не отдельные модели.}
  \note{Сформулировать базовую идею прикладного слоя: все сценарии начинаются с x2(T).}
\end{frame}

% G2
\begin{frame}{Для какого растворителя молекула лучше растворяется?}
  \scriptsize
  Пример: ибупрофен при 298 K, скрининг по 65 растворителям.
  \vspace{0.3em}
  \begin{tabular}{lcccl}
    \toprule
    \textbf{Растворитель} & $\ln x_2$ & $x_2$ & \textbf{mg/mL*} & \textbf{Комментарий}\\
    \midrule
    DMSO & $-1.23$ & 0.29 & $\sim380$ & полярный, высокая загрузка\\
    Ацетон & $-2.01$ & 0.134 & $\sim170$ & летучий кетон\\
    Этанол & $-2.45$ & 0.087 & $\sim120$ & более зелёный спирт\\
    Гексан & $-5.67$ & 0.003 & $\sim3.5$ & низкая совместимость\\
    Вода & $-8.12$ & 0.0003 & $\sim0.04$ & водное ограничение\\
    \bottomrule
  \end{tabular}
  \par\vspace{0.45em}
  \begin{minipage}{0.96\textwidth}
    \scriptsize\color{TGSlate}
    *mg/mL — грубый пересчёт. В UI можно фильтровать по токсичности, цене, температуре кипения и зелёному рейтингу растворителя.
  \end{minipage}
  \note{Смысл отбора растворителей: быстро отсеять среды с низкой совместимостью и выбрать кандидатов для лабораторной проверки.}
\end{frame}

% G3
\begin{frame}{При какой температуре кристаллизовать?}
  \small
  \begin{columns}[T,onlytextwidth]
    \begin{column}{0.52\textwidth}
      Охлаждающая кристаллизация:
      \[
        x_2(T_{hot}=340K)=0.15,\qquad
        x_2(T_{cold}=280K)=0.03.
      \]
      \[
        Y_{theory}=\frac{0.15-0.03}{0.15}=80\%.
      \]
      \begin{itemize}\tightitems
        \item считаем горячую загрузку;
        \item остаток в маточном растворе;
        \item теоретический выход и прокси пересыщения.
      \end{itemize}
    \end{column}
    \begin{column}{0.44\textwidth}
      \plotasset[width=\linewidth]{ideal_sle_example.pdf}
      \smallnote{Ограничение: это не кинетическая модель зародышеобразования/роста и не модель полиморфизма.}
    \end{column}
  \end{columns}
  \note{Показать, что x2(T) напрямую превращается в технологическое окно кристаллизации.}
\end{frame}

% G4
\begin{frame}{Антирастворительная кристаллизация: добавляем антирастворитель}
  \small
  Сценарий: вещество хорошо растворяется в DMSO, слабо — в воде.
  \[
    x_2(\mathrm{DMSO},298K)=0.29,\qquad
    x_2(\mathrm{H_2O},298K)=0.0003.
  \]
  \[
    \mathrm{отношение\ антирастворителя}\sim\frac{0.29}{0.0003}\approx10^3.
  \]
  \begin{columns}[T,onlytextwidth]
    \begin{column}{0.48\textwidth}
      \begin{block}{Что модель даёт}
        \scriptsize
        ранжирование пар растворитель/антирастворитель, температурные окна, грубая движущая сила выпадения.
      \end{block}
    \end{column}
    \begin{column}{0.48\textwidth}
      \begin{block}{Чего модель не даёт}
        \scriptsize
        кинетику смешения, динамику пересыщения, форму кристаллов и барьер зародышеобразования.
      \end{block}
    \end{column}
  \end{columns}
  \note{Антирастворительный сценарий — хороший пример поддержки принятия решений, где нужны не абсолютные истины, а ранжирование вариантов.}
\end{frame}

% G5
\begin{frame}{Можно ли заменить хлороформ на что-то экологичнее?}
  \small
  Замена на более зелёный растворитель:
  \[
    x_2(\mathrm{кандидат})\ge0.8\,x_2(\mathrm{хлороформ})
  \]
  плюс ограничения на температуру кипения, совместимость и безопасность.
  \vspace{0.4em}
  \begin{center}
  \begin{tabular}{lccc}
    \toprule
    \textbf{Растворитель} & относит. $x_2$ & зелёный рейтинг & решение\\
    \midrule
    Хлороформ & 1.00 & низкий & база\\
    Этилацетат & 0.92 & высокий & кандидат\\
    Ацетон & 0.84 & средний & кандидат\\
    Гексан & 0.31 & низкий & отклонить\\
    \bottomrule
  \end{tabular}
  \end{center}
  \note{Это не регуляторный движок зелёной химии, но быстрый фильтр по растворимости для замены токсичного растворителя.}
\end{frame}

% G6
\begin{frame}{BCS-классификация: первичный скрининг}
  \small
  \begin{columns}[T,onlytextwidth]
    \begin{column}{0.52\textwidth}
      \begin{block}{Что такое BCS}
        \scriptsize
        Biopharmaceutics Classification System делит кандидаты по двум вопросам: \textbf{растворится ли доза в жидкости ЖКТ?} и \textbf{пройдёт ли молекула через мембрану?}
      \end{block}
      \begin{block}{Что делает TGNN-Solv}
        \scriptsize
        Модель даёт водную растворимость $S$. Затем приложение сравнивает дозу с доступным объёмом:
        \[
          D_0=\frac{\mathrm{Dose}}{V\cdot S}.
        \]
        Если $D_0\le1$, доза помещается в раствор; если $D_0>1$, растворимость может ограничивать всасывание.
      \end{block}
      \smallnote{Проницаемость TGNN-Solv не измеряет: используется дескрипторный прокси, поэтому это первичная сортировка кандидатов, а не регуляторная BCS-классификация.}
    \end{column}
    \begin{column}{0.44\textwidth}
      \fitfig{
      \begin{tikzpicture}
        \draw[->] (0,0)--(5,0) node[right] {проницаемость};
        \draw[->] (0,0)--(0,3.8) node[above] {растворимость};
        \draw[TGLight] (2.5,0)--(2.5,3.5);
        \draw[TGLight] (0,1.9)--(4.8,1.9);
        \node at (1.1,2.8) {III};
        \node at (3.6,2.8) {I};
        \node at (1.1,0.8) {IV};
        \node at (3.6,0.8) {II};
        \fill[TGOrange] (3.6,0.8) circle (0.08);
      \end{tikzpicture}}
      {\scriptsize
      \begin{tabular}{ll}
        \toprule
        I & высокая раств. + высокая прониц.\\
        II & низкая раств. + высокая прониц.\\
        III & высокая раств. + низкая прониц.\\
        IV & низкая раств. + низкая прониц.\\
        \bottomrule
      \end{tabular}
      }
    \end{column}
  \end{columns}
  \note{Для BCS нужен явный дисклеймер: это первичная оценка пригодности кандидата, а не регуляторная классификация.}
\end{frame}

% G7
\begin{frame}{Растворимость вдоль ЖКТ: где препарат растворится?}
  \small
  \begin{columns}[T,onlytextwidth]
    \begin{column}{0.50\textwidth}
      \begin{block}{Идея}
        \scriptsize
        В разных отделах ЖКТ разные pH, объём жидкости и состав среды. Поэтому одно число ``водная растворимость при 298 K'' превращается в профиль: где доза может раствориться, а где останется твёрдой фазой.
      \end{block}
      \begin{enumerate}\tightitems
        \item TGNN-Solv предсказывает базовую растворимость.
        \item Приложение применяет pH- и средовые поправки.
        \item Доза сравнивается с объёмом компартмента.
      \end{enumerate}
    \end{column}
    \begin{column}{0.46\textwidth}
      \begin{tabularx}{\textwidth}{lccY}
        \toprule
        \textbf{Отдел} & pH & V & \textbf{Смысл}\\
        \midrule
        Желудок & 1.2 & 250 & кислая среда\\
        Двенадцатиперстная & 5.8 & 100 & буфер + желчь\\
        Тощая & 6.2 & 200 & прокси FaSSIF / FeSSIF\\
        Подвздошная & 7.4 & 200 & сдвиг ионизации\\
        \bottomrule
      \end{tabularx}
      \[
        f_{\mathrm{dissolved}}\approx
        \min\!\left(1,\frac{V\cdot S_{\mathrm{corr}}}{\mathrm{Dose}}\right).
      \]
    \end{column}
  \end{columns}
  \keybox{Это не PBPK, не модель абсорбции и не PK/PD; это оценка ``давления растворения'' до абсорбции.}
  \note{Пояснить, что PK-profile в приложении — профиль растворения по средам ЖКТ, а не полноценная фармакокинетика.}
\end{frame}

% G8
\begin{frame}{Отбор растворителей для заданного маршрута синтеза}
  \small
  \begin{columns}[T,onlytextwidth]
    \begin{column}{0.52\textwidth}
      \begin{block}{Что НЕ делает приложение}
        \scriptsize
        Оно не придумывает синтез и не строит ретросинтетический маршрут. Маршрут задаёт химик: шаги, температуры и список допустимых растворителей.
      \end{block}
      \begin{block}{Что оно делает}
        \scriptsize
        Для каждого шага и растворителя модель считает растворимость при горячей и холодной температуре. Хороший кандидат даёт высокую загрузку на реакции и низкую остаточную растворимость при выделении.
      \end{block}
      \[
        \text{температурный размах растворимости}
        =
        \frac{x_2(T_{\mathrm{hot}})}{x_2(T_{\mathrm{cold}})}.
      \]
    \end{column}
    \begin{column}{0.44\textwidth}
      \begin{tabularx}{\textwidth}{lYY}
        \toprule
        \textbf{Шаг} & \textbf{Что задано} & \textbf{Что ранжируем}\\
        \midrule
        $A\to B$ & $T_{\text{реакц.}}$, растворители X/Y/Z & загрузка + выделение\\
        $B\to C$ & $T_{\text{реакц.}}$, растворители U/V/W & окно кристаллизации\\
        \bottomrule
      \end{tabularx}
    \end{column}
  \end{columns}
  \par\vspace{0.45em}
  \begin{minipage}{0.94\textwidth}
    \scriptsize\color{TGSlate}
    Оценка — прозрачная эвристика поверх $x_2(T)$: горячая загрузка, холодное выделение, температурный размах растворимости. Это не отдельная модель машинного обучения.
  \end{minipage}
  \note{Снять возможное недопонимание: приложение помогает выбрать растворители для уже заданного маршрута.}
\end{frame}

% J1
\begin{frame}{Experiment Lab: интерактивная поверхность проекта}
  \small
  Запуск:
  \[
    \texttt{python scripts/launch\_lab.py}
  \]
  \begin{columns}[T,onlytextwidth]
    \begin{column}{0.48\textwidth}
      \begin{itemize}\tightitems
	        \item обучение, эксперименты, HPO;
	        \item Pipeline Studio и Model Architect;
	        \item результаты и графики, Benchmark Studio;
	        \item предсказания, приложения, планировщик;
	        \item документация, воспроизведение, Job Center, окружение.
      \end{itemize}
    \end{column}
    \begin{column}{0.48\textwidth}
      \begin{block}{Ключевой дизайн}
        \scriptsize
	        UI может жить в лёгком окружении, а тяжёлые задачи обучения и предсказания выполняются через настраиваемую Python-команду.
      \end{block}
    \end{column}
  \end{columns}
  \note{Experiment Lab — не отдельная научная модель, а рабочая поверхность, которая связывает обучение, предсказания, приложения и документацию.}
\end{frame}

% J2
\begin{frame}{Демо предсказания: парацетамол в этаноле}
  \scriptsize
  Вход: \texttt{CC(=O)Nc1ccc(O)cc1} + \texttt{CCO}, $T=298.15$ K.
  \vspace{0.35em}
  \begin{tabularx}{\textwidth}{lYlY}
    \toprule
    $\ln x_2$ & $-3.42$ & $x_2$ & 0.033\\
    $T_m$ & 442 K & $\Delta H_{fus}$ & 26\,100 J/mol\\
    $\tau_{12}$ & 1.23 & $\tau_{21}$ & 0.87\\
    $\alpha$ & 0.31 & $\Phi$ & 4.12\\
	    $\ln\gamma_2$ & 0.70 & коррекция & +0.15, шлюз 0.82\\
    \bottomrule
  \end{tabularx}
  \vspace{0.5em}
  \begin{block}{Интерпретация}
    \scriptsize
    Умеренная растворимость. Доминирует кристаллический штраф ($\Phi=4.12$), неидеальность умеренная ($\gamma_2\approx2.0$), коррекция малая.
  \end{block}
  \note{Показать, что предсказание возвращает не только число, но и физическую декомпозицию.}
\end{frame}

% J3
\begin{frame}{Демо температурного скана: кривая $\ln x_2(T)$}
  \small
  \begin{columns}[T,onlytextwidth]
    \begin{column}{0.55\textwidth}
      \plotasset[width=\linewidth]{ideal_sle_example.pdf}
    \end{column}
    \begin{column}{0.41\textwidth}
      \begin{block}{Что смотреть}
        TGNN-Solv: физическая кривая через $\Phi+\ln\gamma_2$.\\
        DirectGNN: прямая базовая модель машинного обучения.\\
        Идеальный SLE: серый аналитический ориентир.\\
        Экспериментальные точки: если доступны.
      \end{block}
      \begin{block}{Ключевой смысл}
        Температурная зависимость нужна для проектирования процесса: выбора $T_{hot}$, $T_{cold}$ и оценки температурного размаха растворимости.
      \end{block}
    \end{column}
  \end{columns}
  \note{Температурная кривая — главный прикладной продукт модели, потому что проектирование процесса требует зависимости от T, а не одиночную точку.}
\end{frame}

% 44
\begin{frame}{Резюме и вклады}
  \small
  \begin{enumerate}\tightitems
    \item \textbf{TIMP}: термодинамическая декомпозиция в передаче сообщений;
    на коротком бюджете: линейная проба $R^2$ $+0.036$, косинус каналов $0.006$,
    штраф MAE $+0.08$.
    \[
      \mathbf{m}_{ij}
      =
      \phi_{\mathrm{disp}}\sqrt{\alpha_i\alpha_j}
      +
      \phi_{\mathrm{polar}}\operatorname{sp}(\delta q_i\delta q_j).
    \]
    \item \textbf{Дифференцируемый SLE-конвейер}:
    \[
      \frac{d\ln x_2^*}{d\theta}
      =
      -\frac{\partial(\Phi+\ln\gamma_2)/\partial\theta}{1+x_2^*\eta}.
    \]
    \item \textbf{Систематическая диагностика}: DirectGNN обошёл лучший RF;
    physics bottleneck $=0.09$ MAE; подстановка истинных параметров показала, что кристаллическая голова
    не является главным bottleneck.
    \item \textbf{Инженерия функции потерь}: поэтапное обучение, доля основной потери $>50\%$, подстановка истинных параметров.
    \item \textbf{Прикладной слой}: от отбора растворителей до BCS-подобной оценки — 4 модуля поверх единого $x_2(T)$, поддерживаемые для TGNN и DirectGNN.
    \item \textbf{Новые гипотезы}: H9 — структурированным представлениям TIMP
    нужен больший бюджет NRTL-головы; H10 — слой взаимодействия недополучает
    градиент в физическом пути; вспомогательная ветка подняла градиент
    попарного внимания до уровня DirectGNN ($0.99\times$).
  \end{enumerate}
  \note{Сжато пройти вклады и вернуть аудиторию к исходному контрольному вопросу.}
\end{frame}

% 45
\begin{frame}{Благодарности / вопросы}
  \begin{center}
    {\Large Спасибо всем за внимание!}\\[1.1em]
    {\large Вопросы?}\\[1.3em]
    \[
      \boxed{
      \ln x_2
      =
      -\frac{\Delta H_{\mathrm{fus}}}{R}
      \left(\frac{1}{T}-\frac{1}{T_m}\right)
      -
      \ln\gamma_2(\tau_{12},\tau_{21},\alpha)}
    \]
  \end{center}
  \note{На вопросах вернуться к трём критериям: представление, температурная экстраполяция, идентифицируемость.}
\end{frame}

\appendix
\section{Дополнение}

\begin{frame}{Дополнение: зачем нужно демпфирование}
  \small
  \begin{columns}[T,onlytextwidth]
    \begin{column}{0.55\textwidth}
      \begin{block}{Обновляем не напрямую, а с under-relaxation}
        \scriptsize
        Вместо полного шага
        \[
          x_{k+1}=g(x_k)
        \]
        используем демпфированный оператор
        \[
          x_{k+1}=h_\lambda(x_k)
          =(1-\lambda)x_k+\lambda g(x_k),
          \qquad 0<\lambda\le 1.
        \]
        В неподвижной точке $x^*=g(x^*)$:
        \[
          h_\lambda'(x^*)
          =1-\lambda+\lambda g'(x^*).
        \]
      \end{block}
      \begin{block}{Что это меняет}
        \scriptsize
        Если $g'(x^*)<0$, полный шаг может давать
        колебания около решения. Тогда уменьшение $\lambda$
        сжимает эффективный наклон:
        \[
          |h'_\lambda(x^*)| < |g'(x^*)|.
        \]
      \end{block}
    \end{column}
    \begin{column}{0.41\textwidth}
      \begin{block}{Пример}
        \scriptsize
        Пусть $g'(x^*)=-0.8$.
        Тогда
        \[
          h'_{1}(x^*)=-0.8,
          \qquad
          h'_{0.5}(x^*)=0.1.
        \]
        То есть oscillation почти исчезает.
      \end{block}
      \keybox{\scriptsize Демпфирование полезно не «всегда», а тогда, когда проблема — колебательный шаг вокруг решения, а не истинная экспансия отображения.}
    \end{column}
  \end{columns}
  \note{Смысл слайда в строгом утверждении: демпфирование уменьшает колебания, когда производная отрицательна или близка к минус единице. Это не магическая кнопка от любой расходимости.}
\end{frame}

\begin{frame}{Дополнение: условие сходимости демпфированного шага}
  \small
  \begin{columns}[T,onlytextwidth]
    \begin{column}{0.56\textwidth}
      \begin{block}{Локальное условие}
        \scriptsize
        Для локальной сходимости нужно
        \[
          |h'_\lambda(x^*)|<1.
        \]
        Подставляем
        \[
          h'_\lambda(x^*)=1-\lambda+\lambda g'(x^*)
          =1-\lambda\bigl(1-g'(x^*)\bigr).
        \]
        Тогда
        \[
          |1-\lambda(1-g')|<1
        \]
        \[
          \Longleftrightarrow
          -1<1-\lambda(1-g')<1
        \]
        \[
          \Longleftrightarrow
          0<\lambda(1-g')<2.
        \]
      \end{block}
    \end{column}
    \begin{column}{0.40\textwidth}
      \begin{block}{Следствия}
        \scriptsize
        \[
          \boxed{g'(x^*)<1}
        \]
        — необходимое условие.

        \vspace{0.35em}
        Если $g'(x^*)<1$, то можно взять
        \[
          0<\lambda<\frac{2}{1-g'(x^*)}.
        \]

        \vspace{0.35em}
        Если $g'(x^*)>1$, простое under-relaxation
        уже не спасает: нужен другой решатель
        или другая параметризация.
      \end{block}
    \end{column}
  \end{columns}
  \note{Здесь полезно явно проговорить границу применимости: under-relaxation работает только пока локальный наклон меньше единицы. Если наклон больше единицы и положителен, проблема глубже.}
\end{frame}

\begin{frame}{Дополнение: почему знаменатель в implicit gradient положителен}
  \small
  \begin{columns}[T,onlytextwidth]
    \begin{column}{0.58\textwidth}
      \begin{block}{Шаг 1: неявная функция}
        \scriptsize
        \[
          F(x^*,\theta)=x^*-g(x^*;\theta)=0,
        \]
        \[
          \frac{dx^*}{d\theta_j}
          =
          -\left(\frac{\partial F}{\partial x^*}\right)^{-1}
          \frac{\partial F}{\partial\theta_j}.
        \]
        Причём
        \[
          \frac{\partial F}{\partial x^*}
          =1-g'(x^*).
        \]
      \end{block}
      \begin{block}{Шаг 2: связь с NRTL}
        \scriptsize
        Для SLE-задачи
        \[
          g'(x^*)=-x_2^*\,\eta,
          \qquad
          \eta=\frac{\partial\ln\gamma_2}{\partial x_2}\Big|_{x_2^*},
        \]
        поэтому
        \[
          1-g'(x^*) = 1+x_2^*\eta.
        \]
      \end{block}
    \end{column}
    \begin{column}{0.38\textwidth}
      \begin{block}{Почему знак положительный}
        \scriptsize
        Для недемпфированного шага локальная сходимость даёт
        \[
          |g'(x^*)|<1
          \;\Longrightarrow\;
          1-g'(x^*)>0.
        \]
        Для демпфированного шага:
        \[
          0<\lambda(1-g'(x^*))<2.
        \]
        При $\lambda>0$ это тоже означает
        \[
          1-g'(x^*)>0.
        \]
      \end{block}
      \keybox{\scriptsize Знаменатель в формуле implicit differentiation положителен именно потому, что мы находимся в локально сходящемся режиме fixed-point iteration.}
    \end{column}
  \end{columns}
  \note{Этим слайдом закрываем самый частый математический вопрос: почему знаменатель не меняет знак. Ответ: это прямое следствие локальной сходимости итератора.}
\end{frame}

\begin{frame}{Дополнение: какие атомы влияют на прогноз?}
  \small
  \begin{columns}[T,onlytextwidth]
    \begin{column}{0.49\textwidth}
      \centering
      \plotasset[width=\linewidth]{attribution_tgnn_mpnn_paracetamol_ethanol.png}\\[-0.2em]
      {\scriptsize TGNN MPNN}
    \end{column}
    \begin{column}{0.49\textwidth}
      \centering
      \plotasset[width=\linewidth]{attribution_tgnn_timp_paracetamol_ethanol.png}\\[-0.2em]
      {\scriptsize TGNN TIMP}
    \end{column}
  \end{columns}
  \vspace{0.35em}
  \smallnote{Интегральные градиенты: красный атом увеличивает предсказанный $\ln x_2$, синий уменьшает, белый почти нейтрален.}
  \note{Это резервный слайд для вопроса об интерпретируемости. Важный смысл: градиент проходит от $\ln x_2$ через решатель, NRTL-голову, взаимодействие молекул и энкодер до атомных признаков.}
\end{frame}

\begin{frame}{Дополнение: где теряется градиент?}
  \small
  \begin{columns}[T,onlytextwidth]
    \begin{column}{0.55\textwidth}
      \plotasset[width=\linewidth]{gradient_flow.pdf}
    \end{column}
    \begin{column}{0.41\textwidth}
      \begin{block}{Диагноз H10}
        \scriptsize
        Энкодер получает градиент, но слой взаимодействия TGNN
        получает среднюю норму около $7\cdot10^{-4}$ против
        $2.3\cdot10^{-2}$ у DirectGNN.
      \end{block}
      \smallnote{Следствие: кристаллическая голова обучает представление вещества,
      а совместимость вещества и растворителя остаётся слабее обученной.}
    \end{column}
  \end{columns}
  \note{Использовать этот слайд, если спросят, почему подстановка истинных кристаллических параметров не помогает и почему нужна вспомогательная ветка на пару молекул.}
\end{frame}

\begin{frame}{Дополнение: когда энкодер учит химию?}
  \small
  \begin{columns}[T,onlytextwidth]
    \begin{column}{0.61\textwidth}
      \centering
      \begin{tikzpicture}[x=1cm,y=1cm,>=Latex]
        \draw[TGSlate!35,line width=1.2pt] (0.4,2.2) -- (7.8,2.2);
        \foreach \x/\lab in {0.6/0,2.2/50,5.5/250,7.6/300} {
          \draw[TGSlate!55,line width=1pt] (\x,2.05) -- (\x,2.35);
          \node[font=\scriptsize,text=TGSlate!80] at (\x,1.82) {\lab};
        }
        \fill[TGBlue!20,rounded corners=4pt] (0.55,2.45) rectangle (2.1,3.05);
        \fill[TGGreen!18,rounded corners=4pt] (2.2,2.45) rectangle (5.45,3.05);
        \fill[TGOrange!18,rounded corners=4pt] (5.55,2.45) rectangle (7.55,3.05);
        \node[font=\scriptsize\bfseries,text=TGBlue!80] at (1.32,2.75) {Фаза 1};
        \node[font=\scriptsize\bfseries,text=TGGreen!70!black] at (3.82,2.75) {Фаза 2};
        \node[font=\scriptsize\bfseries,text=TGOrange!80!black] at (6.55,2.75) {Фаза 3};

        \foreach \x in {1.2,2.8,3.5,4.2,4.9,6.4,7.1} {
          \fill[TGRed] (\x,2.2) circle (1.8pt);
          \draw[TGRed!70,line width=0.7pt] (\x,2.2) -- (\x,1.35);
        }

        \node[draw=TGSlate!35,fill=white,rounded corners=6pt,
          minimum width=2.45cm,minimum height=0.95cm,
          align=center,font=\scriptsize] (freeze) at (1.7,0.7)
          {заморозить\\энкодер};
        \node[draw=TGSlate!35,fill=white,rounded corners=6pt,
          minimum width=2.45cm,minimum height=0.95cm,
          align=center,font=\scriptsize] (probe) at (4.1,0.7)
          {Ridge: эмбеддинг\\$\rightarrow$ дескрипторы};
        \node[draw=TGSlate!35,fill=white,rounded corners=6pt,
          minimum width=2.45cm,minimum height=0.95cm,
          align=center,font=\scriptsize] (track) at (6.45,0.7)
          {$R^2$ по эпохам\\и фазам};

        \draw[->,TGSlate!70,line width=0.9pt] (freeze.east) -- (probe.west);
        \draw[->,TGSlate!70,line width=0.9pt] (probe.east) -- (track.west);

        \node[anchor=west,font=\scriptsize,text=TGRed!85] at (0.55,3.38)
          {красные точки — моменты съёма linear probe};
      \end{tikzpicture}
    \end{column}
    \begin{column}{0.35\textwidth}
      \begin{block}{Как читать}
        \scriptsize
        Вместо одного финального probe мы периодически
        снимаем срез представления по эпохам:
        замораживаем энкодер, восстанавливаем дескрипторы
        и строим кривую $R^2(t)$.
      \end{block}
      \smallnote{Реальная кривая появляется только в отдельном прогоне
      с флагом \texttt{--probe-every}. На этой версии слайда показан
      сам протокол измерения, чтобы не подменять его фиктивными данными.}
    \end{column}
  \end{columns}
  \note{Если возникнет вопрос, почему здесь нет готовой кривой: диагностика уже реализована в коде, но требует отдельного тренировочного прогона с периодическим probe. На слайде показан именно протокол, а не вымышленный результат.}
\end{frame}

\begin{frame}{Дополнение: чувствительность по тестовым примерам}
  \small
  \begin{columns}[T,onlytextwidth]
    \begin{column}{0.50\textwidth}
      \plotasset[width=\linewidth]{sensitivity_boxplot.pdf}
    \end{column}
    \begin{column}{0.46\textwidth}
      \plotasset[width=\linewidth]{sensitivity_heatmap.pdf}
    \end{column}
  \end{columns}
  \vspace{0.25em}
  \smallnote{Главный вывод здесь уже не про отдельные выбросы, а про структуру задачи: кристаллические параметры во всём тесте существенно слабее контактных. При очень низкой растворимости сильнее всего влияет $\alpha$, затем доминируют $\tau_{12}$ и $\tau_{21}$. Связь чувствительности с ошибкой остаётся слабой ($|\rho|\le0.15$), поэтому одной иерархии чувствительности недостаточно, чтобы локализовать источник ошибки.}
  \note{Слева читается глобальная иерархия: по медиане контактные параметры намного жёстче кристаллических. Справа видно, что лидер внутри контактных параметров меняется по диапазонам растворимости: в самом нерастворимом режиме выходит вперёд $\alpha$, дальше доминируют $\tau_{12}$ и $\tau_{21}$.}
\end{frame}

\begin{frame}{Дополнение: что выучили TIMP-каналы?}
  \small
  \begin{columns}[T,onlytextwidth]
    \begin{column}{0.54\textwidth}
      \plotasset[width=\linewidth]{timp_channel_probe_matrix.pdf}
    \end{column}
    \begin{column}{0.42\textwidth}
      \begin{block}{Практический вывод}
        \scriptsize
        В тесте есть только одно уникальное вещество
        с экспериментальными Hansen-метками, поэтому строки с $\delta$ показаны
        через явно помеченный псевдо-Hansen-прокси.

        Чистой декомпозиции пока нет: полярный канал в этом прогоне
        чаще сильнее даже по общим дескрипторам. Это согласуется
        с полярным корпусом и коротким бюджетом обучения.
      \end{block}
      \smallnote{Корректный вывод здесь не «каналы идеально чистые», а
      «каналы различимы, но семантическое разделение пока неполное».}
    \end{column}
  \end{columns}
  \note{Важно не переобещать. На этом split мы не можем честно заявлять восстановление экспериментальных параметров Хансена: меток в тесте почти нет. Поэтому интерпретация осторожная: каналы различаются, но не разделились в чистые физические оси.}
\end{frame}

\begin{frame}{Дополнение: неполярные и полярные молекулы в TIMP}
  \small
  \begin{columns}[T,onlytextwidth]
    \begin{column}{0.62\textwidth}
      \plotasset[width=\linewidth]{timp_channel_ratio_vs_logp.pdf}
    \end{column}
    \begin{column}{0.34\textwidth}
      \begin{block}{Доля дисперсии}
        \scriptsize
        \[
          \rho =
          \frac{\|\mathbf{g}_{disp}\|}
          {\|\mathbf{g}_{disp}\|+\|\mathbf{g}_{polar}\|}
        \]
        Линия регрессии и медианы по корзинам показывают
        умеренную положительную связь с MolLogP.

        Значит дисперсионный канал всё же связан с
        липофильностью, но связь пока не настолько сильна,
        чтобы говорить о полном разделении каналов.
      \end{block}
    \end{column}
  \end{columns}
  \note{Это простой sanity-check: если канал дисперсии не связан с липофильностью, физическая интерпретация TIMP сомнительна.}
\end{frame}

\begin{frame}{Дополнение: распределение весов по блокам}
  \small
  \begin{columns}[T,onlytextwidth]
    \begin{column}{0.62\textwidth}
      \plotasset[width=\linewidth]{weight_violin.pdf}
    \end{column}
    \begin{column}{0.34\textwidth}
      \begin{block}{Зачем это нужно}
        \scriptsize
        Слева показаны центральные 99\% значений весов,
        справа — доли почти нулевых и экстремальных параметров.

        Практический вывод: массового зануления нет, а
        экстремальные веса сосредоточены не в TIMP-каналах,
        а в служебных блоках и хвосте модели.
      \end{block}
      \smallnote{Выбросы не исчезли: они сохранены в CSV-статистике.
      На слайде скрыта только их визуальная доминация.}
    \end{column}
  \end{columns}
  \note{Показывать, если возникнет вопрос о стабильности TIMP или о том, не схлопнулись ли каналы.}
\end{frame}

\begin{frame}{Дополнение: Wilson и UNIQUAC}
  \small
  \centering
  \setlength{\tabcolsep}{5pt}
  \begin{tabular}{lccc}
    \toprule
    \textbf{Модель} & \textbf{Обуч. парам.} & \textbf{LLE} & \textbf{Восстановимость параметров}\\
    \midrule
    NRTL & 3 & да & слабее\\
    Wilson & 2 & нет & лучше\\
    UNIQUAC & 2 & да & лучше\\
    \bottomrule
  \end{tabular}

  \vspace{0.45em}
  \begin{columns}[T,onlytextwidth]
    \begin{column}{0.58\textwidth}
      \begin{block}{Wilson: меньше пространство поиска}
        \scriptsize
        \[
          \frac{G^E}{RT}
          =
          -x_1\ln(x_1+\Lambda_{12}x_2)
          -x_2\ln(x_2+\Lambda_{21}x_1).
        \]
        Два параметра $\Lambda_{12},\Lambda_{21}$ проще восстановить из данных, но модель не описывает равновесие жидкость—жидкость.
      \end{block}
      \begin{block}{UNIQUAC: размер и форма отдельно}
        \scriptsize
        \[
          \ln\gamma_2
          =
          \ln\gamma_2^{comb}(r,q)
          +
          \ln\gamma_2^{res}(\tau_{12},\tau_{21}).
        \]
        $r,q$ отвечают за размер и форму молекул; $\tau$ отвечает за остаточные энергии контакта.
      \end{block}
    \end{column}
    \begin{column}{0.38\textwidth}
      \plotasset[width=\linewidth]{nrtl_gamma_example.pdf}
      \smallnote{NRTL остаётся вариантом по умолчанию; Wilson/UNIQUAC нужны как абляции на восстановимость параметров.}
    \end{column}
  \end{columns}
  \note{Показывать только при вопросе об альтернативных моделях избыточной энергии Гиббса.}
\end{frame}

\begin{frame}{Дополнение: запланированные эксперименты}
  \scriptsize
  \begin{tabularx}{\textwidth}{clYY}
    \toprule
    \# & \textbf{Эксперимент} & \textbf{Тестирует} & \textbf{Статус}\\
    \midrule
    1 & Полный бюджет 50/200/50, 3 случайные инициализации & физический путь против DirectGNN & запланировано\\
    2 & Oracle-оценка & bottleneck кристаллического блока & запланировано\\
    3 & Абляция bridge-потери & вред регуляризатора Regular Solution & запланировано\\
    4 & TGNN + дескрипторы & добавление экспертных признаков & запланировано\\
    5 & TIMP против MPNN & эффект двух физических каналов & запланировано\\
    6 & Wilson vs NRTL & переизбыточная параметризация & запланировано\\
    7 & Температурная экстраполяция & главный тезис физической модели & запланировано\\
    \bottomrule
  \end{tabularx}
  \note{Показать только если спросят о планах: основной доклад должен закончиться результатами и выводами.}
\end{frame}

\begin{frame}{Дополнение: план абляционного исследования}
  \scriptsize
  \begin{tabularx}{\textwidth}{lYY}
    \toprule
    \textbf{Абляция} & \textbf{Что проверяет} & \textbf{Ожидание}\\
    \midrule
    Без перекрёстного внимания & полезность явного взаимодействия молекул & ухудшение\\
    Без коррекции & полезность ограниченного клапана & ухудшение\\
    Без вспомогательных потерь & полезность многозадачного обучения & сильное ухудшение\\
    Без bridge-потери & тянет ли Regular Solution в неверную сторону & возможно улучшение\\
    $T$ в энкодер & нужна ли изоляция температуры & ухудшение\\
    $\gamma_2=1$ (идеально) & вклад неидеальности раствора & сильное ухудшение\\
    \bottomrule
  \end{tabularx}
  \note{Абляции должны проверять причинные утверждения, а не просто перебирать галочки.}
\end{frame}

\begin{frame}{Дополнение: дорожная карта}
  \fitfig{
  \begin{tikzpicture}
    \draw[->] (0,0)--(12.8,0) node[right] {недели};
    \foreach \x/\w in {0/1,3/2,6/3,9/4}
      \draw (\x,0.08)--(\x,-0.08) node[below,font=\scriptsize] {\w};
    \node[anchor=east,font=\scriptsize] at (-0.1,2.4) {Код};
    \node[anchor=east,font=\scriptsize] at (-0.1,1.6) {GPU};
    \node[anchor=east,font=\scriptsize] at (-0.1,0.8) {Анализ};
    \fill[TGBlue!55] (0,2.2) rectangle (2.8,2.6) node[midway,font=\scriptsize,white] {TIMP/Wilson/тесты};
    \fill[TGOrange!65] (3.0,1.4) rectangle (6.0,1.8) node[midway,font=\scriptsize,white] {3 запуска TGNN/Direct/RF};
    \fill[TGOrange!65] (6.2,1.4) rectangle (9.0,1.8) node[midway,font=\scriptsize,white] {TIMP/дескр./Wilson};
    \fill[TGGreen!60] (6.2,0.6) rectangle (8.8,1.0) node[midway,font=\scriptsize,white] {T-экстраполяция};
    \fill[TGPurple!55] (9.1,0.6) rectangle (12.0,1.0) node[midway,font=\scriptsize,white] {финализация таблиц};
  \end{tikzpicture}}
  \note{Это backup-слайд: показывать только если будет вопрос о следующих шагах.}
\end{frame}

\end{document}
